\documentclass[11pt,a4paper,oneside]{book}

% --- Encoding & fonts ---
\usepackage[utf8]{inputenc}
\usepackage[T1]{fontenc}
\usepackage{lmodern}
\usepackage[english]{babel}
\usepackage{microtype}            % better justification, fewer overfull lines

% --- Page geometry ---
\usepackage[a4paper,margin=1in]{geometry}

% --- Line spacing ---
\usepackage{setspace}
\onehalfspacing

% --- Headings ---
\usepackage{titlesec}
\titleformat{\chapter}[display]
  {\normalfont\huge\bfseries\filcenter}
  {}
  {0pt}
  {\Huge}
\titlespacing*{\chapter}{0pt}{20pt}{40pt}

% --- Headers/footers: page number centered in footer ---
\usepackage{fancyhdr}
\pagestyle{fancy}
\fancyhf{}
\fancyfoot[C]{\thepage}
\renewcommand{\headrulewidth}{0pt}
\fancypagestyle{plain}{\fancyhf{}\fancyfoot[C]{\thepage}\renewcommand{\headrulewidth}{0pt}}

% --- Widow/orphan control ---
\widowpenalty=10000
\clubpenalty=10000

% --- Paragraphs: justified with first-line indent (book style) ---
\setlength{\parindent}{1.5em}
\setlength{\parskip}{0pt}

% --- Table of contents: show chapters only (sections excluded) ---
\setcounter{tocdepth}{0}
% Number sections internally but do not print chapter numbers as "Chapter N"
% (titles already read 'Chapter N: ...'). We use \chapter* + manual TOC entry.

\title{Capital That Serves Life}
\author{}
\date{}

\begin{document}

% ===================== TITLE PAGE =====================
\begin{titlepage}
  \centering
  \vspace*{0.32\textheight}
  {\Huge\bfseries Capital That Serves Life\par}
  \vspace{1.2em}
  {\Large\itshape Recovering Moral Economy in an Age of Extraction\par}
  \vspace{2em}
  {\large\itshape Working manuscript\par}
  \vfill
\end{titlepage}

% ===================== TABLE OF CONTENTS =====================
\frontmatter
\pagestyle{plain}
\tableofcontents
\clearpage

\mainmatter

\chapter*{Prologue: The Question That Was Buried}
\addcontentsline{toc}{chapter}{Prologue: The Question That Was Buried}
\section*{The Question That Was Buried}

This is not a neutral history of economics. It is a retrieval. It begins from the conviction that modern economic life has forgotten one of the oldest and most necessary questions human beings can ask: what is an economy for?

This question once stood near the center of moral, philosophical, and theological reflection. It was not considered strange to ask whether wealth formed or deformed the soul, whether lending served necessity or exploited it, whether trade strengthened community or dissolved it, whether money remained connected to real life or became an abstract power over it. To speak about economy was also to speak about household, city, land, obligation, virtue, justice, mercy, and the common good.

Today, the question often sounds naive. We are trained to ask different questions. Does it grow? Does it scale? Does it maximize shareholder value? Is it innovative? Does it attract capital? These are not meaningless questions. A society that cannot produce, coordinate, allocate, build, trade, and innovate will not serve human beings well. Good intentions do not abolish scarcity. Moral seriousness does not remove the need for prices, incentives, institutions, and practical judgment. But these are not ultimate questions.

An economy can grow while degrading the lives of the people inside it. It can become more efficient while stripping away resilience, care, and responsibility. It can innovate in ways that deepen dependency. It can scale systems of addiction, surveillance, speculation, and exhaustion. It can produce wealth while hollowing out the conditions of human flourishing. When that happens, the technical questions keep functioning. The metrics may even improve. But the deeper question has been buried. What is this economy for?

\section*{Capital That Serves Life}

This book is organized around a distinction that older moral traditions understood more clearly than we do: wealth can serve life, or it can feed on it. Capital can serve life when it helps human beings produce what is genuinely needed. It can fund useful work, sustain households, build institutions, support invention, steward land, preserve knowledge, create infrastructure, and strengthen communities. It can help people coordinate their gifts across time, allowing one generation to build what another generation receives.

But capital can also feed on life. It can profit from desperation, dependency, enclosure, scarcity, monopoly, opacity, debt, and domination. It can extract from workers, founders, tenants, borrowers, suppliers, ecosystems, public institutions, and future generations. It can turn productive enterprises into financial instruments, demand liquidity where patience is needed, and claim infinite upside where limits are morally required.

The central claim of this book is not that commerce is bad, profit is immoral, or modern finance should simply be rejected. That would be too easy, and it would be false. The central claim is that we have lost the moral grammar needed to distinguish between capital that serves life and capital that feeds on it.

\section*{A Normative Book in Historical Dress}

This book tells a historical story, but it is not only a history book. It follows a thread across more than two thousand years of thought: the struggle to distinguish useful, life-serving economic activity from accumulation detached from real need, real contribution, and real limits.

The thread begins with Aristotle, who distinguished oikonomia, the ordering of the household toward life, from chrematistics, the art of accumulation for its own sake. Christianity inherited this concern and deepened it. The Scholastics asked when prices were just, when merchants were legitimate, when risk could be compensated, when lending became usury, and when profit corresponded to real service. The Franciscans sharpened this vocabulary in their encounter with poverty, urban markets, merchants, and credit.

Then came the great separation. Economics gradually broke away from theology, moral philosophy, and older accounts of the good. This brought real gains: better tools for understanding prices, incentives, scarcity, coordination, and unintended consequences. But it also carried a cost. As economics became more technical and self-consciously scientific, it became less able to ask openly what economic life is for. The result was not neutrality. It was hidden normativity. Growth, efficiency, shareholder value, and return on capital began to function as moral standards without admitting that they were moral standards.

When a society loses the ability to ask what capital is for, capital does not become neutral. It becomes sovereign.

\section*{How to Read This Book}

This book asks for patience from several kinds of readers. Economists may find some of its moral and theological language unfamiliar. Theologians may find some of its financial detail technical. Entrepreneurs and investors may find parts of the critique uncomfortable. The discomfort is part of the point. Moral claims about economic life become real only when translated into governance, contracts, ownership forms, accounting, and law.

Above all, this book should be read as an act of reconstruction. The past cannot save us. But it can remind us that the present is not inevitable. Our current arrangements are not simply the economy. They are historically produced moral, legal, financial, and technological choices. They can be redesigned. They can be judged. They can be converted.

\chapter*{Chapter 1: The Question We Stopped Asking}
\addcontentsline{toc}{chapter}{Chapter 1: The Question We Stopped Asking}

Every economy answers a question, whether it admits as much or not.

The question is not, at first, a technical one. It has nothing to do with growth rates, with the efficiency of capital movement, or with how many acronyms a policy document can absorb before the reader loses the will to live. Those questions matter, but they are second-order. They become intelligible only after something deeper has been settled.

What is an economy for?

There is nothing sentimental about the question, and no question of hanging moral ornaments on a neutral machine like tinsel on a combine harvester. Every economic system answers it already, through its laws and institutions, through the incentives it sets and the habits it forms, and, most revealingly, through the things it never quite gets round to mentioning.

Goods are only a small part of what an economy moves around. It also forms people, trains desire, and teaches a whole society what success ought to look like. It decides whose labour is to be noticed and whose may be politely ignored, which harms can be externalised onto someone else and which dependencies make for a workable business model. It decides what is sacred and what may be liquidated, and what may be described, with a straight face, as ``efficiency.''

An economy is never merely a machine for producing goods. It is a moral order.

That moral order can teach patience or impatience, reward stewardship or extraction, build trust or normalise suspicion. Land may appear in its accounts as inheritance and shared responsibility, or as an asset class to be speculated upon. A founder may learn that a company is a community of work, or a lottery ticket waiting to be flipped. An investor may learn that capital exists to serve productive life, or that life itself is raw material for capital's return.

When we decline to ask what an economy is for, we do not escape moral judgment; we outsource it. Growth, shareholder value, liquidity, exit multiples, return on capital: these step forward to answer in our place, all the while pretending to be mere measurements.

But a default answer is still an answer.

\section*{Economy as Ordering}

The word ``economy'' comes from the ordering of the household. Before it referred to national output, fiscal policy, abstract systems of exchange, or graphs that rise cheerfully upward while the planet catches fire in the background, it meant the management of a real human community: how it fed and sheltered itself, how labour was divided, how land was handed down, how care was given to those who needed it, and how all of this carried on across the generations.

This older meaning is easy to romanticise, and the temptation should be resisted with both hands. The ancient household was not a cosy domestic idyll waiting patiently for modernity to rediscover it. Households have always been places of hierarchy and exclusion, of coercion and of work that went unpaid because nobody thought to count it. The ancient oikos is no more a model to recover wholesale than a medieval drainage ditch is a model for urban sanitation.

Still, the older meaning preserves something worth holding onto. Economy is not merely exchange. It is order.

To order something is to arrange it toward an end. A hospital is ordered toward healing, a school toward education, a court toward justice, a craft toward making something useful or beautiful. These ends are not decorative extras: they are what makes the institution intelligible in the first place, and what distinguishes a functioning one from a merely busy one. A hospital that turns a handsome profit while failing to heal anyone has failed as a hospital, however solvent. A school can have excellent branding and still fail to educate, in which case it has failed as a school. The same goes for the court that processes cases efficiently without doing justice. Busyness is not purpose.

So what is an economy ordered toward?

Modern societies answer this question sideways, as if trying not to make eye contact with it. The economy, they say, is for growth, or for efficiency, or for satisfying preferences, or for innovation, or for national competitiveness, or for investor returns. Each answer contains a partial truth. An economy that produces too little cannot sustain human life; one that suffocates enterprise becomes stagnant and brittle in its own way. But no partial truth can bear the full weight of the question. Growth in what? Efficiency for whom? Innovation toward what end, and prosperity measured by whose account of the good?

The old language of economy forces these questions into view because it refuses to sever economic activity from the ends of human life. Economic systems are not judged solely by how much they produce, but by what they make possible: whether they help people live well, whether they sustain those who cannot fend for themselves, whether they protect the conditions of trust, and whether they make space for care and friendship, for craft, and for the stubbornly unprofitable forms of life that somehow make life worth living. The question is not being imposed on economics from outside by moralists who have wandered in through the wrong door. Every economy has already answered it, wisely or foolishly, consciously or otherwise.

\section*{The Disappearance of Teleology}

To ask what an economy is for is to ask a teleological question: a question about ends, purposes, and the direction in which activity is aimed. Teleology was once the ordinary grammar of practical reasoning. A knife is for cutting; a bridge is for crossing; a physician's art is for healing; a political community exists for the common good. To understand what something is, you must understand what it is for.

Modern economics has grown uncomfortable with this kind of language. It is comfortable instead with description: of what people choose, of what they prefer, of how they respond when incentives shift, of how prices coordinate information that is dispersed across millions of strangers. These are genuine and important forms of knowledge. Anyone who has watched a central planning bureaucracy try to anticipate the price of onions should be properly grateful for the informational power of markets.

But economics becomes thinner when it refuses to ask whether preferences are rightly ordered to begin with, whether choices are being made under conditions of domination, whether the incentives in play corrupt the very goods they are supposed to organise, or whether the system as a whole is serving human flourishing or sawing through the branch on which it sits.

The disappearance of teleology does not produce neutrality. It produces hidden normativity.

A society that declares ``the economy is for growth'' has already made a claim about ends. So has the company charter that says the firm exists to maximise shareholder value, the fund that defines success as exit value, the platform that runs on engagement as its governing metric. We are never really choosing whether to have ends; we are only choosing whether to name them honestly and subject them to scrutiny, or to allow them to govern us from behind a curtain while pretending to be spreadsheets.

Unnamed ends become more powerful, not less. They govern institutions in the costume of neutral description, and reshape moral imagination in the costume of performance metrics. A means becomes a master, without anyone really noticing.

Money is useful because it enables exchange; that is its proper work. It becomes formative, and dangerous, when it turns into the measure of all things, training us to see the world through the lens of price. We start to confuse what is valuable with what is monetisable, what is productive with what is merely profitable, what is worthwhile with whatever happens to scale. Capital is useful too, because it lets human beings coordinate work across time, bridging the gap between planting and harvest, between research and discovery, between founding a thing and seeing it become stable. But it becomes disordering when it detaches from the purposes it was meant to serve and begins to demand that everything else become liquid, measurable, ownable, subordinate. This is the inversion of means and ends. Older moral traditions recognised it as one of the most persistent dangers in economic life. Modern societies have not escaped the danger; we have only developed better software for it.

\section*{Markets as Moral Systems}

Markets are often described as mechanisms, and this is not wrong. They coordinate information and match supply with demand. They allow decentralised actors to cooperate without ever needing to know one another personally, and they reveal needs and possibilities that no central authority could fully anticipate, even one equipped with excellent intentions, a large staff, and very expensive consultants. The Austrian tradition was right to insist on the epistemic power of market order. That insight belongs within any serious account of moral economy.

But markets are never only mechanisms. They are moral systems, because they depend on moral conditions they cannot generate by themselves.

A market does not exist before law and morality, like a naturally occurring pond into which governments later dump regulations. It is constituted by law and morality from the very beginning. Every market rests on prior judgments: about what may be owned and what may be sold, about what counts as fraud, about who may enter into contracts and on what terms, about how liability is to be assigned, and about what must remain a public good beyond private enclosure. None of these are technical questions with purely technical answers. They are moral and political choices, expressing particular judgments about fairness, dignity, obligation, and the limits of commodification.

This matters because the moral quality of a market cannot be determined just by asking whether its participants chose freely within it. Choices are always made within structures, and structures can make certain choices appear as the only option when they are, morally speaking, a form of compulsion. A borrower ``chooses'' a predatory loan because the alternative is eviction. A worker takes degrading labour because the alternative is hunger; a founder accepts extractive investment because the alternative is watching the company die. Formally, these are choices. Morally, they may be signs of domination.

Freedom is not just the absence of visible force. It depends also on the presence of real alternatives, on background conditions that are minimally just, and on institutions that stop another person's necessity from becoming someone else's business model.

This is also why moral economy cannot be reduced to individual virtue. Decent people inside extractive structures face extractive pressures. An investor may want to be patient and yet find herself running a fund whose terms demand liquidity and an exit. A manager may care deeply about workers and still operate under governance rules that make short-term financial return the only legitimate measure of performance. A platform that begins as a genuinely useful network can, over time, become structurally incentivised to enclose, to surveil, to dominate. The architecture matters. Who owns and who controls, who bears the risk and who receives the surplus, who is protected by law and who remains invisible to the balance sheet: these are not decorative questions. They are the substance of moral economy.

\section*{Profit, Purpose, and Justification}

One of the easiest mistakes in moral economy is to treat profit as either self-evidently good or inherently suspect. Both positions have the advantage of being simple, which is always tempting. Both are wrong.

Profit can be a sign that real value has been created and recognised. A baker who sells bread, a carpenter who restores a home, an investor who patiently funds productive enterprise: each may earn profit in ways that correspond to genuine service. Profit can sustain continuity and reward skill, compensate for risk, fund reinvestment, and make useful work durable over time. This is not trivial. Bread does not bake itself because the baker has a noble heart. Tools wear out, and workers need wages. A moral economy that cannot account for these realities will soon find itself very morally pure and very short of bread.

But profit can also come from extraction: from monopoly power and manufactured dependency, from regulatory arbitrage and asset stripping, from debt traps and speculation on scarcity, from the systematic exploitation of people who happen to have less bargaining power than the person across the table. It can be generated by making other people, or institutions, or ecosystems more fragile, by shifting risk downward while moving surplus upward, by rewarding those who control chokepoints rather than those who create genuine value.

So the question is not whether there is profit. The better questions are more precise, and more uncomfortable. Profit from what activity, and at whose expense? In proportion to what contribution, what risk, what responsibility? Reinvested into what kind of future?

A moral economy does not deny the legitimacy of gain. It asks gain to give an account of itself.

This is why older debates about usury, just price, and merchant legitimacy remain more contemporary than they first appear. They were not primitive attempts to regulate finance before modern theory arrived wearing a better suit. They were attempts to ask whether monetary gain corresponded to a real contribution or came at the expense of someone else's vulnerability. That question has not dissolved with time. It has merely changed its outfit, reappearing in payday lending and sovereign debt, in private equity roll-ups and housing speculation, in platform monopolies and the enclosure of open-source labour by proprietary infrastructure. A society without a moral vocabulary for profit cannot distinguish enterprise from extraction; a society that cannot make that distinction will struggle to notice when it is being consumed by the very system on which it depends.

\section*{Capital as Servant or Master}

Capital is stored capacity, the ability to do something across time. It can take the form of money or tools, land or machines, code or knowledge, networks or institutions, or even trust accumulated through long practice. It lets human beings coordinate effort beyond the immediate present, bridging the time between planting and harvest, between research and discovery, between founding a thing and seeing it stand.

This is why capital is powerful. And this is why capital must be ordered.

When capital serves, it remains accountable to the good it makes possible. It supports productive activity and participates in risk; it accepts proportionate limits and receives return in proportion to genuine contribution. It respects the integrity of the enterprise, community, or commons it funds. When capital rules, the logic inverts. Work becomes a cost to be minimised, housing becomes yield, companies become portfolios. Mission becomes branding, and human beings become inputs to be measured by what they can return on investment.

The difference is not always visible from outside. Two companies might sell the same product, two funds might use identical language about impact and purpose, two platforms might provide apparently similar services, and yet one of each pair may be structured to sustain life while the other is structured to extract from it. This is why moral economy must attend carefully to design: to the actual architecture of ownership, governance, and obligation that determines what a firm or fund can and cannot do, regardless of stated intentions. Good intentions are lovely things. They are also, when unsupported by structure, distressingly biodegradable.

Capital becomes servant or master through structure. That is the central practical insight of this book.

\section*{The Hidden Moral Education of Economic Life}

Economic systems do not merely allocate resources. They educate desire.

A society that rewards extraction teaches people to become extractive. Founders learn to optimise for valuation rather than durability. The same logic teaches investors to seek control without continuing responsibility, encourages workers to think of themselves as personal brands, and trains institutions to treat their mission as a communications strategy. Over time, people internalise the system's assumptions, and start to see themselves and each other through the categories the economy provides.

If the economy assumes that human beings are primarily self-interested maximisers, institutions will be built around that assumption, and those institutions will train people to behave more like self-interested maximisers. What begins as a simplifying model becomes a social pedagogy. Our models do not merely describe us. They recruit us.

The same logic can run in the other direction. Institutions can teach stewardship: patience and fidelity, the discipline of craft, responsibility toward those who come after. They can protect productive communities from liquidation, honour forms of value that price cannot capture, and make it possible to keep promises across generations. They can build the conditions in which people are more likely to be good, rather than merely more efficient.

This is why the question ``what is an economy for?'' cannot be separated from the question ``what kind of people are we becoming?'' An economy ordered toward human flourishing must care not only about outputs but about formation. It must ask what habits, dependencies, and desires its structures are producing, and whether people are being trained into freedom or into mere appetite, into responsibility or into evasion, into stewardship or into extraction. That is not moralism. It is realism about what human beings actually are: creatures shaped by practices and incentives, by institutions and dependencies, and not least by what they love.

\section*{The Cost of Forgetting}

When societies stop asking what an economy is for, they do not become more practical. They become less able to recognise their own disorder.

They notice the symptoms well enough. Inequality and loneliness, ecological destruction and housing crises, debt burdens and institutional fragility, political anger and declining trust, monopolistic concentration, burnout, the erosion of local resilience. These are not subtle developments. They are the kind of symptoms that would cause any sensible physician to stop talking about minor lifestyle adjustments and order a proper examination.

But without the deeper question, the symptoms appear disconnected. Each becomes a discrete policy problem, or a market failure, or a technological shortcoming, or a management issue. Some of them are exactly those things. They are also signs of something more fundamental, however: economic life detached from any shared account of what it is for. The economy produces more, but not always what is needed. It connects more people without always deepening anything between them. It innovates faster without always innovating toward wisdom. And it creates more liquidity, sometimes by dissolving the very bonds that make life livable.

Without the deeper question, both critique and reform stay shallow. We argue about distribution without asking what is being distributed, about growth without ever asking what is growing, about freedom without checking whether people are in fact free from domination. An economy can be full of activity and still be profoundly disordered. The question ``what is this economy for?'' is therefore not a luxury for the philosophically inclined. It is diagnostic. It reveals whether the parts of economic life are rightly ordered to the whole, or whether the whole has, while no one was looking, been subordinated to one of its parts.

\section*{Recovering the Question of the Good}

To recover the question of the good is not to imagine that everyone in a pluralist society will agree on ultimate ends. They will not. Anyone who has spent ten minutes at a neighbourhood association meeting knows this, and that is before we get to religion, metaphysics, political economy, or the contentious question of whose turn it is to maintain the shared hedge.

Modern societies contain many religions and moral traditions, many philosophies, many competing visions of human flourishing. Any serious recovery of moral economy must reckon with that plurality honestly. It cannot pretend that a single comprehensive doctrine can be imposed on economic life without doing violence to the freedom and complexity of common life.

But pluralism does not require silence about the good. It requires disciplined public reasoning about the goods that make common life possible at all: dignity and justice, trust and stewardship, freedom from domination, care for those who cannot fend for themselves, ecological responsibility, the chance to do meaningful work, and the capacity of communities to shape their own futures. These goods are contested at the margins. They will not make politics tidy or economics simple, nor will they make human beings notably less exasperating. But they are not empty. Without them, economic freedom becomes nothing more than the freedom of the powerful to set the terms for everyone else.

The task, then, is not to impose a single comprehensive doctrine on economic life. It is to recover enough moral language to ask whether our economic institutions are serving life, and enough institutional imagination to redesign them when they are not.

That is where the rest of this book begins. The chapters that follow turn to the traditions that first gave Western moral economy its grammar. Aristotle will not offer complete answers; his world was not ours, and many of his assumptions must be firmly set aside. But he saw something that remains decisive: economic activity can be ordered toward life, or it can become the art of accumulation without limit. That distinction is the beginning of moral economy. To understand capital that serves life, we must first understand why the ancients feared wealth that had forgotten what it was for.

\chapter*{Chapter 2: Aristotle and the Birth of Moral Economy}
\addcontentsline{toc}{chapter}{Chapter 2: Aristotle and the Birth of Moral Economy}

We begin with Aristotle. He does not give us a complete answer, but he asks the right first question.

His world cannot be ours. It was deeply hierarchical, dependent on slavery, organised around assumptions about gender and citizenship that must be rejected outright. There is no innocent classical past waiting patiently for us to restore it, like a marble statue under a sheet. The marble, if we are honest, has some very serious cracks.

But Aristotle saw something that modern economic life often struggles to name: acquisition is not one thing. Wealth-getting can be ordered toward life, or it can become detached from it. Money may remain a useful instrument, but it can also become an object of desire in itself. And exchange, while it can meet real needs, has a tendency to turn into a technique for converting money into more money, without reference to anything beyond itself.

That distinction, between wealth that serves life and wealth that feeds on it, is the founding insight of moral economy.

Aristotle did not ask, in the first instance, how wealth could be maximised. He asked where wealth belonged within a rightly ordered life, and he placed economic activity inside ethics and politics, not alongside them in a separate department with its own jargon and a suspicious number of charts. The question of wealth was inseparable from the question of virtue, the question of property from the question of the household, and the question of exchange from the question of the city. In Aristotle, economy is not yet ``the economy.'' It is part of the art of living well.

That is why he matters.

\section*{Oikonomia and Chrematistics}

Aristotle's most important distinction for our purposes is between oikonomia and chrematistics.

Oikonomia, household management, is the ordering of a real human community toward the needs of life: how it feeds and shelters itself, how labour is divided, how land is handed down, how care is given, and how all of it carries on across the generations. The household is not the highest form of human association for Aristotle; the polis is higher. But the household is necessary. Human beings cannot pursue virtue, friendship, contemplation, or political life if the basic conditions of survival are not first sustained. It is difficult to deliberate nobly about justice while starving, freezing, or trying to repair the roof with no tools and no bread.

Economic activity therefore begins as a servant. It is limited by the needs it serves, and those needs have a natural boundary. There is such a thing as enough food, enough tools, enough land for a decent and virtuous life.

Chrematistics is something different. The word refers to the art of wealth-getting, of which Aristotle identifies two forms. The first is natural: it acquires what life genuinely requires. The second is unnatural: it seeks wealth without limit, without reference to any need beyond accumulation itself.

Here Aristotle identifies a danger that will run through every chapter of this book. When wealth is pursued as an end rather than a means, desire becomes unlimited. Ordinary needs are bounded. Hunger is satisfied by food, the need for shelter by a roof, the need for shoes by a pair of shoes (though modern retail has done heroic work to complicate this last one). The desire for money, treated as an end in itself, knows no such satisfaction. More can always be wanted and always justified, by any number of names: prudence, opportunity, growth, ambition, fiduciary duty, or simply ``what the market requires.'' The activity becomes infinite because the end has been severed from any finite good.

This is the beginning of the distinction between capital that serves life and capital that feeds on it. Aristotle does not use that language, but the structure of the distinction is already fully present. Crucially, this is not a critique of material goods as such. Aristotle is no ascetic. He knows that human life requires external goods, that virtue needs a material context, and that poverty can deform human possibility as surely as excess can. The problem is not wealth itself. The problem is wealth detached from its proper end.

\section*{The Household, the City, and the Good Life}

For Aristotle, the household exists for the sake of life. The city exists for the sake of the good life.

This is among the most important teleological claims in the history of political thought. Human beings do not associate merely to survive. They associate in order to live well, in order to pursue justice and virtue and friendship, to deliberate together about the common good. The city is not just a defensive alliance, nor an administrative convenience, nor a tax jurisdiction with ambitions. It is the community in which human excellence becomes fully possible.

Economic activity must therefore be subordinated to politics, and politics to the good life. This hierarchy strikes most modern readers as almost entirely inverted. We are accustomed to governments judged by growth rates, by investment confidence and market stability. Our cities reorganise themselves around real estate values and the attraction of capital. Universities are repositioned as innovation pipelines, housing is treated as an asset class, and human life itself, implicitly and often explicitly, is asked to justify its existence in economic terms.

Aristotle would call this disorder, not prosperity: a civilisation that has put the furniture on the ceiling and congratulated itself on the novelty of the arrangement.

Wealth, for Aristotle, is several steps removed from the highest human ends. It is necessary but subordinate, useful but not ultimate. And because it is powerful but not ultimate, it is dangerous. Money can deform desire and turn means into ends. It can tempt the soul toward excess and train people to measure success by accumulation, rather than by the quality of their character or the depth of their contribution to the common good.

This is why Aristotle's economic thought cannot be separated from his ethics. Virtue is the formation of desire according to reason: a virtuous person desires the right things, and desires them in the right way and for the right reasons. Every craft and discipline is intelligible because it has an end. Medicine is for health, shipbuilding for ships, strategy for victory. When an activity loses sight of its end, it becomes disordered. An economy that has forgotten what it is for does not thereby become neutral. It becomes disordered without knowing it.

\section*{Use Value and Exchange Value}

Aristotle also offers an early version of the distinction between use value and exchange value. This distinction will travel through centuries of moral and economic thought, reappearing in Marx, in institutional economics, in debates about financialisation, and in contemporary arguments about platform monopolies and the digital commons.

Every object, Aristotle observes, can be used in more than one way. A shoe can be worn, which is its proper use. But it can also be exchanged. This is not necessarily wrong. Exchange allows households to meet needs they cannot satisfy on their own. If one household has surplus grain and another has surplus oil, trade can serve both. Commerce can be a form of cooperation.

But exchange also introduces a danger. Once goods are viewed primarily through their exchangeability, their relation to concrete use becomes secondary. The shoe is no longer seen first as something made to protect a foot; it is seen as something with a tradeable value. As exchange develops and money enters as a common measure, this abstraction deepens. A shoe still has a natural use, grain still feeds, a house still shelters. But money has no natural limit in the same way. Its use is precisely to be exchangeable, to stand in for anything. And precisely because it can represent any good, it can tempt human beings to desire it without reference to any particular good.

Money can become desire made abstract.

This is why Aristotle worries about forms of wealth-getting that operate through exchange purely for the sake of monetary increase. When exchange serves use, it remains connected to life; when it serves only accumulation, it detaches from life's natural purposes. The disorder begins when things cease to be judged by the ends they serve, and are judged instead by their capacity to generate more exchangeable value. Housing then becomes a speculative asset rather than shelter, a company an exit vehicle rather than a community of production, land an opportunity rather than a place to live and pass on. When exchange value begins to crowd out every other form of valuation, Aristotle's warning becomes contemporary.

\section*{Why Aristotle Feared Money Detached from Reality}

Aristotle's suspicion of money is not a rejection of its usefulness. It is a fear of money's power to detach wealth from reality.

Natural wealth, food and tools and land and household goods, consists in things needed for life. These things are finite because the needs they serve are finite. Money is different. It is artificial; it exists by convention; it has value because a community treats it as valuable. Precisely for that reason, it can become dangerously unstable. Money can represent real goods and real work, real trust and real obligation. But it can also become detached from all of these, a claim floating free of the productive reality it was meant to serve.

This is why Aristotle condemns usury. When money breeds money, when a lender profits simply by virtue of having money to lend, without risk or labour or productive contribution, the unnatural structure of chrematistics becomes most visible. The Greek word for interest is tokos, meaning offspring, and Aristotle finds the metaphor morally revealing. Living things reproduce according to nature; money does not grow wheat or bear children or heal bodies or build houses or sustain cities. To make money generate more money simply by being money is, for Aristotle, the purest form of unnatural accumulation.

Modern readers will object, rightly, that capital can be productive. Money invested in tools or training, in infrastructure or research, can enable real growth. Interest can compensate for risk, opportunity cost, and administrative effort. Later Franciscan and Scholastic thinkers will work through precisely these questions with considerable sophistication. Those later developments matter, and Aristotle is not the final word on any of this. But his warning retains its force because it identifies a permanent structural danger: monetary claims can become detached from real contribution. A financial claim may continue extracting long after any genuine contribution has been repaid. A lender may profit from another person's desperation, an investor demand control disproportionate to risk, a platform monetise a network it did not create alone, a landlord capture value produced by public infrastructure and community life. In all these cases, the Aristotelian concern returns in modern form: money has become detached from the living realities it was meant to serve.

\section*{Virtue and the Discipline of Desire}

Aristotle's economic thought is inseparable from his account of virtue, because the deepest problem of accumulation is not technical. It concerns the formation and deformation of desire.

Virtue, for Aristotle, is not mere rule-following. It is the habituated capacity to desire and choose well, to love what is genuinely good in the right measure and for the right reasons. The virtuous person is not someone who reluctantly does good while inwardly longing for vice, and glancing wistfully at the cake, the crown, or the quarterly bonus. Virtue is what happens when the soul has been well trained, through practice and community and reason, so that its desires are rightly ordered rather than merely restrained.

Economic life shapes this training whether it intends to or not. A society that honours moderation, generosity, and public service forms different desires than one that honours accumulation, status display, and domination. The household that teaches sufficiency forms different children from one that teaches endless comparison. And the same goes for markets: one that rewards useful work forms different entrepreneurs from one that rewards manipulation, enclosure, and the exploitation of captive users.

This is why Aristotle is not satisfied to ask whether people are free to choose. He asks what kind of people they are becoming, and whether the economy is making them better or worse at the art of living well. A disordered economy does not merely distribute goods unjustly. It forms disordered souls. It trains desire away from the good and rewards the restless and acquisitive over the patient and excellent, until sufficiency becomes progressively harder to recognise as a virtue rather than a failure.

In this respect, Aristotle's critique runs deeper than many modern critiques of capitalism. He is not only asking whether wealth is distributed fairly. He is asking whether the pursuit of wealth deforms the soul and the city; whether it trains us to want the wrong things and measure ourselves by the wrong standard. Every economy already answers that question, through its incentives and its rewards and its silences. The only question is whether it answers wisely or destructively.

\section*{Property, Use, and the Common Good}

Aristotle's treatment of property is another place where he resists easy classification. He is not a communist in any modern sense. He criticises Plato's proposal to abolish private property, arguing that private ownership encourages responsibility, enables generosity, and supports household stability: people tend to take better care of what is theirs, and one cannot give generously without first having something to give.

But Aristotle does not treat property as absolute. Property may be privately held, he argues, but its use should serve the common good. Ownership carries moral expectations; it is not merely a right of exclusion. Wealth should be used in ways that support virtue, friendship, and the public life of the city. The propertied person is not simply a possessor, but a steward, accountable to the community within which property has its meaning.

This distinction between ownership and use will become enormously important in later thought: in Catholic Social Teaching, in Protestant stewardship traditions, and in contemporary debates about mission-locked governance and fiduciary duty. It helps us see that the moral question is not simply whether property is private or collective. Private property can serve life when it supports responsibility, continuity, and productive use; it becomes extractive when it enables domination, manufactured scarcity, and the conversion of shared goods into private power. Collective property can protect commons and democratic stewardship; it can also become destructive, dissolving responsibility and concentrating power without accountability. The legal title does not settle the moral question. Who owns matters, but so do the obligations attached to ownership, the limits restraining its use, and the forms of life it makes possible or forecloses. Aristotle gives us a vocabulary for asking these questions without reducing everything to either absolute private accumulation or state control. Property is a moral relationship within a community ordered toward the good life.

\section*{The Problem of Unnatural Accumulation}

At the centre of Aristotle's economic thought is the problem of accumulation without limit.

Natural acquisition is bounded because it seeks what is sufficient for life. Unnatural accumulation is unlimited because it treats wealth as an end in itself; and a means elevated to an end loses the natural boundary that an end would impose. A person who wants food can be satisfied by eating. The person who seeks money as such can always seek more, because money is not a final good. It is a universal means, and universal means have no built-in stopping point.

This needs careful handling. It would be too simple, and wrong, to say that all growth is unnatural or that all accumulation is disordered. Households need savings, firms need reserves, communities need infrastructure. Innovation often requires patient investment long before any return appears. Any number of institutions, from monasteries and family businesses to cooperatives, schools, research labs, and evergreen funds, need some capacity to endure bad seasons. The question is not whether accumulation exists, but whether it is ordered. Whether, in other words, it remains accountable to the purposes that justified it.

Accumulation becomes morally dangerous when sufficiency disappears as a category, when limits start to be treated as failure, and when every good is subordinated to the expansion of financial claim. A company may be profitable, stable, and genuinely useful, and still be judged inadequate because it cannot grow fast enough. A founder may want to build something durable, yet find herself pressured toward exit because investors require liquidity. The same logic distorts housing: a city may have enough homes for shelter in principle, but the homes have become speculative vehicles, and so they are unaffordable. Without the category of sufficiency, every limit appears as an obstacle. Reserves become underutilised capital, commons become untapped assets, and every human vulnerability becomes a market opportunity. That is not prosperity; it is disorder that has learned to call itself growth.

\section*{The Limits of Aristotle}

Precisely because Aristotle is useful, we must be honest about his limits.

His moral economy is built on a social world we cannot accept. The household he describes includes relations of rule that subordinate women, children, and enslaved people. His vision of citizenship excludes most of those whose labour makes the city possible. His leisure rests on hierarchy, and his account of the good life is profound but reserved, in practice, for a privileged few. A thinker who could theorise the common good while accommodating slavery has already revealed something important about the gap between moral vocabulary and moral practice.

We should not romanticise premodern moral economy on the back of this. A society can ask teleological questions and still answer them unjustly; it can speak of virtue while excluding most people from any share in it. But Aristotle's own framework contains the resources for criticising his exclusions. If the city exists for the good life, we must ask who is permitted to share in that good. So too with economic life: if it should serve human flourishing, whose flourishing counts? And if wealth should support virtue and community, then the question is whether workers and women, migrants and the poor are being treated as participants in the common good, or merely as instruments for someone else's.

A modern recovery of moral economy must universalise the question Aristotle asked but failed to universalise in practice. An economy ordered toward life must be accountable to the flourishing of all who participate in it and all who are affected by it: the vulnerable and the excluded, future generations, the natural world on which all economic activity ultimately depends. Aristotle gives us the grammar of ends. We must widen, dramatically, the circle of those whose ends matter. That is not a rejection of his importance. It is the only way to receive him honestly.

\section*{From Aristotle to Moral Economy}

Aristotle's enduring contribution is that he refuses to let wealth define its own purpose. Economic activity must be judged by a standard outside itself: the good life of human beings in community. He distinguishes acquisition ordered toward real need from accumulation without limit. He shows that money can become detached from the reality it represents, and that exchange, however useful, becomes dangerous when it overwhelms use. He shows that property carries moral obligations. And he shows that the inversion of means into ends is one of the oldest and most persistent dangers in human civilisation.

These insights do not solve modern problems directly. Aristotle has nothing to say about how to design an evergreen fund or govern open-source infrastructure, or how to regulate a platform monopoly or structure repayment caps. This is understandable. His experience with software governance was, by all available evidence, limited. But he gives us the first question we need in order to approach those problems with any moral seriousness: what is this economic activity for?

Whenever capital supports real production, shared flourishing, and institutional durability, it participates in oikonomia properly understood. When it seeks accumulation without limit, extracts from dependency, or detaches monetary return from real contribution, it moves toward chrematistics. A housing market that treats shelter as secondary to yield is chrematistic; so is a fund that cannot recognise enough. The opposite pattern includes a lender who bears genuine risk and accepts proportionate limits, a company that serves real needs while sustaining its workers and community, a form of capital that helps human beings coordinate their gifts across time. These participate in economy. The terms are ancient. The pattern is thoroughly contemporary.

This is the distinction we will follow through the centuries ahead.

\section*{The Question Aristotle Leaves Behind}

Aristotle leaves us with a question that is both ancient and urgent: how can the necessary pursuit of material provision be kept subordinate to the good life?

It is a question that every subsequent tradition in this book will wrestle with in its own way. Christianity will deepen the stakes by placing wealth under the judgment of charity, mercy, and salvation, so that avarice becomes not merely a civic nuisance but a spiritual emergency. The Scholastics will ask how commerce can be morally legitimate without becoming exploitative, a question that turns out to be considerably harder than it looks. The Franciscans will discover, perhaps to their surprise, that their commitment to poverty forced them to think more carefully about money than almost anyone else. Salamanca will develop sophisticated accounts of value and market coordination, while wrestling with the uncomfortable fact that the same reasoning travelled rather comfortably alongside empire. And so on, through Protestant vocation and Adam Smith and the long unravelling of moral economy into technical discipline.

All of these traditions are wrestling with a problem Aristotle named first. He teaches us to ask whether wealth is serving life or slowly replacing it; to distrust any accumulation that has no resting point; to recognise that economic life must, in the end, be judged by the kind of human beings and communities it forms.

He does not give us the whole moral economy.

He gives us its beginning.

\chapter*{Chapter 3: Christianity Encounters Commerce}
\addcontentsline{toc}{chapter}{Chapter 3: Christianity Encounters Commerce}

Aristotle's moral vocabulary for economic life was powerful, but it was also bounded. He wrote within the horizon of the Greek city, a world of households and citizens and slaves, where exchange was largely local and the unit of analysis was the polis. He had no framework for the great sweep of empire or for long-distance trade across languages and cultures, no framework for debt instruments that stretched across generations, nor for the spiritual weight a new religion would shortly place on the souls of the wealthy and the claims of the poor.

That expansion of the moral field came with Christianity.

What Christianity added to Aristotle's framework was not primarily new economic analysis. It did not arrive with a better theory of price formation, a revised account of trade flows, or a handy appendix on medieval liquidity management. Its contribution was more radical. Christianity transformed the stakes. Wealth became not only a civic question but a question of salvation, and the poor were no longer merely a social fact but a theological claim upon the conscience of every person who had more than they needed. Debt carried the weight of forgiveness, property was now held in stewardship before God, and commerce had become a site of moral formation in which one could grow closer to, or further from, the life one was called to live.

The entire frame of economic judgment shifted. And it would never entirely shift back.

\section*{Roman Law and Early Christianity}

The Roman world into which Christianity was born had a highly developed legal culture around property, contract, and debt. Roman law gave commerce a reliable structure. Ownership was well defined, contracts were enforceable, and disputes between parties could be adjudicated with considerable sophistication. In many respects, this was genuinely valuable. Predictable legal order makes cooperation possible: merchants and lenders and buyers all need to know that agreements will hold. Commerce without law is not freedom. It is usually a queue forming outside the office of the strongest person in town.

But Roman law also assumed and reinforced the hierarchies of Roman society. Ownership was bound up with status, and debt could be enforced through bondage. The power relations embedded in commerce were largely taken for granted by the legal order that regulated them. The law could tell you whether a contract was valid; it could not, on its own, tell you whether the relationship the contract created was just.

Early Christian thinkers recognised this gap and pressed hard against it. Legal validity was not moral sufficiency. A transaction could be enforceable and still be an occasion of sin; a contract could be signed freely and still create a relationship of domination. Figures such as Ambrose of Milan, John Chrysostom, and Basil of Caesarea wrote with striking directness about the moral obligations of wealth. Chrysostom, preaching in Antioch and Constantinople at the end of the fourth century, did not speak gently about wealthy Christians who ignored the poor at their doors. For him, the extra cloak in your wardrobe belonged to the person who was cold, not as a charming optional donation but as a debt of justice. Basil asked, with more sharpness than comfort generally allows, who the rich man is who does not give, and answered: a thief.

These were not marginal cranks shouting from the edge of Christian life. They were bishops speaking from its centres. The early Church did not reject commerce or property outright. It insisted that both were morally dangerous and morally accountable: that wealth was a gift held in trust, that labour had a dignity beyond its market price, and that every transaction occurred within a moral order law alone could not sustain. The distinction between what is legally permitted and what is morally ordered would become one of Christianity's lasting contributions to moral economy.

\section*{Monasticism and the Economics of Enough}

If early bishops provided the theology, monasticism provided something rarer: a working experiment in economy ordered toward life rather than accumulation.

Monastic communities of every description, the Benedictines and Cistercians, the Franciscans and Dominicans, and many others across the medieval centuries, organised production and labour, land and surplus and distribution according to purposes quite different from those of ordinary commercial life. The Rule of Benedict, written in the sixth century, is in part an extraordinarily practical document about almost every aspect of communal life: about work and time, about food and clothing and tools, about sleep and obedience and the welcoming of guests, about correction and the careful management of a community's resources. It is also, in the same breath, a document about the formation of desire. It is about learning to want the right things, to take satisfaction in enough, and to find in ordered common life something more durable than individual accumulation could offer.

This is one of the great surprises of monastic life. From the outside, it can appear to be a flight from economics. In reality, monasteries were often intense exercises in economic design. Someone had to grow the food and maintain the buildings, organise the labour, preserve the tools, welcome the travellers, and care for the sick. Someone had to copy the manuscripts, manage the land, store the grain, and decide, with impressive regularity, all manner of practical questions that would have been perfectly recognisable to any modern operations manager, had such a person been available and willing to take the tonsure.

Monks and nuns renounced personal property. But their communities did not renounce productivity. Cistercian monasteries became among the most efficient agricultural enterprises in medieval Europe, experimenting with crop rotation and water management and wool production. Benedictine houses preserved manuscripts and educated students, hosted travellers and cared for the sick. The monastic economy was not anti-economic. It demonstrated that production, coordination, and even surplus could exist within a moral framework: that wealth could be present without becoming sovereign, work disciplined without being exploitative, and surplus managed without being idolised. The question of what all this activity was for could remain a living question, rather than a forgotten one.

Monastic institutions also became important sources of financial innovation. Communal funds and charitable lending pools, structured forms of redistribution, developed within and around monasteries in ways that would later influence the Monti di Pietà and the Scholastic debates about credit and usury. The institutions of moral economy were not only ideas. They were built, sometimes awkwardly and sometimes brilliantly, often imperfectly, by people who had decided to take the question of enough seriously enough to organise their lives around it.

\section*{Usury and the Moral Grammar of Lending}

No economic question occupied early and medieval Christianity more persistently than usury, and the reasons tell us something important about what the tradition understood to be at stake.

The prohibition on interest, inherited from Jewish law, reinforced by patristic teaching, and codified in canon law, was not simply a misunderstanding of economics waiting for a later age to arrive with compound interest tables and an air of superiority. It was a moral diagnosis. When a lender charged for the use of money, the concern was not merely that money is sterile, though Aristotle's argument about money breeding money remained influential. The deeper concern was pastoral and structural: interest makes another person's necessity profitable. The greater the borrower's desperation, the more the lender can charge. A system built on interest can therefore become one that is structurally rewarded by the suffering of those who have no alternative. That is not a medieval fantasy. It describes payday lending today, the restructuring of sovereign debt, and the pricing of medical debt with uncomfortable accuracy.

What the tradition recognised, and what modern financial analysis often struggles to name, is that the moral problem with usury is not only the number on the loan. It is the relationship the loan creates. Does the return correspond to genuine risk, shared in partnership, or is it guaranteed extraction in another costume? Does the arrangement serve the borrower's need or exploit it? And does the debt remain proportionate to what was received, or compound into a claim that can never be discharged? A loan that restores the borrower to agency is one thing; a loan that converts misfortune into a recurring revenue stream is something else entirely.

Christian thinkers refined these concerns considerably over time, distinguishing compensation for genuine loss from gain extracted from necessity, separating consumption loans from productive partnerships, and telling legitimate risk-sharing apart from guaranteed extraction. These distinctions were not always clean, and they could be abused. Clever people have always been able to dress moral evasions in respectable clothing, and medieval Europe did not lack tailors. But the effort to distinguish mattered. The tradition was not trying merely to ban lending. It was trying to ensure that lending served life rather than feeding on it, a moral seriousness about credit that purely transactional frameworks cannot sustain.

\section*{Charity, Debt, and Obligation}

Christianity's contribution to moral economy was not only restrictive, not merely a set of prohibitions on avarice and usury. It was also constructive. It introduced a positive moral vocabulary for obligation and redistribution, for mercy, and for the duties that wealth creates.

Debt, in the Christian moral imagination, was never merely a legal relationship between two contracting parties. It had moral texture on both sides. Borrowers were responsible for repayment, but lenders were responsible for how they lent and on what terms, and for whether they were profiting from desperation rather than genuinely sharing risk. Wealth was not simply a possession but a form of accountability. The parable of the talents, read in this tradition, was not primarily about investment returns, however much modern business sermons have tried to make it sound like a divine endorsement of portfolio management. It was about whether what one had been given had been used in service of life, or hoarded in self-protection.

Charity, in its older sense (closer to the Latin caritas, a form of love that seeks the genuine good of the other), was the moral force that Christianity believed could reorder the relationship between wealth and life. It was not optional generosity, not a private virtue to be practised at one's discretion once the important business had concluded. It was an obligation arising from the simple fact of another person's need. The surplus held by the wealthy was not entirely theirs. In a morally ordered world, it was held in trust for those who lacked what was necessary.

This did not produce a simple politics. The early Church was not a proto-socialist movement, and its relationship to property and commerce was always complex. But it produced something important: a persistent insistence that the legitimacy of ownership was not exhausted by legal title, that wealth carried responsibilities beyond the transaction that created it, and that the measure of an economic arrangement was not only what the parties had agreed to, but what kind of community it helped build or undermine.

\section*{The Moral Danger of Avarice}

Among the seven deadly sins, avarice occupied a distinctive place in medieval moral thought. It was not simply the sin of having too much, as if holiness could be measured by inventory control. It was the sin of loving money in the wrong way, of treating a means as an end and allowing the accumulation of wealth to become the organising principle of one's life and choices.

The danger of avarice was not merely personal. It was social and structural. A community shaped by avarice would lose the capacity to recognise what enough meant. The wealthy would be honoured for their wealth rather than for their virtue, their generosity, or their service. Private fortune would be treated as a public good, even when that fortune represented the redistribution of common goods into private hands. And institutions would be designed that rewarded extraction and called the result prosperity.

This is why avarice, for medieval thinkers, was the vice that most directly threatened the common good. Not because desire itself was wrong, but because desire organised around money had no natural stopping point. Unlike gluttony or lust, both at least bounded by the body's limits, avarice could intensify indefinitely. It could grow more sophisticated as more resources came its way, and dress itself in the language of prudence, of responsibility, of stewardship. It could corrupt institutions and legal structures while maintaining every appearance of respectability. It was, in Dante's moral geography, among the most serious forms of disorder, precisely because it was so capable of making itself look like order.

The persistence of this concern, running from the earliest Church Fathers through Aquinas, through the Reformation debates about usury and vocation, and into modern Catholic Social Teaching, suggests that it names something real about the structure of economic temptation. A system that rewards unlimited accumulation without asking what the accumulation is for does not remain neutral. It becomes a school for avarice, training the desires of those who participate in it toward the very disorder the tradition diagnosed.

\section*{Toward the Scholastics}

By the high Middle Ages, the moral concerns that early Christianity had raised, about avarice and usury, about the obligations of wealth and the duties of lenders, about the rights of the poor, had accumulated into a rich but unsystematic body of thought. The world, meanwhile, had grown considerably more complex. Towns were expanding, trade routes multiplying, and merchants were becoming more economically important and, consequently, harder to dismiss as a morally peripheral category. Credit was more necessary, more varied, and harder to analyse with the blunt instrument of a blanket prohibition on interest.

The moral vocabulary was rich, but it needed to be made systematic, applied with precision to the actual conditions of a commercial economy, not merely invoked as a general warning against greed whenever a merchant walked too confidently past a cathedral door. That work fell to the Scholastics.

What they inherited from Aristotle and from Christianity was a conviction that economic life was morally serious, that wealth was not self-justifying, and that the question of what commerce was for could not be answered by commerce alone. What they added was method: the patient, analytical effort to distinguish, case by case and principle by principle, between the forms of exchange and credit and profit that served justice and those that violated it. They did not reject the market. They tried to think inside it, morally, with enough precision to be useful. The result was one of the most remarkable episodes in the history of economic thought: a period in which the questions of moral philosophy and the practical demands of commercial life were held together rather than separated, each making the other more demanding and more interesting.

Aristotle had given moral economy its grammar. Christianity had given it its conscience. The Scholastics, as we shall see, would attempt something even harder: to give it a method.

\chapter*{Chapter 4: The Scholastics Discover the Market}
\addcontentsline{toc}{chapter}{Chapter 4: The Scholastics Discover the Market}

Christianity gave economic life a conscience. The Scholastics gave it a method.

By the time medieval theologians and canon lawyers turned seriously toward commerce, the world around them had changed in ways that made the older moral vocabulary feel insufficiently equipped for the job. Europe was no longer primarily a world of rural estates and localised exchange. Towns were growing, trade routes lengthening, money circulating more widely; credit was becoming more necessary, and contracts more complex. The merchant, once a figure viewed with deep suspicion by classical philosophers and Christian preachers alike, had become economically indispensable. It is one thing to warn against avarice in a village economy, quite another to do so when ships and fairs, credit instruments and partnerships and long-distance supply chains have begun to make everyone's moral life more complicated.

The Scholastics did not respond by rejecting commerce. That is among the most important facts about them. They inherited ancient and Christian suspicion toward greed and usury and unnatural accumulation, but they did not conclude that market activity itself was corrupt. They did something more difficult. They tried to distinguish. They asked when buying and selling were just, and when profit was deserved; when risk mattered morally, and when fraud or coercion corrupted exchange; and when money-making crossed from legitimate gain into exploitation. They asked not only whether a transaction had occurred, but what kind of relationship it created.

In doing so, they began to discover the market as a moral field. Not a self-justifying mechanism, nor a sphere exempt from theology and law and virtue, but not merely a swamp of sin either, though the medieval marketplace no doubt had its share of both mud and inventive dishonesty. The market became, in their analysis, a place where human beings could cooperate or deceive, serve or exploit, share risk or manufacture dependency. And it was a place where the difference between these things could and must be reasoned about carefully.

This was a genuine revolution in moral reasoning. The Scholastics did not create economics as a modern discipline. They did something different and, in some ways, more ambitious: they developed a way of thinking about economic activity before economics had separated itself from ethics. Their answers were not always consistent, and they evolved across generations. But the achievement lies in the seriousness of the attempt. They tried to think economically without surrendering moral judgment.

\section*{Commutative Justice}

At the centre of Scholastic economic thought is the idea of commutative justice.

Justice, for Thomas Aquinas and the tradition around him, is the virtue by which each person receives what is due. It concerns right order in relationships. Distributive justice has to do with how a community allocates common goods, burdens, and honours among its members. Legal justice has to do with what individuals owe the common good. Commutative justice is concerned with the exchanges between persons (buying and selling, lending and renting, hiring, compensating, restoring) and asks whether those exchanges are rightly ordered between the parties.

This is already a profound shift from treating market exchange as morally exhausted by consent. Consent matters, of course. A coerced exchange is unjust, and a deceived party has not truly agreed. But consent alone is not enough. One can formally agree to an exchange under necessity, ignorance, or manipulated conditions and still be wronged. The tenant who signs a lease because the alternative is homelessness has technically consented. So has the founder who accepts an extractive term sheet because the company will otherwise die. So has the user who clicks ``I agree'' on page seventeen of a privacy policy written in the dialect of a haunted law firm.

But commutative justice insists that the form of agreement does not exhaust the morality of agreement.

Justice in exchange requires proportionality. The thing given and the thing received should stand in some reasonable relation to each other. Neither party should use the vulnerability, the ignorance, or the necessity of the other as an opportunity for gain. This does not mean every transaction must be mathematically equal. The Scholastics understood perfectly well that economic life contains uncertainty, customary variation, and real differences in knowledge and location, in timing and circumstance. But it does mean that exchange cannot be morally purified simply because both parties signed.

Many of the most extractive structures in modern economies rely on exactly this kind of formal consent: the supplier who entered the contract, the country that accepted the financing, the open-source maintainer who chose the licence, the worker who accepted the terms. The Scholastic question cuts beneath this procedural surface to ask what kind of relationship the exchange actually created. Did it preserve reciprocity, or formalise extraction? Did it serve genuine need, or monetise vulnerability? Did it protect agency, or surrender it?

Commutative justice is not a medieval curiosity. It is one of the lost foundations of moral market reasoning.

\section*{Aquinas and the Just Price}

The just price is one of the most misunderstood ideas in the history of economic thought.

It is sometimes imagined as a fixed number determined by theologians according to an abstract moral formula. A price set from outside the market, indifferent to supply, demand, and economic reality, perhaps announced solemnly by a friar with a ledger and a severe expression. That caricature misses almost everything interesting about the tradition.

The just price was not an externally imposed correct answer. It was an attempt to ask when a price reached in exchange could be considered fair, given the full range of morally relevant factors: human need and custom, scarcity, cost, quality, risk, and the reasonable estimation of the community. For Aquinas, buying and selling exist for mutual advantage. Exchange is legitimate because one person needs what another has, and both can benefit. But because exchange can be corrupted by deception and by unequal need, it must be governed by justice. A seller who knowingly overcharges beyond what is reasonable violates the equality proper to exchange; so does a buyer who knowingly underpays. Value varies with circumstance, estimation, and scarcity, so a just price is better understood as a morally permissible range than as a single point. Within that range, bargaining may be legitimate. Outside it, exchange becomes exploitative.

What matters most in this tradition is the implication it draws: the market may help reveal value, but it does not automatically sanctify price. A price can emerge from supply and demand and still be unjust if the conditions generating it are themselves morally disordered. Famine prices may reflect real scarcity, or they may equally reflect deliberate hoarding. Housing prices reflect genuine demand, no doubt, but they also reflect speculation, manufactured shortage, financialisation, and monopoly control. Wages reflect bargaining power, and bargaining power can be unjustly distributed.

Aquinas gives us permission, and an obligation, to ask what stands behind the price. This does not mean prices can be replaced by moral intuition. Prices carry real information, coordinate activity, reflect scarcity, and reveal preference. Ignoring them can be genuinely disastrous. A moral economy that treats prices as meaningless will not become just; it will become confused, and probably hungry. But the just price tradition insists that price is never only information. It is also a moral relation. A price can tell us what someone has the power to extract; it cannot tell us, by itself, what someone is due. That distinction matters everywhere capital touches life.

\section*{Merchant Legitimacy}

One of the most consequential Scholastic developments was the moral rehabilitation of the merchant.

In earlier moral imagination, merchants were viewed with deep suspicion. They bought in order to sell at a higher price, seeming to profit without transforming the good itself, without growing or making it, without healing or teaching anything. Their gain appeared parasitic, located in the gap between producer and consumer rather than in any visible act of creation. The classical and early Christian suspicion of the merchant reflected a genuine moral intuition: if profit requires no contribution, then it can only come at someone else's expense.

The Scholastics answered this intuition carefully rather than dismissing it. Sometimes, they argued, the merchant contributes quite a lot. He transports goods from where they are abundant to where they are scarce, bearing the costs and risks (storage and travel, spoilage and theft, delay, uncertainty) that neither producer nor consumer wishes to carry. He possesses knowledge of distant markets and qualities, of timing and demand, that has real value, and he makes goods available to people who would otherwise lack them. Anyone who has ever tried to get fresh food, or medicine, or spare machine parts to the right place at the right time should be reluctant to sneer too quickly at intermediaries.

In such cases, profit is not automatically suspect. It may correspond to genuine service. This is a crucial step toward a moral theory of productive capital: gain can be legitimate when it reflects real contribution, whether of labour or skill, of risk-bearing or coordination, of knowledge or patience or responsibility. The merchant's profit is not justified merely because he successfully obtained it. It is justified when it corresponds to the service he actually provides.

The distinction is delicate, and it remains so today. A payment processor can reduce friction and enable cooperation; it can also extract rents from dependency once it controls the chokepoint. A venture investor can provide patient risk capital and genuine support; he can also impose exit pressure that subordinates a company's mission to his own return timeline. The same is true of a platform: it can solve real coordination problems, or it can enclose the field of coordination and charge tribute to those who must pass through it. The Scholastic rehabilitation of the merchant does not licence any of these outcomes. It provides a test: what real service justifies the gain? If no genuine service can be named, or if the gain is wildly disproportionate to the service, the activity has moved from productive intermediation toward extraction.

\section*{Risk, Fairness, and Reciprocity}

Economic life takes place in time, and time introduces uncertainty. Goods spoil and ships sink, debtors default and markets shift, and harvests fail. The future, in its usual inconsiderate way, refuses to be fully known in advance.

To participate in commerce is to face the possibility of loss, and the Scholastics had to confront the moral question of whether bearing that risk could justify gain. The traditional prohibition on usury had treated lending as morally simple: money is sterile, and charging for its use profits from another's need without contributing anything. But as commerce grew more complex, this simplicity became harder to sustain. A merchant who placed capital in a risky trading venture was not simply lending money to a desperate borrower. He might genuinely lose everything, and he bore uncertainty alongside the enterprise. If the venture succeeded, was some share of the profit not legitimately his?

The Scholastics began to see that the moral analysis of return on capital depended critically on whether the risk was real, whether the gain was proportionate to it, whether the terms were honest, and whether the arrangement was built on genuine partnership rather than on the exploitation of necessity. This distinction, between sterile extraction and fruitful participation, is difficult, and it remains difficult today. Investors routinely invoke risk to justify return, sometimes rightly: without capital willing to bear genuine uncertainty, many useful enterprises would never begin. But risk can also be deployed rhetorically to justify claims that far exceed any proportionate share of the actual danger borne. An investor may receive control rights well beyond the capital at risk; a fund may demand high returns while shifting downside fragility onto workers, suppliers, and the public; a platform may bear little entrepreneurial uncertainty once its network position is entrenched, and yet continue extracting from those who now depend on it.

Reciprocity, as the Scholastics understood it, requires that those who claim the upside bear proportionate responsibility for the downside. Return must correspond to contribution. One party's security should not be purchased through another party's fragility. This is the moral logic behind repayment caps and anti-usury covenants in non-extractive finance: capital may receive return, but not infinite and perpetual claim; risk may be compensated, but it should not be converted into permanent domination. The financial claim must remain ordered to the productive activity it helped make possible. This is Scholastic reasoning in contemporary dress.

\section*{Commerce as Morally Embedded Cooperation}

One of the deepest insights the Scholastics offer is that commerce is morally ambivalent because cooperation itself is morally ambivalent.

Human beings cooperate in markets because they need one another. No household or craftsperson, no farmer or merchant or city can provide everything alone. Exchange allows gifts and skills and goods and knowledge to circulate. It helps strangers serve each other across distance and difference. It can create peaceable interdependence and reward attention to another's need. These are genuine goods, and the Scholastics were not blind to them.

But cooperation can shade into domination, and interdependence into dependency. Knowledge can become a tool of manipulation, and scarcity a kind of leverage. A contract can formalise coercion in legal language while making the underlying power relation invisible. A market can look like freedom from a distance and feel like captivity from inside. This is why the Scholastics refused to let commerce define its own morality. Markets had to be embedded in virtues and law, in custom and conscience and the common good. They had to be surrounded by practices that could restrain fraud and exploitation and avarice, and interpreted by moral reason that could see what the contract form concealed.

Embeddedness does not mean centrally fixing every price or morally micromanaging every transaction. It means that exchange takes place inside a larger moral order that it did not create and cannot sustain by itself. The market is not the highest court of appeal. Economic outcomes can be judged by standards the market did not generate. Legality is not the same as justice, and profitability is not the same as legitimacy. This is precisely what modern economic thought often struggles to hold onto. Once the market is imagined as autonomous, as a self-sufficient system governed by its own internal laws, moral critique appears external and intrusive. Ethics becomes an afterthought, a constraint, or a matter of personal preference. The Scholastics assumed the opposite. Morality is not added to economic life from outside. Economic life is already moral because it consists of human acts in relation to other people. Buying and selling, lending and borrowing, investing and owning are all acts of relationship. They can honour what is due, or they can violate it. That is not a theological gloss on economics. It is a description of what economic activity actually is.

\section*{The Moral Meaning of Contract}

Contracts were central to Scholastic moral reasoning because they are the primary way economic obligations are formalised, and (it must be said) the primary way extraction becomes respectable.

A contract is more than an agreement. It is a structured relationship across time, defining expectations and duties, risks and remedies and claims. It can protect fairness by making promises explicit and enforceable. But it can also encode injustice when one party uses superior power to impose terms the other cannot reasonably refuse. The form of agreement does not exhaust the morality of agreement.

A contract may be technically valid and still be morally compromised by fraud or hidden information, by grave imbalance, or by the exploitation of necessity. Once an exploitative relationship is written down, signed, and enforceable, it acquires a kind of cleanliness. The underlying power disappears into procedure: a hard bargain becomes a document, a coercive structure becomes a clause, and extraction puts on a tie and asks to be called ``standard market terms.''

Modern economies are dense with precisely this procedural cleanliness: terms of service no user reads or understands, debt covenants that gradually transfer control, employment contracts that restrict future livelihood, venture term sheets that bury control implications in technical language, licences that allow commons-based labour to be enclosed by private platforms. In each case, the form of agreement is present; the substance of justice may not be. The Scholastic instinct is to look beneath the form to the nature of the relationship the contract creates. Was the agreement truly informed? Was it made under genuine freedom, rather than under necessity? Did it distribute burdens and risks proportionately, and preserve the purpose of the underlying activity, rather than converting it into an instrument of extraction? These questions are not merely historical. They are the questions that contract design must answer in the present. The Scholastics did not have our instruments, but they understood that principle completely.

\section*{Usury and the Habit of Distinction}

No issue forced Scholastic economic thought to become more precise than usury, because no issue more clearly demanded that similar-looking practices be distinguished by their actual moral structure.

The inherited prohibition was broad and severe: charging for the use of money was morally suspect, especially when lending to those in need, because it made another's necessity profitable without any visible labour, risk, or productive contribution. The concern was not irrational. Debt can enslave, and interest can compound beyond any possibility of repayment. A lending system can be structurally rewarded by the desperation of borrowers in ways that make exploitation not just possible but incentivised. The history of debt, ancient and medieval and modern, contains enough horror to justify the original alarm.

But commercial life was growing more complex, and the simple category of interest was proving too blunt. A merchant who placed capital into a risky trading venture was not the same as a moneylender charging a starving family. A person who lent money and genuinely bore the risk of loss was not obviously the same as someone who extracted guaranteed gain from a borrower with no alternative. So the Scholastics began distinguishing. They distinguished loans from partnerships, consumption from productive investment, compensation for genuine loss from illicit gain extracted from necessity, and money as sterile possession from capital as genuine participant in productive risk.

These distinctions were imperfect, sometimes strained, and capable of being abused. Clever people have always found ways to dress extraction in acceptable clothing, and the history of financial innovation is in part a history of finding new forms for old moral problems. But the effort itself matters enormously. The habit of distinction, the refusal to treat superficially similar practices as morally equivalent, is one of the most important intellectual tools moral economy possesses. A society that cannot distinguish productive risk-sharing from guaranteed extraction cannot protect itself from the latter while preserving the former. Slogans cannot do this work. Only careful, patient analysis can.

The Scholastics' central question, when does return on money correspond to real participation in productive risk, and when does it become extraction, is one of the central questions of this book. It returns in modern finance with new instruments and new language, but the same underlying moral structure. Does venture capital enable genuine innovation, or subordinate the enterprise to exit? Does private equity improve productivity, or strip value? Does a platform earn revenue for service, or collect tribute from dependency? These questions cannot be answered without exactly the kind of moral precision the Scholastics were developing.

\section*{Against Naivety and Cynicism}

The Scholastic approach is worth defending against two temptations that still distort economic debate, pulling in opposite directions.

The first is naivety: the belief that economic life can be governed by good intentions alone. This temptation produces moral language without institutional competence. It condemns markets without understanding what they actually do well, and dreams of justice while ignoring the mechanisms by which goods are produced, allocated, and sustained. It underestimates scarcity and incentives, coordination problems, and the impressive human capacity to act in self-interest while sincerely believing otherwise. Moral economy expressed as sentiment, without economic literacy, tends to produce policies that either accomplish nothing or make things worse while feeling virtuous. That is not an ideal outcome, though it has supplied a great deal of material for disappointing committee meetings.

The second temptation is cynicism: the belief that because economic life is genuinely complex, moral judgment is irrelevant or childish. This temptation treats legality, profitability, or market price as sufficient tests of legitimacy. It assumes that whatever sophisticated actors agree to must be acceptable, and hides power behind technical language while mistaking complexity for moral exemption. It produces a world in which any arrangement, however extractive, can be defended by pointing to its contractual form, to its market conformity, to its compliance with existing regulation. The Scholastics resisted both. They were morally serious without being economically naive. They understood that the world of exchange required practical reasoning, not purity fantasies; but they refused to let those constraints abolish conscience.

A contemporary moral economy needs exactly this balance. It must be able to read a term sheet and analyse a fund structure, to understand a software licence and follow the logic of a capital stack. But it must also be able to ask whether those instruments serve justice. Technical literacy without moral judgment produces captured analysis; moral judgment without technical literacy produces ineffective critique. The Scholastics offer a model of what it looks like to hold both together.

\section*{From Aquinas to the Franciscans}

Aquinas is central to this story, but he is not its end.

His framework of commutative justice and just price, of legitimate exchange and moral limits on profit, gives the tradition a powerful structure. But the commercial world continued to develop, and new questions pressed in ways that Aquinas's framework was not fully equipped to answer. What exactly is capital? Can money placed in enterprise become genuinely productive, not merely lent at interest but invested in a way that creates new value? What is the moral meaning of entrepreneurial foresight? When does investment participate in production rather than merely profiting from another's work?

These questions became especially urgent in the Franciscan tradition. This may seem surprising, given that the Franciscans are best known for voluntary poverty. But it was precisely because they took poverty seriously that they were forced to think carefully about money. They lived and worked among the merchants and debtors, the craftsmen and the poor of the growing medieval towns. They saw commerce from both sides: its capacity to sustain life, and its power to exploit. And they understood that simply condemning economic complexity would not help the poor, who needed access to credit, functioning markets, and institutions that could protect them from the most predatory forms of extraction.

Peter John Olivi, Bernardino of Siena, and the institutions of the Monti di Pietà would each push moral economy further, toward a more sophisticated account of productive capital, of entrepreneurial judgment, and of what it would actually take to design finance so that it could serve the poor without destroying itself. If Aquinas gives us the architecture of justice in exchange, the Franciscans begin to show how capital itself can be subjected to moral analysis and institutional design. That is where we turn next.

\chapter*{Chapter 5: Franciscans, Merchants, and the Birth of Capital}
\addcontentsline{toc}{chapter}{Chapter 5: Franciscans, Merchants, and the Birth of Capital}

The Franciscans seem, at first, like unlikely guides to the birth of capital.

Francis of Assisi did not set out to justify commerce, to refine financial contracts, to defend the dignity of merchants, or to produce a theory of capital formation that would delight later historians of economic thought. He stripped himself of property, embraced dependence, and organised his life around Christ and fraternity, around prayer and preaching and solidarity with the poor. The Franciscan imagination was marked by renunciation, not accumulation. Its founding gesture was not the signing of a partnership agreement but a man standing nearly naked before the bishop of Assisi and handing back even the clothes his father had given him.

And yet this is precisely why the Franciscans matter for the history of moral economy. Because they took poverty seriously, they had to take money seriously. They could not treat finance as a neutral technical field, nor as something spiritually dangerous and therefore beneath analysis. They lived in the growing towns of medieval Europe, close to merchants and artisans, to debtors and widows, to workers and the urban poor. They saw how credit could help households survive and how debt could destroy them; they saw that the poor suffered not only from a lack of charity, but from badly structured economic institutions. From this proximity came a paradox: the tradition most associated with poverty became one of the traditions most capable of thinking carefully about capital.

The Scholastics had taught us to ask whether exchange is just. The Franciscans sharpened the question: can capital itself be ordered toward life? And if so, how?

\section*{Poverty and Economic Intelligence}

Franciscan poverty was not simply an ethical preference for having less. It was a spiritual and social witness: a challenge to the possessive logic of wealth, a proclamation of dependence on God, and a form of solidarity with those who could not insulate themselves from need. But poverty also clarified vision in ways that comfort tends to obscure.

Those who stand near the poor see aspects of the economy the comfortable can afford to ignore. They see the cost of debt when a harvest fails, when work disappears, when illness comes and rent is still due. They see the difference between a loan that preserves a household through a difficult season and a loan that gradually transfers everything the household owns. Small amounts of credit can be life-saving; small forms of extraction can compound into ruin. For the poor, bad financial design is not theoretical. A contract can mean the loss of tools, and a missed payment can mean shame or hunger, displacement or the collapse of a family's fragile stability. The moral difference between a fair loan and a predatory one may determine whether a household survives the year.

This is why Franciscan economic thought cannot be reduced to medieval attitudes toward money. It is closer to a pastoral economics of vulnerability, one that asks how economic structures touch the lives of those least able to absorb harm. That question has not dated. A startup founder with no alternatives can be pressured into destructive investment terms just as surely as a medieval artisan could be forced into a ruinous loan. A worker without savings can be trapped in degrading labour, and a small supplier dependent on a single platform can be squeezed by changing fees and algorithms. Vulnerability is still monetisable. Franciscan moral economy begins by refusing to treat it as someone else's business opportunity.

\section*{Olivi and Productive Capital}

Peter John Olivi is one of the most important figures in this story because he helps articulate something earlier moral theology had struggled to name: capital can be more than sterile money.

The older condemnation of usury rested partly on Aristotle's claim that money, as money, is barren: a medium of exchange, a measure, a convention, but not something that reproduces naturally. To demand payment for the mere use of money appeared to make money breed from money, detached from any real labour or production. Christian tradition deepened the concern, because such lending could exploit the needy while violating both charity and commutative justice. But commercial life complicated the picture, as commercial life has a habit of doing whenever moral categories become too neat.

Consider money that is not simply hoarded or lent for consumption, but placed into productive activity: used to purchase tools, inventory, equipment, or transport; deployed to enable a merchant to undertake a risky trading voyage, a craftsperson to expand production, or a trader to bring goods from where they are abundant to where they are scarce. The person providing this capital faces genuine uncertainty and possible loss. Medieval theologians did not suddenly throw open the doors and announce that all financial engineering was now sanctified. But Olivi saw that money ordered toward productive enterprise is not morally identical to money sitting idle. When money participates in real economic activity, it takes on a different character: not merely currency, but a claim embedded in future-oriented work. Profit from capital might correspond to real economic contribution when that capital genuinely participates in productive risk and service.

Here the moral distinction that modern finance desperately needs begins to take shape. Return that corresponds to productive participation is not the same thing as return that extracts from need; sharing risk is not the same as imposing debt; enabling enterprise is not the same as controlling it; and patient capital is not the same as predatory capital. Olivi does not give us non-extractive finance in any modern sense, but he helps make it thinkable by posing the right question: when does capital serve as a partner in production, and when does it become a master over it?

\section*{Entrepreneurial Foresight and Its Limits}

Olivi's account of productive capital is inseparable from his analysis of foresight.

Commerce takes place in time, and a merchant does not merely exchange goods in a static present. He anticipates future conditions: where goods will be needed, what they may be worth, how long transport will take, what risks may arise. He acts under genuine uncertainty, and that uncertainty is part of what distinguishes his activity from simple possession. A merchant who understands that grain is abundant in one region and scarce in another may create genuine value by moving it. A trader who anticipates future demand may coordinate production before need becomes visible, and an entrepreneur who builds useful infrastructure before revenue is assured may serve a future good that does not yet exist. In each case, foresight serves provision.

But foresight can also become extraction. It can profit from anticipated scarcity without relieving it, exploit informational advantage, corner a market, create artificial scarcity, or position itself to benefit from another's coming desperation. The moral issue is therefore not foresight itself but its purpose and structure. A founder who builds open infrastructure before there is a market for it may be exercising life-serving foresight, and a fund supporting such work with patient, mission-aligned capital may participate in that service. But a platform that predicts where dependency will arise and positions itself as the unavoidable intermediary may be practising extractive foresight: intelligent, perhaps, but not just. Economic intelligence requires moral interpretation. To foresee is not automatically to serve.

\section*{Risk, Reciprocity, and the Limits of Capital's Claims}

Risk is one of the main ways capital claims moral legitimacy, and that claim contains a real truth. Without willingness to bear genuine uncertainty, many useful projects would never begin. New enterprises and trade routes, technologies and forms of shared infrastructure all require investment before results are guaranteed. Someone must carry the possibility of loss. The Franciscan and Scholastic traditions recognised this, and did not treat all gain from uncertain enterprise as immoral.

But risk is not a magic word that sanctifies every claim. It must be examined carefully, because it can also be deployed rhetorically to justify claims that far exceed any proportionate share of the actual danger borne. An investor may receive control rights well beyond the capital at risk; a fund may demand high returns while shifting downside fragility onto workers, suppliers, and the public; a lender may price risk in ways that guarantee default, then seize the collateral as the intended outcome; a platform may take genuine early risk, then use network dominance to extract rents from participants who no longer have any realistic alternative.

The moral question is not whether risk existed at some point. It is whether the continuing financial claim remains proportionate to continuing contribution and vulnerability. This is the logic behind repayment caps: capital may deserve return for risk, time, and contribution, but that claim must end or transform once it has been satisfied. Past risk does not justify perpetual extraction. Capital does not acquire infinite rights simply because it once took a chance, any more than a person who once helped you move house earns the right to live indefinitely in your kitchen. That principle sits at the heart of what Olivi was reaching toward: the idea that capital's legitimacy depends on its ongoing participation in the activity it funds, not merely on historical entitlement.

\section*{Bernardino of Siena and the Dignity of Commerce}

The Franciscan tradition did not stop with Olivi. Bernardino of Siena, the great fifteenth-century preacher, developed a powerful account of the merchant's dignity, and it is worth pausing on the fact that a Franciscan was making this argument.

Bernardino understood that commerce could genuinely serve the common good when practised with justice, skill, and responsibility. Merchants moved goods and provided needed services, assumed risks and employed workers, and coordinated economic life across distances that individuals alone could not span. Their work could be honourable when it served real needs and respected moral limits. But Bernardino did not simply bless profit. He distinguished good merchants from bad ones, and the distinction turned not on wealth but on virtue: on diligence and knowledge, on reliability and prudence, on truthfulness, and on concern for justice. Commerce required more than cleverness. It required character.

This has direct application to how we think about entrepreneurship today. Founders are often celebrated simply for disruption and scale, for fundraising ambition and the ability to attract capital. The cultural mythology is familiar: raise money, grow fast, break things, appear on a conference stage, and perhaps wear shoes that imply both wealth and indifference to it. But moral economy asks a different set of questions. What kind of innovation, what kind of disruption? What kind of dependence is being created, and what kind of power is being concentrated? A founder building useful tools, or dignified work, or resilient institutions, or open infrastructure participates in merchant dignity. A founder building addiction, surveillance, manufactured scarcity, or regulatory evasion does not become morally legitimate simply by being entrepreneurial. Commerce can be a genuine vocation, as Bernardino saw, but only when it remains accountable to the common good.

\section*{The Monti di Pietà and the Problem of Institutional Design}

The Monti di Pietà were among the most important institutional experiments in late medieval and early modern moral economy, and they matter not as historical curiosities but as a case study in a problem that has never really been solved: how do you design a financial institution that serves the vulnerable without eventually becoming indistinguishable from what it was created to replace?

The Monti emerged in response to a practical crisis. The poor needed access to credit for any number of reasons (illness or unemployment, failed harvests, tools, rent, temporary distress), but available lending was often exploitative, with rates high enough to transform a short-term emergency into a long-term trap. The Monti attempted a third way: charitable lending institutions, often publicly supported and religiously inspired, providing small loans against pledges at low cost, with the explicit aim of preventing exploitation and preserving household stability.

What makes the Monti genuinely instructive is that moral critique had to become operational. The question was not merely whether usury was wrong. It was how to design a lending institution that could actually meet the needs predatory lenders were meeting, without becoming predatory itself. That meant statutes and governance structures, collateral rules, accounting systems and reserves, lending rates and audits, and explicit limits on what the institution could do. Good intentions were not enough. The institution had to be built. And building it revealed a permanent tension: the fund must survive in order to serve. If the charitable capital is depleted, it cannot help anyone. If administration is unfunded, the institution collapses. If lending rates are too low, the institution becomes unsustainable; if too high, morally compromised. These are not abstract dilemmas. They are the recurring design problems of every non-extractive financial institution. Anyone who thinks ethical finance can escape accounting, governance, reserves, and default risk has not yet had the pleasure of running an ethical financial institution.

\section*{The Moral Problem of Sustainability}

The Monti also force us to confront a harder truth: the language of sustainability can become the moral vehicle for mission drift.

The original logic is sound and necessary. A charitable fund must not be depleted if it is to continue serving. Charging enough to cover administrative costs, absorb losses, and preserve capital may be not merely permitted but required. The poor are not served by an institution that collapses after a generation. A lender that disappears in a mist of good intentions does not provide much long-term justice.

But the same argument, applied without discipline, can travel far from its origins. The institution charges enough to survive, and then enough to build reserves, and then enough to expand, and then enough to compete with other lenders, and eventually enough to attract outside capital and to satisfy stakeholders who have come to expect a surplus. At some point the original mission has become a story told about an institution whose real logic has changed beneath it. This is not a hypothetical sequence; it describes the recurring pattern of institutional drift across many sectors. A nonprofit becomes dependent on funders whose priorities reshape its work. A social enterprise accepts investment requiring growth incompatible with its purpose. An ethical lender raises rates in order to scale. A mission-driven company exits to a buyer whose governance cannot protect what mattered.

The Monti remind us that moral finance must guard against the internal logic of institutional survival at least as vigilantly as it guards against external predation. Survival is necessary, but it is not the highest good. A fund exists to serve its mission. If preserving the fund requires betraying the mission, then the fund has become an end in itself, and ends in themselves, as Aristotle observed, are the beginning of disorder. This is why non-extractive financial design needs not only maximum return caps and repayment formulas, but explicit mission tests: governance duties alongside governance rights, clarity about why reserves exist, and a clear account of what must never be done in the name of sustainability.

\section*{Moral Accounting}

The Monti also reveal something that modern finance has largely forgotten: accounting is a moral practice.

A ledger is not merely a technical tool. It embodies a way of seeing, a judgment about what counts and what must be preserved, what may be spent and what is owed, what is surplus and what requires explanation. An institution that tracks only financial return will see borrowers differently than one that also tracks household preservation, avoided exploitation, the causes of default, and mission integrity. The categories in which we account shape the decisions we are capable of making.

The moral ledger cannot replace the financial ledger. But the financial ledger is incomplete without it. For contemporary non-extractive finance, the practical questions are not difficult to state, even if they are demanding to answer. Are repayments coming from genuine free cash flow, or from distress and underinvestment? Is the enterprise becoming more resilient, or more fragile? Are workers, users, and communities strengthened or depleted by the relationship with capital? Does the capital structure protect the mission beyond the current leadership? Accounting should make these questions harder to evade. That is what it means for a ledger to carry moral weight.

\section*{Productive Capital and the Common Good}

The Franciscan contribution helps us state more clearly what productive capital actually means, and it is more demanding than the term usually implies.

Productive capital is not simply capital that generates profit. Extractive capital also generates profit, often very reliably. Indeed, one of the more annoying things about extraction is that it can be financially excellent right up until it has consumed the conditions of its own existence. Productive capital, in the morally serious sense, is capital that participates in the creation, preservation, or circulation of real goods in a way that serves life and remains proportionate to its contribution. It accepts limits. It does not use vulnerability as the foundation for domination.

This definition is more demanding than ordinary financial language allows. Modern usage tends to call capital ``productive'' if it increases output or return. Moral economy asks a harder question: productive of what, for whom, and at what cost? A factory can increase output while degrading workers and exhausting land; a platform can produce convenience while destroying local capacity and capturing dependency; a fund can produce returns while weakening the companies it owns. Production alone is not sufficient. Capital serves life when the productive activity itself is rightly ordered, when it strengthens rather than depletes the community of production, when it preserves future agency rather than consuming it, and when it enables households, workers, users, and institutions to flourish over time. The common good is not an external constraint on productive capital. It is the standard by which productive capital must ultimately be judged.

\section*{From Charity to Institution}

One of the most important movements in this chapter is the transition from charity to institution, from the personal moral response to the structural one.

Charity is indispensable. Christian moral economy cannot exist without love of neighbour, care for the poor, almsgiving, and mercy. But charity alone is not enough when the structures of economic life produce recurring harm. If the poor are trapped by usurious lenders, almsgiving may help after the damage is done; a better lending institution can prevent some of the damage in the first place. If founders are forced into extractive capital, philanthropy may support a few alternatives; better financing structures can change what options are available at all.

The Monti embody this shift. They do not merely tell rich people to be generous. They build a financial mechanism that changes the credit options available to the poor. Moral concern becomes institutional form. A society ordered toward life cannot rely only on heroic personal virtue. It must build structures that make justice more ordinary, institutions in which the easier path is not always the extractive path, and in which stewardship is durable beyond the moods and personalities of individual stewards. That is the promise of steward ownership, of mission locks and evergreen funds, of repayment caps: attempts to institutionalise moral intention so that it does not evaporate under financial pressure. Charity responds to need. Institution changes the conditions under which need arises. Both are necessary. But without institution, charity is always repairing what extraction keeps breaking.

\section*{The Franciscan Ambiguity}

The Franciscan story is powerful, but it is not pure, and we should not pretend otherwise.

No tradition that enters seriously into the world of money remains untouched by ambiguity. The very act of designing morally disciplined finance creates the temptations it was designed to resist. Exceptions become loopholes. Sustainability becomes expansion. Administrative necessity becomes margin. Low-interest lending begins to resemble the practices it was created to oppose. Moral language can legitimise institutions that have drifted, with little fanfare, from their founding purpose, providing cover that a purely commercial institution would not enjoy.

This ambiguity should not produce cynicism. It should produce vigilance. The lesson is not that ethical finance is impossible but that it must expect drift and design against it, assuming that every justification can be abused and building counter-mechanisms accordingly: audits and governance checks, transparency requirements, caps and mission locks, and the practice of returning periodically to the founding questions. Are we still serving whom we said we would serve? Are we still an alternative to extraction, or have we become a more comfortable version of it? The Franciscans help us ask these questions because their own tradition lived in the tension between poverty and institutional design, between renunciation and financial architecture, between moral purity and practical compromise. That tension is not a flaw in the story. It is the story.

\section*{Capital as Moral Category}

The title of this chapter speaks of the birth of capital, but the phrase requires care.

Capital in the broad sense (tools and land, money and stored grain, ships and buildings) existed long before the Franciscans. What begins to emerge more clearly in Franciscan and late Scholastic thought is not capital as a material fact but capital as a moral category. Capital is not merely wealth. It is wealth placed into time: wealth ordered toward future activity, toward the possibility of production, preservation, and service. And because it is oriented toward the future, it can be asked what future it is oriented toward. It can preserve and produce, circulate and serve; it can also dominate and extract, enclose and enslave.

Without a moral category of capital, finance becomes opaque. Money moves, returns are calculated, contracts are enforced, and institutions grow, but the deeper question disappears: is this capital serving life, or feeding on it? The Franciscans do not let that question disappear. That is their lasting contribution.

The next stage of the story moves into the School of Salamanca, where the discussion broadens further. Salamanca's thinkers will develop more sophisticated accounts of subjective value, market price, supply and demand, money, and international trade. But Salamanca will also introduce a deeper ambiguity, because the intellectual world that reflected on market freedom and price was also the world of empire. Productive capital and extractive capital do not always appear as opposites. Sometimes they travel together, carried by the same trade route, spoken of in the same language of provision and opportunity. Moral economy must therefore learn to discern whole systems, not only individual transactions. That is what Salamanca will demand of us.

\chapter*{Chapter 6: Salamanca --- Market Coordination and Moral Ambiguity}
\addcontentsline{toc}{chapter}{Chapter 6: Salamanca --- Market Coordination and Moral Ambiguity}

If the Franciscans sharpened moral distinctions around productive capital, the School of Salamanca broadened the entire horizon of economic reasoning.

Earlier Scholastics had focused primarily on the ethics of individual transactions. Was this price just, this loan usurious? Did this merchant deserve his profit? Did this contract preserve reciprocity, or dress extraction in respectable clothing? These were essential questions, and they remain so. But Salamanca widened the lens. The question was no longer only whether a single exchange was fair, but how a network of exchanges organises resources across space and time, and what moral principles should govern that organisation.

The School of Salamanca flourished in sixteenth-century Spain, gathering thinkers at the University of Salamanca and in the Dominican and Jesuit intellectual traditions associated with it. Francisco de Vitoria, Domingo de Soto, Martín de Azpilcueta, Luis de Molina, Juan de Mariana: these were not economists in the modern sense. They did not have academic departments, regression models, policy briefs, or the consolation of being misunderstood by journalists. They were theologians, jurists, and moral philosophers living through one of the most dramatic expansions of commerce in human history. Trade routes stretched to the Americas and Asia, silver flooded in from New World mines, and prices shifted in ways that puzzled and harmed ordinary people. Money and empire, conquest and trade and moral theology, collided with unusual force.

Salamanca did not invent modern economics. It did not assume perfectly competitive markets, atomised rational agents, or the mathematical formalism that would characterise later classical and neoclassical thought. But it did develop concepts that anticipate subjective value theory and supply-and-demand reasoning, monetary theory, and the coordination properties of price, all while remaining embedded in a moral framework that refused to treat market outcomes as self-justifying. That combination is what makes Salamanca important, and what makes it troubling when its ambiguities are examined honestly.

\section*{Subjective Value and the Moral Meaning of Price}

One of the most significant intellectual developments in the Salamanca tradition was the recognition of subjective value.

Earlier Scholastic discussions of the just price had tended to ground value in the objective properties of goods: their labour content, their usefulness, their cost of production. Aquinas had emphasised proportionality and the equality proper to exchange. The Salamantine thinkers pushed further, observing that the value of a good is not simply inherent in the object itself. It arises from the intersection of human need and estimation, of scarcity and circumstance. The same measure of grain is worth more in a famine than in a harvest surplus; a medical supply worth little in ordinary times may be priceless in a plague. Value is relational and contextual, not fixed by nature, by decree, or by the secret opinion of a committee of theologians in excellent hats.

This was a genuine advance in understanding. It helped explain why prices vary across time and place without requiring either divine intervention or moral failure. The Salamantines recognised that dispersed human knowledge (about local scarcity and future needs, about risks and transport costs and conditions invisible to distant observers) is embedded in the prices that emerge from exchange. They did not fully articulate what Hayek would later call the price system as a mechanism for communicating dispersed knowledge, but they were seeing the same phenomenon from a different angle.

The moral implications, however, were not straightforward. If value is subjective, then a price that reflects genuine scarcity and real need may be just even if it seems high. A merchant who correctly anticipates where goods will be needed, who bears the cost of transport and risk, and who sells at a price reflecting these conditions may deserve his profit, even if buyers find the price painful. But the same logic of subjective value can be turned in a darker direction. If value is whatever someone will pay under their circumstances, then the desperate borrower's willingness to accept any interest rate reflects the ``true value'' of the loan, and the tenant with nowhere else to go reveals by her acceptance the ``just price'' of the rent. Subjective value, taken alone, can become a sophisticated rationalisation for whatever the market will bear, including the extraction of maximum return from maximum vulnerability.

The Salamantine tradition recognised this danger and tried to hold the line. Azpilcueta and others insisted that justice remained contextual in a way that included, not excluded, moral evaluation of the conditions generating the price. A price reflecting genuine scarcity and real service may be just. A price reflecting manufactured scarcity, monopoly control, or the exploitation of emergency need is not sanctified simply because someone agreed to it under duress. The market reveals information. It does not, by itself, produce justice. That distinction, clearly stated but not always successfully maintained, is one of the most important and contested legacies of the Salamantine school.

\section*{Supply, Demand, and the Coordination of Markets}

The Salamantine thinkers were also among the first to articulate something approaching systematic supply and demand reasoning, and they did so in a way that kept moral questions visibly in play.

They observed that scarcity drives price upward while abundance pushes it down. They saw that desire and urgency shape what buyers will pay, and that availability of substitutes, risk of spoilage, transport costs, and timing all influence the terms on which exchange occurs. These were not merely empirical observations. They were the building blocks of a theory of how markets coordinate the movement of goods from where they are abundant to where they are needed: a process that could serve genuine human welfare when it functioned well.

But they also observed that the same dynamics that make markets useful can make them dangerous. Scarcity is sometimes natural: a bad harvest, a delayed shipment, an unexpected frost. The world, unfortunately, does not always organise itself around human convenience. But scarcity can also be manufactured. A merchant who withholds grain during a famine to force a higher price is not coordinating a market; he is exploiting one. The resulting high price reflects not genuine conditions of supply and demand in any morally meaningful sense, but the calculated exploitation of necessity. Here the Salamantine inheritance from earlier Scholastic thought is clear: the moral status of a market price depends not only on the number itself, but on what produced it and whose vulnerability it is trading on.

This gives Salamantine market reasoning a texture that later classical economics would largely lose. Supply and demand analysis in its modern form treats market prices as information signals, coordinating actors toward efficient outcomes. That insight is real and valuable. What it does not easily accommodate is the question of whether the conditions generating those signals are themselves just. Salamanca kept that question alive, insisting that the market's coordination function is valuable but not morally self-sufficient, and that moral analysis of price must include analysis of the conditions that produced it.

\section*{Money, Inflation, and Azpilcueta's Insight}

One of the most striking contributions of the Salamanca school was its engagement with monetary phenomena. The figure most central to this development is Martín de Azpilcueta, sometimes called the Doctor Navarrus, whose work on the value of money anticipates what would much later be called the quantity theory.

Azpilcueta observed that the massive influx of silver from the Americas was driving up prices across Spain and Europe. This was not immediately obvious to everyone. Prices were rising, certainly, but the connection to the money supply required careful observation and reasoning. His insight was that money itself has no fixed intrinsic value. Its value, like that of other goods, depends on its abundance or scarcity relative to the goods it circulates. When the quantity of money increases without a corresponding increase in goods, prices rise; when money is scarce, prices fall. The relationship is not mechanical, and is subject to many complicating factors, but it is real.

What made this more than a technical observation was its moral significance. Price inflation caused by monetary expansion was not simply a neutral economic phenomenon. It redistributed wealth in ways that were often deeply unjust. Landowners paid in rents set by old contracts, workers paid at customary wages, creditors owed fixed sums: all found the real value of what they received declining. Those who held goods, or who could adjust their prices quickly, benefited. Debtors benefited and creditors lost. The poor, who held little property and had the least ability to adjust, often bore the greatest real burden.

The Salamantine analysis of money was therefore not separable from questions of justice. Debasement of coinage and manipulation of currency, inflationary monetary policy, the flooding of markets with undervalued money: these were not merely technical economic events. They were acts with moral consequences, capable of enriching some at the expense of others through mechanisms invisible to those who suffered them. Money is a tool, and like all tools it can be used well or badly. It can facilitate exchange, or become a hidden instrument of redistribution and extraction. The Salamantine recognition that monetary mismanagement converts money from a servant of exchange into an instrument of power anticipates debates that remain unresolved today.

\section*{Empire, Conquest, and the Moral Limits of Market Reasoning}

The deepest and most troubling dimension of the Salamanca tradition is its entanglement with empire.

The same intellectual world that developed sophisticated market theory was also the world of Spanish conquest in the Americas. The thinkers who analysed just price and monetary inflation were contemporaries of the encomienda system, of the forced labour of Indigenous peoples, of the destruction of Indigenous economies, and of the extraction of silver that financed European expansion. Francisco de Vitoria, the most important founding figure of the school, engaged directly with the question of whether the Spanish conquest was morally legitimate, and his answer was, at minimum, deeply ambivalent. He challenged some of the most extreme justifications for conquest, argued that Indigenous peoples had genuine rights of property and self-governance, and questioned whether military force could be justified by the refusal to accept Christianity. These were serious interventions. They mattered. They placed limits on the fantasies of total imperial entitlement.

But Vitoria and his colleagues also produced reasoning that, in practice, could be used to accommodate conquest. The argument that peaceful trade conferred a natural right of entry to foreign territories, and that interference with such trade could justify military response, was a principle that could serve commercial empire as readily as it could constrain it. The Salamantine tradition developed a sophisticated account of natural law, of just war, and of the rights of peoples, yet developed this account within a world in which European power was expanding violently, and the institutional structures of that expansion were not fundamentally challenged.

The Salamantine thinkers were not simply apologists for empire; several of them were among the most serious critics of specific forms of colonial violence in their time. But the moral vocabulary they developed, of market freedom and price coordination, of natural rights of exchange and the legitimacy of commercial activity across borders, could travel alongside imperial power in ways its authors did not fully control. Market reasoning, even when developed with genuine moral seriousness, can be extracted from its moral context and deployed in service of power. A language of freedom can describe the freedom of the powerful to set terms for everyone else. This is not a reason to dismiss Salamantine thought. It is a reason to receive it critically: to take seriously what it saw clearly, while remaining honest about what it failed to prevent. The tradition that understood market coordination was the same tradition that could not fully resist the moralisation of conquest. That should make us humble about our own blind spots.

\section*{Coordination Without Moral Neutrality}

The deepest contribution of Salamanca to this book's argument is perhaps this: markets can coordinate, but coordination is not the same thing as justice.

The Salamantine thinkers saw more clearly than most of their predecessors that prices, when they function well, carry real information. They reflect scarcity and desire, urgency, and comparative advantage. They allow strangers to cooperate across distance without needing to know each other's circumstances in detail, and they provide incentives for production, conservation, and movement of goods toward need. These are genuine goods, and any moral economy that ignores them will make costly mistakes. It is not morally serious to design institutions as though information, incentives, and scarcity do not exist. Reality has a way of returning one's memo with corrections.

But the Salamantine tradition at its best also saw that none of this made market outcomes morally self-sufficient. A price can coordinate efficiently and still be unjust; a market can clear and still distribute harm unequally; trade can expand and still reproduce domination. The coordination properties of price are real and valuable. They are not a substitute for moral evaluation of the conditions that generate those prices, of the distribution of power among those who set them, of the background institutions that determine who enters markets on what terms, and of the social costs that markets externalise onto those with the least power to resist.

This is why the Salamantine tradition is not simply a precursor to modern economics, even though it anticipates many of modern economics' core concepts. Modern economics largely accepted the coordination insights while setting aside (or treating as external to the discipline) the moral evaluation those insights were originally embedded in. The result was a more powerful analytical tool and a less morally literate one. Salamanca's legacy suggests that the separation was a choice, not a necessity. One can understand how markets coordinate and still ask whether they are coordinating justly. One can take price seriously as information and still ask what conditions produced that information, and at whose expense.

\section*{What Salamanca Leaves Behind}

The School of Salamanca leaves us with a richer and more complicated picture of moral economy than we had before it, and a more honest one.

It gives us a more sophisticated account of how value arises from human need and circumstance rather than from fixed properties of objects. It shows that prices can serve genuine coordination functions, and that this matters morally, because markets that coordinate well can serve life in ways central planning cannot easily replicate. It offers early insights into the monetary dimension of economic life and the ways money can be used to redistribute wealth invisibly and unjustly. And it forces us to confront the irreducibly difficult question of how market reasoning relates to imperial power: how the same intellectual tradition can produce both genuine moral insight and moral accommodation of domination.

No clean summary does this tradition justice. Its thinkers were serious and subtle, and sometimes wrong in ways that had devastating consequences for people on the other side of the world. That combination is uncomfortable but instructive: intellectual sophistication does not automatically produce moral clarity, and a tradition can see markets with remarkable precision while failing to see what those markets were built upon.

The tradition that follows, in Protestant thought, in the moral psychology of Adam Smith, and in the eventual separation of economics from moral philosophy, will inherit all of this complexity. It will develop the coordination insights while often losing the moral embeddedness, and build a more powerful analytical economics while building a less morally attentive one. Understanding how that happened, and what was lost in the process, is what the next chapters are about.

\chapter*{Chapter 7: Vocation, Stewardship, and the Protestant Ethic}
\addcontentsline{toc}{chapter}{Chapter 7: Vocation, Stewardship, and the Protestant Ethic}

The Reformation changed the question.

The traditions we have been tracing (Greek and Christian, Scholastic and Franciscan and Salamantine) had mostly been asking whether commerce could be just, whether wealth could be ordered toward life, whether exchange might be made to serve need rather than exploit it. These were essential questions. But now came a different and deeper one: whether ordinary work itself could be a place of divine calling.

Protestantism brought the language of vocation into the world of daily labour, household responsibility, craft, and commerce. It challenged the older hierarchy that placed religious life above ordinary life: monastery above workshop, contemplation above making, celibacy above household, poverty above stewardship. For Luther and those who followed him, the farmer and the mother, the magistrate and the printer, the teacher and the merchant could all serve God through the duties of their station. Work was not merely survival or accumulation. It could be a calling. The ordinary world was not spiritually irrelevant. It was the arena in which faith became obedience, service, discipline, and care.

This was a profound moral development, and a dangerous one.

When vocation remains ordered toward God and neighbour, toward household and common good, it deepens moral economy. It dignifies ordinary work and disciplines idleness, resists clerical elitism, and places economic life under divine judgment. But when vocation detaches from grace and Sabbath, from humility and limits, it can mutate. Work becomes identity, discipline becomes anxiety, reinvestment becomes accumulation, and calling becomes career. The desire to serve God in the world becomes the compulsion to prove oneself through worldly success. This is where Max Weber enters the story, and where the plot thickens considerably.

\section*{Luther and the Dignity of Ordinary Calling}

Martin Luther's doctrine of vocation begins as a pastoral and theological revolution.

In late medieval Christianity, religious life often appeared as the higher path. Monks and nuns, priests, and contemplatives were understood to be pursuing a more perfect state than the ordinary baptised person could achieve within the world. Luther challenged this hierarchy at its root. For him, the Christian does not need to flee the world in order to serve God. A father caring for a child, a farmer tilling his land, a cobbler making shoes, a magistrate judging disputes: all can participate in God's providential care for creation through the ordinary duties of their station. Vocation is not confined to explicitly religious work. It is the form love takes in one's concrete responsibilities.

This matters for moral economy because it dignifies ordinary labour. Work becomes service to neighbour, not merely a means of survival. The cobbler serves God not by putting little crosses on shoes but by making good shoes; the farmer serves by feeding others; the merchant, if honest, serves by making goods available. Economic life becomes a field of obedience. This deepens the anti-extractive logic of moral economy in a specific way. If work is service, then exploitation is not merely inefficient or legally problematic. It is a violation of vocation. To treat workers as disposable inputs is to deny the dignity of their calling. Luther's doctrine therefore resists both monastic contempt for ordinary labour and the capitalist reduction of labour to commodity.

But Luther's view has real limits. Because vocation is tied to one's station, it can become conservative in problematic ways, sanctifying existing hierarchies, encouraging obedience within unjust structures, and dignifying work without transforming the conditions under which work is performed. Vocation can liberate ordinary work from spiritual inferiority while simultaneously spiritualising endurance under domination. A moral economy must receive Luther's dignity of ordinary calling while refusing to turn vocation into a justification for unjust conditions. The worker's calling does not excuse the employer's extraction.

\section*{Luther on Usury, Calvin on Interest}

Luther was fierce against usury; fiercer, in some ways, than the Scholastic tradition he inherited. For him, usurers were not clever financial technicians navigating difficult moral terrain with elegant instruments and a regrettable public image. They were predators. They profited from distress, evaded the demands of love, and used legal form to conceal moral violence. Luther saw debt not merely as a contract but as a field where the vulnerable could be trapped by the powerful.

His critique is more direct and pastoral than the Scholastic tradition's. Where Aquinas worked through categories of commutative justice and proportionality, Luther asked a simpler and sharper question: what does this transaction do to my neighbour? Does it help him, or burden him; does it preserve his household, or exploit his need? For Luther, the Christian cannot hide behind technicalities. Love strips away the disguises of legalism.

Calvin's approach was more nuanced and, in some ways, more consequential. He did not simply bless lending at interest. He distinguished more carefully between different kinds of lending and different kinds of borrowers. A loan to a wealthy merchant engaged in productive enterprise is not morally identical to a loan to a poor person in distress. Where capital genuinely participates in commerce, generates real gain, and does not exploit vulnerability, moderate interest may be legitimate. But Calvin insisted on limits. The poor should not be burdened by interest. The lender must not neglect equity and charity. Rates must not be excessive, and the arrangement as a whole must serve the common good. The legality of interest does not eliminate moral judgment; if anything, it makes such judgment more necessary.

Calvin's more flexible approach is a genuine advance in moral precision. But it carries a risk he could not fully control. Once all interest is no longer treated as automatically usurious, the moral task becomes more demanding, and more easily evaded. What begins as a careful exception becomes ordinary expectation; what is justified under conditions of productive enterprise gets applied to consumption and housing, to healthcare and distress; what is tolerated under moral limits expands until those limits are forgotten. Modern capitalism largely failed to maintain the discipline that Calvin's exception required. The exception swallowed the rule.

\section*{Work as Calling}

The Protestant doctrine of vocation changed the moral imagination of work in ways whose effects are still with us.

Work became more than survival, more than social status, more than preparation for almsgiving. It became the ordinary place where faithfulness was enacted. One served God not by escaping daily responsibility but by fulfilling it well and in love. This dignified ordinary labour and challenged idleness; it encouraged literacy and skill; it emphasised that household and trade, profession and civic office were all arenas of moral responsibility.

For moral economy this matters deeply. If work is calling, then economic life must be judged by whether it allows people to exercise that calling well. A workplace that destroys health and fragments attention, that prevents rest and reduces human beings to replaceable units, is not merely unfortunate. It is vocationally disordered. A company that makes useful things, treats workers with dignity, honours craft, and sustains the community can be a place of vocation. A company built on addiction, surveillance, deception, or the systematic degradation of work cannot be redeemed simply because some people inside it work hard. Industry, like piety, can be misdirected.

The dignity of work depends not only on the worker's intention, but on the purpose and structure of the work itself. It is not enough to tell people to find meaning in their jobs. Some jobs are badly ordered. Some industries are built around harmful ends. Vocation requires institutional conditions: fair compensation and truthful management, protection from exploitation, work that contributes to real goods, and space for family and community, for worship and rest. Without these, vocation is a pious word placed over economic coercion.

\section*{Stewardship}

Stewardship is among the most important economic categories in Christian thought, and among the most vulnerable to corruption.

A steward is not an absolute owner. A steward receives responsibility for something that ultimately belongs to another, and must manage, preserve, cultivate, and use what has been entrusted for the purposes of the true owner. In Christian terms, creation belongs to God. Human beings are not sovereign possessors but accountable trustees. This changes the meaning of property. Ownership becomes responsibility; control, service; wealth, trust. Capital becomes something to be governed in light of its purpose rather than exploited for private gain.

Stewardship therefore belongs at the heart of moral economy. It resists both possessive individualism and careless collectivism. It does not deny that particular persons or institutions may hold property. But it insists that holding property makes one morally accountable: accountable for how it is used, for what it makes possible, for whom it affects, and for what comes after. The investor must answer for influence, the entrepreneur for the mission, the manager for workers and customers, the technologist for the dependencies their infrastructure creates. A person does not cease to be responsible simply because the harm is mediated through a capital structure, a supply chain, a software platform, or a series of extremely professional-looking documents.

Stewardship also introduces time in a way that financial short-termism resists. A steward does not consume everything now. A steward preserves for future use, resists liquidation when preservation is required, and distinguishes between surplus that may be distributed and capital that must be held for the mission. Capital that serves life is stewarded capital: held under purpose, governed by limits, accountable to those affected by it. Capital that feeds on life is anti-stewardship. It treats possession as licence and liquidity as superiority, the future as collateral and exit as success, even when exit destroys the purpose for which the enterprise existed.

But stewardship has its own danger: it can be reduced to personal virtue while structures remain extractive. A person may speak of stewardship while holding shares in companies governed by extraction; a founder may speak of stewardship while accepting capital that forces exit. Stewardship must therefore become institutional, built into ownership structures, governance documents, fiduciary duties, and capital terms, or it will be overpowered by the architecture around it. If ownership rewards extraction and capital demands exit, personal stewardship will not survive for long. Good intentions are not load-bearing beams. This is why later chapters turn to steward ownership and mission locks, to evergreen finance and repayment caps: not as secular replacements for stewardship, but as attempts to give stewardship legal and financial form.

\section*{Puritan Discipline and the Shell of the Ethic}

The Puritan intensification of vocation introduced another important element: disciplined reinvestment.

If work is calling, idleness is suspect. If consumption can become vanity, restraint is virtuous. If wealth is entrusted by God, it should not be wasted in luxury. The result is a culture of discipline. One works carefully and spends modestly, reinvests surplus and avoids frivolity, examines oneself, and orders one's life rationally around productive activity. There are worse social habits than showing up on time, keeping accounts, repairing the roof, and not spending the winter stores on decorative nonsense.

But discipline is morally unstable when separated from joy, mercy, rest, and grace. Work can become a way to prove election or worth or superiority. Reinvestment can become accumulation, thrift can become hardness toward the poor, and productivity can become a spiritualised form of anxiety. Success can become evidence of virtue; failure, evidence of moral defect. This is the pattern Weber identified. In certain strands of Protestant asceticism, worldly activity was charged with intense spiritual significance, and consumption was restrained not because wealth was rejected but because wealth was to be rationally ordered. The unintended result could be the accumulation of capital without its traditional release through luxury. If this religious discipline then weakens while the disciplined behaviour remains, the practice survives its justification. The shell remains after the soul has departed. Work continues, reinvestment continues, accumulation continues; but Sabbath fades, mercy fades, and limits begin to look irrational. That is the Protestant ethic corrupted.

\section*{What Weber Got Right and What He Missed}

Weber's argument about the Protestant ethic and the spirit of capitalism remains powerful because it identifies a moral and spiritual transformation beneath economic change. Capitalism requires not only markets and capital but a particular kind of person: disciplined and calculating, future-oriented, willing to work systematically, to reinvest surplus, and to organise life rationally around productive activity. Weber saw that certain forms of Protestant asceticism helped form this kind of person. He also saw that modern capitalism eventually no longer needed the theology that generated its spirit. The religious roots could wither while the economic discipline remained, and people would continue to work, save, and accumulate not because they felt called by God but because they were trapped in a system that demanded it. The ethic became a cage.

This is one of Weber's most enduring insights. Economic systems outlive the moral worlds that formed them, and practices survive after their purposes have been forgotten. Modern capitalism is populated not only by hedonists but by exhausted moral strivers: people who work compulsively, optimise themselves, reinvest and measure and delay rest, and sacrifice relationships and contemplation, all while experiencing this as duty rather than pleasure. The system captures conscience, not only resources.

But Weber underemphasised something crucial. The Protestant ethic is not simply the origin of the capitalist spirit. It is a corrupted form of something better. Vocation corrupted becomes careerism, stewardship corrupted becomes accumulation, discipline stripped of grace becomes anxiety, and thrift stripped of mercy becomes cruelty toward the poor. These concepts became dangerous not because they were wrong but because they were severed from the theological and moral order that gave them their proper shape. The sharper thesis is this: the Protestant ethic was not stewardship born. It was stewardship intensified, and then, in modern capitalism, stewardship corrupted.

There is therefore much to recover. The dignity of ordinary work, responsibility in one's station, disciplined use of resources, resistance to waste, the idea that all of life belongs under calling: these are real goods. But they must be re-embedded within their restraints. Sabbath interrupts production. Grace declares that worth is not earned through output. Mercy prevents discipline from becoming cruelty, and the common good prevents private stewardship from becoming private accumulation. Without these, Protestant discipline becomes one of the spiritual engines of extraction rather than one of its opponents.

\section*{Sabbath and the Doctrine of Enough}

Sabbath may be the most neglected economic doctrine in the modern world.

Sabbath is not merely a religious day off, a quaint remnant of agricultural rhythms that the smartphone has mercifully rendered obsolete. It is a limit placed on production and consumption, on mastery and anxiety. It interrupts the claim that more work is always better, that time must always be monetised, that land must always produce and workers must always be available. It declares that creation is gift before it is resource, that human beings are received before they are productive, and that rest is not justified by future efficiency. The poor and the servants, the animals and the strangers and the land are all included in the Sabbath's rhythm of release. The economy does not have total claim over life.

For a book about capital that serves life, Sabbath is indispensable. Capital that cannot rest will eventually devour. A company that cannot accept enough will extract; a fund that cannot tolerate patient time will force exit; a platform that cannot stop optimising engagement will colonise attention; a society that cannot interrupt productivity will sacrifice bodies and families and communities and contemplation to output. Sabbath is the theological form of enough, and enough is the category that extractive capitalism structurally cannot recognise.

A recovered Protestant contribution to moral economy must place Sabbath back beside vocation. Work is a calling, but it is not a god. Stewardship is a responsibility, but not a mandate for domination. Reinvestment can be prudent or compulsive; capital can be useful or sovereign. Without Sabbath to interrupt the logic of accumulation, the Protestant ethic is an engine with no brake.

\section*{The Corruption of Stewardship}

The corruption of stewardship happens slowly, and it always begins with something defensible.

Capital is accumulated to serve a purpose. A family saves to preserve the household, a business builds reserves against uncertainty, a fund raises capital to support useful enterprise. These are not corrupt beginnings. Indeed, they may be necessary and wise. A family without savings is vulnerable, a business without reserves fragile, and a fund without capital is mostly a concept with a bank account and aspirations.

But over time, the means can overtake the end. The family fortune becomes more important than the family's responsibilities. The business exists to preserve itself rather than serve customers, workers, and community. The institution serves its own continuation more than the good for which it was founded. Stewardship drifts into control, control into self-preservation, and self-preservation into extraction, while the language remains noble: responsibility, prudence, sustainability, long-term value.

This is why Protestant moral economy must be especially vigilant about self-deception. Stewardship is a beautiful word, which means it can conceal possession; vocation can conceal ambition; discipline can conceal anxiety or contempt. The corrective is not cynicism but accountability. Stewardship requires visible limits and governance that can say no to the steward, transparency about purpose, and practices of rest, release, and generosity. A steward who cannot be questioned has become an owner in the absolutist sense, and absolute owners, as every tradition in this book has recognised, tend toward extraction.

\section*{The Threshold of Commercial Society}

By the time we reach early modern commercial society, the traditions traced in this book have begun to converge into something rich and unstable. What they share, from Aristotle through Christianity and Scholasticism, through the Franciscans and Salamanca, through Luther and Calvin, is the conviction that economic life cannot be left to justify itself, that wealth requires ordering, and that the question of what commerce is for cannot be answered by commerce alone.

What they differ on, sometimes sharply, is how to answer that question under conditions of expanding markets, growing complexity, and the relentless ingenuity of people who would prefer not to have their financial arrangements subjected to theological scrutiny. That tension does not disappear. It intensifies.

Adam Smith stands at this threshold. He will inherit a world in which commercial society is no longer marginal. Markets are expanding, trade is becoming more complex, specialisation deeper, wealth growing. Moral philosophy must now ask not only whether individual acts of exchange are just, but whether commercial society as a whole can sustain the moral conditions it requires. Smith is often remembered as the founder of modern economics, but before he wrote about the wealth of nations he wrote about moral sentiments. He understood that markets depend on sympathy and judgment, on esteem and restraint, and on the desire to be worthy of respect. He did not imagine economic life without moral psychology.

The Protestant ethic brings us to that threshold with both a promise and a warning. The promise: that ordinary economic life can be a field of calling. The warning: that calling without limits can become a cage. Both will matter enormously for what comes next.

\chapter*{Chapter 8: Adam Smith and the Moral Psychology of Markets}
\addcontentsline{toc}{chapter}{Chapter 8: Adam Smith and the Moral Psychology of Markets}

Adam Smith stands at the threshold.

He is often remembered as the father of economics, the apostle of self-interest, the thinker who finally liberated markets from moral supervision. In this simplified version, Smith taught us that private greed produces public good, that markets work best when left alone, and that economic life can be understood without moral judgment. This Smith has served as a patron saint for a great deal of mischief.

The real Adam Smith was a moral philosopher before he was a political economist. His first great work was not The Wealth of Nations but The Theory of Moral Sentiments. He thought deeply about sympathy and judgment, about vanity and self-command, about virtue and corruption, and about the desire to be worthy of love and respect. He did not imagine human beings as isolated calculators of utility, wandering through life with tiny abacuses where their souls should be. He saw them as social creatures who long to be seen, approved, admired, and morally justified. Smith understood that commercial society depends on habits and sentiments it does not automatically create, and that it intensifies temptations it also cannot automatically restrain.

He belongs to both worlds at the threshold. Behind him lies the moral economy tradition: Aristotle's warnings about accumulation, Christian concern for avarice, Scholastic commutative justice, Franciscan analysis of productive capital, Salamanca's market reasoning, and Protestant vocation. Ahead of him lies modern economics, with its increasing abstraction, its specialisation, and its eventual separation from theology and moral philosophy. Smith himself did not make that separation. But he made it possible. Understanding why requires reading the Smith who has been forgotten alongside the Smith who became famous.

\section*{Sympathy and the Social Self}

Smith begins moral philosophy not with rules or contracts or utility, but with sympathy.

By sympathy, Smith does not mean pity alone. He means the human capacity to enter imaginatively into the situation of another person. We cannot literally feel another's pain or joy, but we can imagine what it would be like to stand where they stand; through this imaginative movement we become capable of approval and disapproval, of compassion and resentment, of admiration, and of moral judgment.

This produces an anthropology with profound economic implications. Economic actors are never merely preference-maximisers. They are persons seeking recognition and dignity, honour, belonging, and moral self-approval. Their choices are shaped not only by prices and incentives but by imagination and status, by shame and admiration, by habit, and by conscience. No one in a market is simply executing a transaction. Each person acts within a social world of esteem and judgment, and both parties leave the exchange slightly different from how they entered it.

This makes commercial society morally powerful in a double sense. It does not merely distribute goods; it distributes honour. It teaches people what to admire. A society that admires wealth without asking how wealth was acquired forms one kind of person. A society that admires stewardship, craftsmanship, truthfulness, and service forms another. The economy educates desire through admiration, and that education may be more consequential than any explicit moral teaching. Smith understood this with unusual clarity. Vocation can dignify work, but commercial society can redirect vocation toward status. Discipline can serve stewardship, but it can also serve the desire to be admired as successful. The market is not only an allocation system. It is a school of admiration, and what a society chooses to admire, its people tend to become.

\section*{The Impartial Spectator}

At the centre of Smith's moral theory is the figure of the impartial spectator.

Human beings judge themselves by imagining how a fair and well-informed observer would see them. We step outside our immediate passions and ask whether our conduct could be approved by someone not captured by our self-love. This imagined spectator disciplines vanity and anger, greed, and self-deception. It gives conscience a social and imaginative form. The impartial spectator is not identical to God, to law, or to public opinion. Public opinion can be corrupted, and crowds can admire what is unworthy. The powerful can surround themselves with flatterers, consultants, and public relations teams, which is much the same thing but with invoices. The image matters because it reveals how moral judgment operates within social life: we become capable of judging ourselves by learning to see ourselves from beyond ourselves.

A merchant tempted to deceive a buyer must ask how his conduct would appear if seen clearly. A lender profiting from distress must ask whether his gain would be approved by someone not benefiting from the transaction. An investor imposing extractive terms must ask whether the claim is proportionate to any genuine service. The impartial spectator asks the question self-interest avoids: what does this look like when I am not the centre of the story?

Modern markets often weaken this capacity structurally. Financial abstraction makes consequences invisible; complex ownership structures diffuse responsibility; supply chains separate decision-makers from workers; digital platforms separate designers from users' lived experience. The farther away consequences become, the harder it is for any spectator, however well-intentioned, to see them clearly. This is one of the moral dangers of financialisation. It creates systems in which many people can plausibly say they did not directly harm anyone. They merely optimised a portfolio or enforced a covenant, followed fiduciary duty, met market expectations, or improved shareholder returns. The harm is real, but the responsibility has been distributed so elegantly that no one can find where to stand in order to see it. Smith's moral psychology suggests that restoring the impartial spectator requires institutional design, not only personal conscience. Accounting must reveal externalised harm, governance must represent affected parties, and ownership forms must preserve mission. Moral imagination needs structural support.

\section*{Wealth and Moral Corruption}

Smith is deeply concerned with the way wealth distorts moral judgment, and it is one of his most important warnings.

Human beings have a tendency to admire the rich and powerful. We are drawn to their ease and apparent freedom, we imagine their lives as happier than they are, and we treat their success as evidence of wisdom while excusing their vices more readily than the faults of the poor. We confuse wealth with worth. When this happens at the level of a whole society, the results are serious. Capital no longer has to justify itself. The difference between productive and extractive gain becomes blurred. A steward and an extractor may receive the same applause, provided both arrive with enough money and an appealing narrative.

A society that admires wealth without moral discrimination will reward the appearance of success more than the reality of service. It will allow the powerful to define prudence and responsibility in ways that protect their own position, and it will make critique appear as envy while making the poor seem morally deficient rather than structurally disadvantaged. This is a subtle but devastating form of corruption, one that does not require open vice, only misdirected admiration. Modern capitalism intensifies this problem. Billionaires become public philosophers, exit value becomes proof of genius, and market capitalisation becomes a proxy for social importance. Philanthropy launders reputation, and companies that degrade attention, privacy, or labour can still be celebrated as innovative if they generate enough wealth. Smith would not have been surprised by any of this. He knew that spectators are not naturally impartial. A moral economy must therefore work at disciplining admiration itself, learning to honour right use rather than mere accumulation, service rather than scale, repair rather than disruption, wisdom rather than mere intelligence.

\section*{The Invisible Hand and Its Misuse}

No phrase in Smith has been more abused than ``the invisible hand.''

In popular memory, it has come to mean that self-interest automatically produces social good if markets are left alone: a principle used to sanctify market outcomes, minimise moral concern, and treat public intervention as presumptively suspect. This version of Smith has enjoyed a long and remarkably convenient life, especially among people for whom invisible hands tend to move in profitable directions. But it is a serious misreading.

Smith does believe that self-interest can generate unintended social benefits. The butcher, brewer, and baker provide dinner not from benevolence but from regard to their own interest. A commercial society can channel ordinary self-love into mutual service without requiring universal benevolence. This is a real and important insight that moral economy should not dismiss. Systems must work with ordinary human motives. But Smith knew perfectly well that self-interest does not always serve the common good. He warned explicitly that merchants and manufacturers often conspire against the public, seeking to rig markets in their favour through monopoly, mercantilist privilege, and political capture. The invisible hand is not a substitute for justice. It functions, when it functions, within a moral and institutional order that presupposes competition and law, trust, and some alignment between private gain and public benefit. When those conditions fail, self-interest produces not coordination but extraction.

Smith's real contribution is more subtle than the slogan suggests. He shows that economic systems can sometimes convert self-interest into mutual service, and that this is genuinely valuable. But he also gives us the framework for asking what institutional conditions enable that conversion and what conditions cause it to fail. A platform monopoly is not redeemed by having arisen through market processes; a housing bubble does not serve the common good because buyers and sellers responded to incentives; a surveillance business model is not made humane by the existence of formal consent. Smith gives us no permission to stop judging. He gives us reason to judge more carefully.

\section*{Division of Labour and the Narrowing of the Person}

Smith's famous praise of the division of labour is paired with a warning that his admirers have often overlooked.

Specialisation increases productivity in ways that benefit everyone. When tasks are divided, skill improves, time is saved, and goods become more widely available. The pin factory is Smith's classic illustration: many workers, each performing a simple operation, together produce far more than isolated individuals could manage alone. Pins may not sound spiritually thrilling, but a society that can make useful things efficiently is not to be despised.

But the same division that increases productivity can narrow the worker. A person who spends his entire working life performing a few simple operations may gradually lose intellectual, moral, and civic capacities. Work can become dulling and deforming. The worker may become less able to judge, deliberate, imagine, or participate fully in public life. Smith does not treat productivity as self-justifying. He asks what productive systems do to the people inside them. A society may become wealthier while making workers less whole.

This warning remains urgent. Modern specialisation can fragment responsibility in ways that parallel the pin factory's narrowing. An engineer may optimise metrics without seeing human effects; a warehouse worker may move through algorithmically managed tasks with diminishing autonomy; financial analysts may model returns without encountering the communities their models affect; content moderators may absorb trauma in hidden corners of digital infrastructure while value is captured elsewhere. The question is not whether specialisation is useful (it is) but whether the organisation of work allows human beings to flourish or progressively diminishes their capacity to do so. Smith saw that commercial society creates an obligation to care for the formation of people whose capacities have been narrowed by economic structure. A moral economy must go further and ask how firms can be designed so that work remains genuinely human.

\section*{Commercial Society and Virtue}

Smith does not simply condemn commercial society. He sees genuine virtues in it, and his realism is part of what makes his analysis useful.

Commerce can encourage prudence and reliability, honesty, and attention to the needs of others. A merchant who repeatedly cheats will lose trust and trade. Commercial interdependence can soften manners by making people useful to each other, and markets can move people away from some older forms of hierarchy, violence, and arbitrary domination. These are real goods. A world without markets is not automatically more humane. It may be more coercive, more stagnant, or more violent. Human beings have proved themselves admirably capable of injustice under every institutional arrangement so far attempted.

But commercial society also forms vices: vanity and envy, restlessness, indifference to the poor, cowardice before wealth, and the tendency to hollow out character while maintaining the appearance of respectability. It can reward pleasing others without serving them, and make people dependent on reputation while eroding the deeper qualities reputation is supposed to represent. Smith's value lies in holding both sides simultaneously. Commercial society is neither salvation nor damnation. It is an order of human cooperation with specific moral possibilities and specific moral dangers, and the question is always what kind of market, under what conditions, governed by what institutions, ordered toward what ends, and forming what kind of person.

\section*{Smith Against Monopolists and Rent-Seekers}

Smith is frequently invoked by defenders of business. He deserves to be invoked at least as often by critics of monopoly, rent-seeking, and corporate capture.

He was deeply suspicious of merchants and manufacturers when they sought political privilege. He understood that business interests regularly speak in the language of public good while pursuing private advantage, and he warned that proposals from those who profit from a particular trade should be examined with great caution. Monopoly and cartel behaviour, regulatory capture, restrictions on competition: none of these is an expression of market freedom. They are corruptions of it. Rent-seeking occurs when economic actors seek gain not by producing value but by controlling access, manipulating rules, capturing institutions, or extracting tribute from a privileged position. A platform that controls access to users, a patent holder who controls a life-saving drug, a dominant infrastructure provider who controls the terms on which others must operate: these are modern forms of what Smith identified as the persistent tendency of commercial power to rig the game in its own favour.

This matters for the argument of this book because a critique of extraction need not be anti-market. In many cases it is pro-market in the deeper sense, defending productive enterprise against rent-seeking, open competition against capture, and genuine service against tribute. Non-extractive capital can be presented not as hostility to markets but as a recovery of markets from the financial and institutional forms that have distorted them. Smith helps make that argument credible precisely because he was not an anti-market moralist. He was a moralist who took markets seriously enough to identify their internal enemies.

\section*{The Poor and the Problem of Visibility}

Smith's sympathy is tested most sharply by poverty.

He understands that the poor are often invisible to the moral imagination of the comfortable. Society admires the great and neglects the lowly. The suffering of ordinary people does not command attention the way the lives of the rich do. This asymmetry distorts judgment across the board. Wealth receives deference, poverty disappears from view, and moral analysis becomes calibrated to the concerns of those with enough status to be seen. If markets allocate attention as well as goods, then the poor are structurally disadvantaged in both. Their needs may be urgent but not profitable; their labour essential but not honoured; their vulnerability may become the basis for someone else's business model.

An economy that cannot see the poor cannot judge itself truthfully. A financial model that ignores vulnerability will mistake extraction for return, a market analysis that ignores power will mistake coercion for choice, and an innovation story that ignores hidden labour will mistake enclosure for genius. Smith's sympathy, properly extended, becomes a method of economic critique, asking persistently whose suffering is hidden by the price, whose labour is concealed by the brand, and whose community bears the cost of the investment.

\section*{Smith and the Older Moral Economy}

Smith is not a medieval Scholastic, and it would be false to read him as one. He does not reason like Aquinas, does not share the theological framework of the early Church Fathers, and does not condemn interest in the older way. He belongs to the Scottish Enlightenment, and that matters.

But he remains genuinely connected to the older moral economy tradition in ways his simplified legacy obscures. He shares Aristotle's conviction that wealth can distort the good life, the Scholastics' insistence that exchange depends on justice and that markets are not self-justifying, Salamanca's appreciation of market coordination and the informational role of prices, and the Protestant recognition that discipline and ordinary work are morally significant. He gathers much of the older tradition into a new language of commercial society, and his importance lies precisely in holding together what would later be separated.

Modern economics inherited parts of Smith while forgetting others. It took the butcher, brewer, and baker but left behind the impartial spectator. It took the invisible hand while leaving behind Smith's warnings against merchants who conspire against the public, and it took the division of labour as productivity while leaving behind the division of labour as deformation. It took self-interest but left behind the desire to be praiseworthy. Smith was narrowed into a symbol of market autonomy. The real Smith belongs to moral economy.

\section*{From Moral Philosophy to Political Economy}

Smith's work marks a genuine transition, and that transition carried real gains. Political economy as a distinct field of inquiry allowed more careful analysis of production and trade, of labour and wages, of money, and of the conditions of national wealth. It revealed patterns that moral exhortation alone could not see: that good intentions can produce bad outcomes, that restrictions can harm those they intend to protect, that markets can coordinate knowledge no authority possesses. A moral economy should welcome this analytical power. Economic literacy is not the enemy of moral judgment; it is one of its preconditions.

But the emergence of political economy also created a temptation: to treat the economic field as increasingly autonomous. Once economic phenomena could be studied in their own terms, it became easier to forget that those terms remained morally embedded. Once wealth could be studied scientifically, it became easier to avoid asking what wealth was for. Smith himself did not fully make that separation, but he made later separation possible because his political economy was so powerful. A thinker who understood the moral psychology of markets became remembered as the founder of a discipline that often bracketed moral psychology; a critic of monopolists became a mascot for business interests; a theorist of sympathy became a prophet of self-interest. The narrowing of Smith is one of the clearest signs of the great separation that follows him.

What Smith gives moral economy is irreplaceable: a realistic account of self-interest that does not require fantasies of universal benevolence; a theory of sympathy that insists economic life must be judged from beyond the standpoint of the self; a warning about the corruption of admiration; and a bridge showing that one can understand markets deeply without surrendering moral judgment. The Smith who matters is not the Smith of slogans but the moral philosopher of commercial society, the one who asked whether wealth is worthy of admiration, whether self-interest is disciplined by justice, whether the poor are visible, and whether commerce strengthens or weakens the moral life of those who inhabit it.

After Smith, that integrated vision begins to fracture. Political economy becomes more independent, moral philosophy more separate, and theological claims recede. The question ``what is an economy for?'' does not disappear immediately, but it begins to lose its place at the centre of the inquiry. The next chapter turns to that loss: the great separation in which economics breaks away from the moral and theological traditions that once gave it limits, language, and purpose.

\chapter*{Chapter 9: The Great Separation}
\addcontentsline{toc}{chapter}{Chapter 9: The Great Separation}

Adam Smith still stood with one foot in moral philosophy. After him, the balance began to shift.

Political economy became more confident as its own field. Production and wages, trade and taxation, rent, money, and growth all became subjects of systematic analysis. This development brought genuine intellectual power. It allowed thinkers to see patterns that older moral theology had often missed, exposed the unintended consequences of well-meaning rules, and showed how incentives shape behaviour, how specialisation increases productivity, how prices communicate scarcity, and how poorly designed intervention can harm the very people it seeks to help. Economics became powerful because it saw some things more clearly when it loosened itself from inherited moral categories.

But the loosening became a detachment. The older question, what is this economic activity for, became increasingly awkward. Economics sought firmer ground in observation and logic, in mathematics, and in prediction, increasingly distinguishing positive questions from normative ones: what is, rather than what ought to be. The result was the great separation: the gradual uncoupling of economic analysis from the older language of virtue and teleology, of theology and justice, of stewardship, and of the common good.

The separation brought real gains. It also brought losses that have not been adequately reckoned with. This chapter is about both.

\section*{Economics Breaks from Theology}

For much of the tradition this book has followed, economic reasoning belonged inside a wider moral order. Aristotle placed household management under ethics and politics; Christian moralists placed wealth under charity and mercy, under sin and salvation; the Scholastics judged exchange through commutative justice; the Franciscans asked whether capital served real production or exploited need. Protestant thinkers brought vocation, stewardship, and Sabbath into economic life, and Smith held political economy together with moral psychology. Economics was not yet autonomous.

The move toward a more independent economic science had understandable motives. Older moral frameworks could be rigid and paternalistic, hierarchical, or simply blind to complexity. Theological condemnation sometimes outran institutional understanding. Moralists could denounce markets without grasping what coordination requires, and rulers could use moral language to justify privilege or control, which is one of humanity's less charming habits. The economic world was becoming more complex, and inherited moral categories were often too blunt for the task of understanding how prices actually move, how incentives actually work, and how trade can generate benefits invisible from a local moral viewpoint.

But the distinction between positive and normative questions, genuinely useful as it is, can also mislead. Economic description is never wholly innocent. What we choose to measure, model, assume, and ignore already reflects judgments about what matters. A model that treats preferences as given has bracketed questions about desire formation; a firm theory built around shareholder value has not avoided teleology, it has chosen one and called it science. When economics broke from theology, it did not become free from ultimate questions. It became less explicit about them. The danger is not that economics became analytical. The danger is that it began to forget that analysis always takes place within a moral horizon.

\section*{Utilitarianism and the New Moral Arithmetic}

One of the most important forces in the great separation was utilitarianism.

Instead of grounding moral judgment in virtue or natural law, in divine command or the common good, utilitarianism proposed that actions and institutions be judged by their consequences for welfare, pleasure, pain, or utility. This made moral reasoning seem more calculable, more public, and more accessible across religious disagreement. One did not need to appeal to contested theological claims. One could ask what produced the greatest good for the greatest number. The strengths were real. Utilitarianism could challenge inherited privilege, expose cruelty hidden behind custom, support reform in law and public health, and force moral reasoning to care about consequences at scale. There is something morally bracing about asking not only whether an institution is venerable, but whether it actually makes people less miserable.

But utilitarianism also narrowed moral vision. Once the good becomes something to be aggregated, other questions become harder to ask. What may never be done, even for aggregate benefit? What forms of dignity cannot be traded off, and what goods are incommensurable? What forms of life are degraded even if preferences are satisfied? What happens when the preferences being counted have themselves been shaped by unjust systems? In economic terms, utility can replace richer accounts of human flourishing. People become bearers of preferences, and the deeper question of whether desires are well ordered becomes harder to ask. A society may maximise preference satisfaction while training people to prefer what harms them. It may produce aggregate gain while sacrificing vulnerable communities. Utilitarianism also struggles fundamentally with the category of enough. If more utility is always better, then limits must be justified by their consequences. Sabbath, restraint, and sufficiency lose their intrinsic moral force. The result is not necessarily cruelty. Often it is humane managerialism. But humane managerialism still cannot answer the question: what is an economy for?

\section*{Industrial Capitalism and the Abstraction of Labour}

The great separation was not only intellectual. It was material.

Industrial capitalism changed the scale, speed, and organisation of economic life. Production moved into factories, labour became specialised and measured and supervised, and workers were separated from tools, land, and household production. Time became wage time. Ownership and labour grew more sharply divided, and economic relations more mediated. A worker might no longer see the final product; an owner might no longer know the workers; a consumer might no longer know who made what they bought. The economy did not become less human. It became human at a distance, which is a much easier condition under which to forget people.

This brought real material gains. Industrialisation increased output and lowered the cost of many goods, expanded infrastructure, and contributed to improvements in health and material security. But perhaps most consequentially for moral economy, it created systems in which responsibility was harder to locate. Harm no longer arose primarily from individual unjust transactions. It arose from the structure of production itself. Who is responsible when wages are too low to live? Is it the employer, the market, the consumer, the investor, the structure of property? No single act of commutative injustice was easy to identify. Extraction had become systemic.

The abstraction of labour made this easier to hide. When labour is treated as a cost, the worker appears primarily as an input; when productivity is measured by output per unit of labour, the human person disappears behind efficiency; when ownership becomes a financial claim, the enterprise can be seen less as a community of persons and more as an asset. The moral economy tradition had always insisted that economic life consists of human acts in relation to other people. Industrial capitalism made those relations progressively harder to see, a problem that has only deepened in the age of global supply chains and digital platforms.

\section*{The Abstraction of Value}

Alongside the abstraction of labour came the abstraction of value.

Older moral economy thought about value in relation to use and need, justice, custom, and common estimation. Modern economics increasingly sought more formal and abstract accounts: price and utility, willingness to pay, marginal productivity. This was understandable. Economic systems needed concepts capable of explaining prices across many goods, places, and conditions. But abstraction becomes dangerous when it forgets what it abstracts from. Goods that cannot be easily priced become harder to see. Harms that are externalised appear outside the model. Forms of value that require time, relationship, or common life are discounted.

A forest valued for timber and water regulation, for biodiversity and sacred meaning and future generations, may enter the model only as monetisable services; the rest becomes invisible. A home valued as shelter, memory, and neighbourhood stability is treated by the housing market primarily as an asset, and when asset value dominates, the home's other meanings are systematically overridden. A company understood as a community of work and a mission-bearing institution is valued by financial markets primarily as expected future cash flows; that measure is not meaningless, but if it becomes sovereign, the enterprise's other forms of value become fragile.

The abstraction of value makes coordination easier and domination easier in the same motion. Once value is translated into a single metric, forms of life that resist the metric can be reorganised, traded, optimised, or sacrificed. Capital prefers commensurability because commensurability makes things comparable, liquid, and investable. But not all goods should be made liquid. Not all values can be faithfully translated into price without something essential being lost. A model that counts cash flows and ignores mission drift makes mission invisible. A model that counts user engagement and ignores addiction turns addiction into growth. One that counts rental yield while ignoring community displacement turns displacement into someone else's problem. The great separation is partly the story of metrics becoming masters.

\section*{The Rise of Homo Economicus}

Perhaps the most revealing symbol of modern economic abstraction is homo economicus: the rational, self-interested actor who seeks to maximise utility under constraints.

As a simplified model, homo economicus can be genuinely useful. Human beings do respond to incentives, do make tradeoffs, and do often pursue their own interests. Models need simplification. No one can build a model that includes every human motive, all childhood memories, seven forms of shame, three kinds of loyalty, and a sudden craving for soup. The problem arises when the model becomes an anthropology, when the simplification is mistaken for a description of what human beings actually are.

Human beings are not merely preference-maximisers. They are lovers and parents and citizens, workers and debtors, founders and stewards, and sinners. They seek meaning and honour, belonging and beauty, truth and forgiveness and admiration. They act from habit and fear, from love and shame, from duty and loyalty and hope. They often do not know what they want. They can prefer what harms them. Smith knew this; Aristotle knew this; so did Christian moralists and Protestant accounts of vocation. The model of homo economicus brackets most of it. Bracketing is not always wrong: it allows analytical focus. But if economics brackets too much for too long, the simplification becomes performative. If investors are assumed to seek maximum financial return, funds get designed around that expectation; if founders are assumed to seek exit, venture capital structures push toward exit; if firms are assumed to exist for shareholder value, governance structures make alternative purposes fragile. The model makes the world more like itself.

\section*{The Case for Separation}

Before judging the great separation too harshly, we should make the case for it honestly.

Older moral economies had real failures. They could be paternalistic and static, hierarchical, anti-commercial, and blind to the genuine benefits of market coordination. They could protect insiders and guilds under the language of justice, condemn practices they did not understand, confuse custom with morality, and preserve unjust social orders by sacralising them. Economic analysis needed room to observe how prices actually move, to understand incentives without first condemning self-interest, and to show that restrictions justified in moral language could harm the poor. The separation also made public reasoning easier in plural societies. If people disagree deeply about theology, virtue, and the highest good, economic policy cannot always wait for metaphysical agreement. Societies need ways to discuss taxation and trade, welfare, money, and regulation across deep moral disagreement. Economics offered a shared language of incentives and costs, of benefits and tradeoffs and outcomes. This was useful and remains useful.

A recovered moral economy should not discard modern economics. It should not pretend that Scholastic categories alone can guide central banking, antitrust, digital infrastructure, climate transition, or global trade. It should not replace analysis with righteous sentiment or condemn market outcomes it dislikes without understanding the mechanisms at work. The great separation taught us that intentions are not enough, that systems have consequences beyond moral aspiration, that scarcity is real, that incentives matter, and that tradeoffs cannot be wished away. These are genuine gains. A moral economy that ignores them will fail.

\section*{Why the Losses Were So Serious}

And yet the gains were purchased at too high a price when economics came to treat moral and theological questions as external to its core.

The problem was not that economics became analytical. The problem was that analysis became detached from purpose. Incentive analysis often bracketed the formation of desire; price became a substitute for judgment; self-interest became anthropology; growth became telos. The deepest questions were pushed to the margins. What kinds of people does an economy form? What limits should it honour? When does return become extraction, and when does voluntary exchange conceal domination? When does capital stop serving enterprise and begin ruling it? These were assigned to ethics, politics, religion, or personal opinion. But economic life did not stop answering them. It simply answered them through institutions, metrics, and capital structures, and the answers appeared technical rather than moral, which made them harder to challenge.

When shareholder value governs the firm, when GDP governs public success, when exit value governs startup finance, when engagement governs digital platforms, these are not technical descriptions. They are moral choices made in the language of data. The great separation made it easier for capital to present its claims as neutral, inevitable, or scientific. It made it easier to hide extraction behind efficiency, to treat economic power as market outcome rather than moral relation, to regard institutional design as technical rather than ethical, and to forget that the economy is a human order rather than a natural force. The gains of economics should have been integrated into moral economy. Instead, too often, they displaced it.

\section*{The Hidden Normativity of Neutrality}

Modern economics often presents itself as positive rather than normative: it analyses behaviour, predicts outcomes, and evaluates tradeoffs, while leaving normative judgment to others. There is wisdom in this humility. Economists should not pretend that technical expertise gives them final authority over the good life. History has not suffered from a shortage of experts willing to tidy up human existence on everyone else's behalf.

But neutrality can become a mask. Assumptions are never neutral, and neither are metrics. The decision to model the firm as a nexus of contracts rather than a community of work is not neutral. The decision to treat nature as input, labour as cost, and capital as primary claimant is not neutral. The decision to measure welfare through preference satisfaction is not neutral. Even the refusal to speak morally has moral consequences. If economists decline to judge whether a market outcome is just, the outcome still stands. If they decline to ask whether bargaining conditions are coercive, consent remains formal. If they decline to ask whether capital's claim is proportionate, return appears as a market fact. Neutrality often stabilises the existing order.

This does not mean economics is propaganda. It means that any discipline studying human action must be honest about its moral horizons. Economics can tell us a great deal about means. But when it refuses to discuss ends, it does not eliminate ends; it leaves them to whoever has the power to build them into institutions. In modern capitalism, that power often belongs to capital. This is why the great separation matters: once moral language is removed from the core of economic reasoning, capital can more easily define success in its own terms. Return becomes performance, growth becomes health, exit becomes success, and market price becomes truth. The question of whether capital serves life becomes harder to ask in any authoritative register.

\section*{What Was Lost: Teleology, Limits, and the Firm}

The deepest loss was teleology.

Older economic thought, for all its limits, asked what economic activity was ordered toward. The household was for life; the city, for the good life; property, for use and stewardship; exchange, for mutual need; work, for vocation and service. Modern economics often avoids such claims. It may ask what individuals want without asking whether wants are worthy, may study how firms maximise without asking what firms are for, may model how markets clear without asking whether the goods being traded should be governed by market logic. The loss of teleology produces confusion about limits. If we do not know what something is for, we do not know when it has enough. A company ordered toward mission can judge when growth endangers purpose. But if the operative end is only more (more growth, more return, more engagement, more scale) then limits appear irrational. There is always more return to seek and more data to capture, more attention to monetise, more dependencies to exploit. The system continues to call this progress because it lacks a language for disorder.

A related loss was the moral language of limits on capital. The older traditions had shared a suspicion of gain without service, accumulation without limit, and lending that profited from need. They feared usury and avarice, monopoly and unjust price, and wealth without responsibility. Modern finance became extraordinarily skilled at multiplying claims (debt and equity, derivatives and securitisation, leverage and management fees, liquidation preferences, platform rents, data extraction) without a moral account of when claims become excessive. The older usury debates were not technically adequate for modern finance, but they preserved a question still urgently needed: when does a financial claim become morally disproportionate? A moral economy needs proportionality: the principle that return should correspond to contribution and risk, to time and responsibility, and should not be infinite simply because contract permits it.

The third loss was the moral meaning of the firm. In older moral economy, economic activity was understood through households and crafts, through guilds and communities. The modern firm is more powerful than any of these predecessors, shaping work and technology and wages, political influence, culture, and ecological impact on a scale that earlier traditions could not have imagined. Yet modern economic theory has struggled to give the firm a rich moral meaning. If it is seen primarily as a nexus of contracts, its communal character fades. If seen primarily as an asset owned by shareholders, then workers, customers, communities, and mission become secondary. Shareholder primacy is not merely a governance theory; it is a teleology of the firm. It answers ``what is the company for?'' with ``financial return.'' That answer is now built into law and compensation culture, into capital markets and exit expectations. A recovered moral economy must challenge it, not by denying the legitimacy of return, but by insisting that the firm is a human institution whose governance must be judged by more than efficiency and shareholder value.

\section*{What Was Lost: The Commons}

The fourth loss was the commons.

Older societies preserved many forms of shared use, customary rights, communal obligation, and common goods. Modern economic thought simplified the landscape into market and state, private and public, and in doing so, made the commons harder to see. Many of the goods most necessary for human flourishing are common or commons-like: language and knowledge and culture, ecological systems, public trust, scientific research, civic institutions, open-source software, digital protocols, public health, and shared infrastructure. These goods are not well governed through private ownership or state administration alone. They require stewardship. But when value is recognised primarily through price, private ownership, or trade, commons appear as inefficiencies or under-monetised assets. Capital then sees opportunity: enclose and meter and license, extract, platformise, monetise. Open-source software becomes training data, community labour becomes product input, and public knowledge becomes private moat.

The moral economy tradition gives resources to resist this pattern. Aristotle asks whether shared goods serve life or become instruments of accumulation. Christian thought insists that ownership is accountable to use and need. The Scholastics judge exchange by justice rather than price alone. Protestant stewardship demands the preservation of entrusted goods for their proper purpose, and Smith warns against concentrated power and rent-seeking. A recovered moral economy must bring the commons back into view, especially as the digital age creates new forms of shared productive capacity that are rapidly being enclosed.

\section*{From Separation to Sovereignty}

The great separation did not merely produce a more autonomous academic discipline. It helped create conditions under which capital could become increasingly sovereign.

When economics no longer asks what capital is for, capital answers for itself. When firms are defined by shareholder value, capital gains governing priority. When prices are treated as sufficient signals of value, capital can buy what should be stewarded. When human beings are modelled as preference-maximisers, capital can shape preferences and call the result demand. When moral limits appear external to economic logic, capital can treat them as constraints to be optimised around. Capital becomes sovereign not by declaring itself king, but by becoming the default measure of institutional success.

Universities reorganise around funding and market relevance. Hospitals reorganise around financial throughput, cities around real estate value, media around attention monetisation, and public infrastructure around private capital and proprietary systems. The logic spreads because so many institutions must ultimately justify themselves in capital's language. The older moral economy tradition sought to keep capital subordinate: subordinate to household and city, to virtue and charity, to justice and vocation, to the common good, and to God. The great separation made that subordination harder to articulate, and capital could then appear as neutral resource, as rational allocator, as objective market judgment, or as necessary discipline.

The next part of this book examines what happens when that process deepens: what extraction actually means, how corporations became sovereign, how financialised capitalism has reorganised economic life, and what institutions might yet recover the older question. What is capital for, and what must it never be permitted to do?

\chapter*{Chapter 10: What Extraction Means}
\addcontentsline{toc}{chapter}{Chapter 10: What Extraction Means}

We have now reached the hinge of the book.

The first part of the journey recovered a language. Traditions stretching from Aristotle through Christianity and the Scholastics, the Franciscans and Salamanca, Protestant thought, and Adam Smith were each, in their different ways, circling the same problem: the difference between wealth that serves life and wealth that feeds on it. The great separation showed us what was gained when economics became more analytical, and what was lost when it became less able to ask what economic life is for.

Now we must name the thing these traditions were circling.

Extraction.

The word is used often but not always carefully. It can become a slogan, meaning whatever the speaker dislikes: profit and capitalism, markets and finance, inequality, rent, bad vibes in a boardroom. Some of these are connected to extraction. None is identical with it in every case. If the concept is too broad, it becomes useless. Commerce becomes extraction, profit becomes extraction, investment becomes extraction. At that point, we have not made a moral distinction. We have simply thrown a net over half of economic life and called the result analysis.

But moral economy depends on distinction. It must be able to say not only that something is wrong, but why this kind of gain differs from legitimate service, risk-bearing, and stewardship. Extraction is not the same as profit, which can be legitimate when it corresponds to real contribution, service, risk, and coordination. It is not the same as markets, which can serve life when they coordinate useful activity and allow people to cooperate freely. Nor is it the same as growth, which can serve life when it means more food and shelter, more health and knowledge and shared capacity. The task of this chapter is to give extraction a working definition: not a final metaphysical formula, engraved on stone and guarded by economists in robes, but a moral-economic diagnostic that can help us recognise when capital has stopped serving life and has begun feeding on it.

\section*{A Working Definition}

Extraction is the appropriation of value from persons, communities, institutions, ecosystems, or productive enterprises in ways that are disproportionate to contribution and sustained by dependency, domination, enclosure, scarcity, opacity, or structural power.

Several things are embedded in this definition. Most obviously, extraction involves appropriation: something is taken, captured, claimed, or siphoned away. It is not merely inequality of outcome, but a relation in which one party gains by drawing value from another source. The value extracted may take many forms (money and labour, land value and data, attention, knowledge, ecological capacity, institutional trust, or future cash flow). Modern extraction often operates precisely by converting non-financial forms of value into financial claim. It turns attention into advertising revenue, community into platform value, public infrastructure into private moat, and ecological stability into someone else's unpaid bill.

Most importantly, extraction is disproportionate to contribution. This is the moral heart of the matter. A person or institution may legitimately receive value for real service, labour, risk, coordination, creativity, or stewardship. Extraction begins when the claim exceeds the contribution, especially when the excess is sustained by power rather than service. It occurs through structures that limit the other party's freedom, visibility, alternatives, or bargaining power: through dependency, domination, enclosure, scarcity, opacity, or structural power.

Crucially, extraction can be legal. Many extractive arrangements are perfectly lawful, contractually agreed, institutionally normalised, professionally managed, and publicly celebrated. They come with signatures and dashboards, with compliance review, and sometimes with very handsome annual reports. The legal form of consent does not settle the moral question. Extraction can also occur inside productive systems. A company may produce useful goods and still be governed extractively; a platform may provide real services and still extract from dependency. Extraction therefore requires judgment, not from sector labels or moral branding, but from asking how the gain is produced, who bears the cost, who controls the relationship, whether alternatives are real, and whether the activity strengthens or depletes the life from which value is drawn.

Extraction is gain without right relation.

\section*{Productive and Extractive Capital}

The distinction between productive and extractive capital is not a distinction between good people and bad people, nor between virtuous sectors and vicious ones. Manufacturing can be extractive, and finance productive; technology can serve life, and agriculture can exploit. A nonprofit can drift toward extraction, and a commercial business can practise stewardship. The world, annoyingly but accurately, refuses to sort itself into morally convenient buckets.

The distinction concerns the relationship between capital, contribution, power, and purpose. Productive capital participates in the creation, preservation, or circulation of real goods. It supports useful work, shares risk in proportionate ways, strengthens the productive capacity of the enterprise or community, and accepts limits. Extractive capital captures value without proportionate contribution. It gains through control and dependency, through scarcity and legal privilege, through informational advantage, or through distress. It may support productive activity initially, but over time its claims dominate, weaken, or redirect that activity toward the needs of capital itself.

A simple example clarifies the difference. A lender provides funds to a small manufacturer to purchase equipment on terms repayable from future cash flow, at a reasonable rate, with flexibility in hardship, and without giving the lender control beyond what is necessary to protect repayment. If the equipment improves productivity and repayment remains proportionate, the capital serves life. Now imagine a different structure: the same loan includes covenants likely to be breached under ordinary stress, upon which the lender can impose fees, seize assets, or force restructuring. The manufacturer's vulnerability becomes the lender's opportunity. The financial claim does not merely support production; it positions itself to extract from distress. Both arrangements are loans. Only one is life-serving. The difference is not sentiment. It is architecture: the actual structure of who controls and who bears risk, who receives surplus and who can force exit, and who remains responsible after extracting value.

\section*{The Genealogy of the Distinction}

The distinction between productive and extractive capital is modern in language but ancient in structure. What this book has been tracing, chapter by chapter, is a long and varied tradition of people noticing the same problem from different angles.

Aristotle gives us the first grammar: oikonomia is acquisition ordered toward life, and chrematistics is wealth-getting without limit. The moral danger lies in the inversion of means and ends, money being useful as a tool but becoming the purpose. Christianity deepens the concern through avarice and usury, through mercy and charity, making wealth dangerous not only because it destabilises the city but because it can deform the soul and bind the vulnerable to the powerful. The Scholastics formalise the distinction through commutative justice and the just price: profit must correspond to service, contract does not exhaust justice, and a voluntary transaction can still be disordered if it exploits necessity or unequal power. The Franciscans clarify that capital can be productive when used in enterprise and foresight, but that this recognition requires sharper discipline, since capital that can be productive can also disguise extraction as participation.

Salamanca expands market reasoning while revealing that coordination is not the same as justice. Protestant vocation and stewardship bring economic activity under calling, but when detached from Sabbath and grace, discipline and reinvestment can become restless accumulation. Smith gives us moral psychology: wealth corrupts admiration, self-interest requires justice, and rent-seekers manipulate law for private gain. Modern economics then provides analytical power while often weakening the moral vocabulary of limits. As teleology recedes, the distinction between productive and extractive capital becomes harder to state. Extraction remains visible in suffering but harder to name in any authoritative register.

The genealogy matters because it shows that extraction is not a new concern invented by modern critics of capitalism. It is the old problem of disordered wealth in contemporary form: money detached from use, gain detached from service, return from risk, ownership from responsibility, growth from flourishing.

\section*{From Chrematistics to Veblen}

To connect the ancient and modern versions of the distinction, Thorstein Veblen is especially useful.

Veblen distinguished between industry and business in a way that echoes Aristotle's distinction between economy and chrematistics. Industry concerns the productive process: making goods and applying knowledge, solving problems, coordinating labour, and meeting material needs. Business, in Veblen's critical sense, concerns pecuniary gain: the strategic pursuit of profit through control, restriction, monopoly, and the manipulation of ownership claims. The logic of production and the logic of pecuniary gain can diverge, and when they do, business can actively sabotage industry.

A factory may be capable of producing useful goods, but owners restrict output to maintain prices; a technology may be capable of serving many, but access is enclosed to maximise revenue; a company may have productive workers and loyal customers, but financial owners extract cash, cut capacity, and sell assets to increase returns; a platform solves a coordination problem and then uses its position to collect rents from those who now depend on it. In each case, capital that should serve production instead interferes with production in order to increase financial claim. Value is not created by making more or serving better; it is captured by controlling the conditions under which others must produce, exchange, or access what they need.

Veblen also helps us see why extraction can be rational inside the wrong system. If the governing aim is pecuniary gain, then restricting output, manufacturing scarcity, or lobbying for advantage may be intelligent. The problem is not irrationality. It is the rational pursuit of a disordered end. Modern extraction often looks sophisticated rather than primitive, modelled in spreadsheets and written into contracts, justified by fiduciary duty, and celebrated as strategic positioning. This is precisely why we need moral economy. Without it, we may mistake extraction for intelligence.

\section*{Rent, Usury, and Enclosure}

Three older categories help illuminate the grammar of extraction: rent, usury, and enclosure. Each concerns a different way of gaining from control rather than contribution.

Rent, in its problematic sense, is income derived from control of access to something needed, rather than from ongoing productive contribution. Not all rent is illegitimate; someone who maintains a building or stewards land may receive legitimate compensation. But rent becomes extractive when it arises from control over scarcity, monopoly, or exclusion without proportionate service. Usury, in the older moral sense, is gain from lending that exploits need or treats money as capable of generating claim without real participation in risk, labour, or production. Not all interest is usury, but usury names the danger that financial return becomes detached from contribution and sustained by the borrower's vulnerability. Enclosure is the conversion of shared or common goods into controlled assets from which access can be monetised. Not every boundary is unjust, but enclosure becomes extractive when shared capacity, public knowledge, ecological systems, or digital commons are privatised in ways that allow one actor to profit from what many created or need.

These three patterns often overlap and reinforce each other. A platform encloses a network and then charges rent for access. A private equity firm acquires a productive company, loads it with debt, extracts fees, and sells the assets. A technology company trains its models on shared knowledge, encloses the resulting infrastructure, and charges access back to the public whose labour helped create the value. A dominant marketplace controls sellers' access to buyers, then raises fees once dependency is established. Rent, usury, and enclosure are not relics. They are living forms of extraction: ways capital gains by controlling the conditions of life and production rather than serving them.

\section*{Extraction and Power}

Extraction is inseparable from power, and this connection is often hidden by the language of voluntary exchange.

A contract may appear equal because both parties sign. A platform may appear consensual because users click. An investment may appear chosen because a founder accepts the terms. But power shapes the field in which choices are made. A founder with three months of runway negotiates differently from a fund with diversified capital; a tenant facing eviction bargains differently from a landlord with multiple properties; a small supplier dependent on one platform bargains differently from the platform that controls access to customers. Extraction occurs when power turns another's lack of alternatives into a revenue stream.

This is why formal consent cannot be the only moral test. The question is not simply whether someone agreed, but what made this agreement their best available option, and who benefits from the absence of better options. This does not mean every unequal exchange is extractive. Inequalities of skill, information, capital, and urgency are common in economic life, and asymmetries do not automatically make exchange unjust. But they create moral responsibility. The more power one party holds, the more careful it must be not to convert that power into disproportionate gain. Power does not make profit illegitimate. It makes profit accountable.

\section*{Extraction and Dependency}

Dependency is not inherently bad.

Human beings are dependent creatures: on parents and neighbours, on farmers and teachers and doctors, on institutions and ecosystems, and, if we are honest, on many people whose names we will never know and whose labour we are inclined not to notice. Societies are webs of interdependence, and markets themselves are forms of organised interdependence. The problem is not dependency as such. The problem is dependency without reciprocity, accountability, or exit.

Life-serving interdependence strengthens the parties involved. It allows cooperation and specialisation, trust, and flourishing. Extractive dependency weakens one party in order to secure the other's claim. It makes exit difficult, monetises lock-in, creates recurring need, and profits when the dependent party cannot leave, cannot understand, cannot substitute, or cannot refuse. A software platform may begin by solving a real problem. Users join because it is useful, developers build around it, and businesses integrate it. Over time, switching becomes costly, data accumulates, workflows harden, and network effects deepen. The platform can then change terms and raise prices, alter visibility, impose fees, or capture data. The service may still be real, but the dependency has become monetisable: extraction through lock-in. A moral economy must distinguish interdependence from dependency capture. Capital serves life when it increases the real agency and resilience of those it touches. It feeds on life when it profits from keeping them dependent.

\section*{Extraction, Opacity, and Externalisation}

Extraction thrives in opacity. If people cannot understand the terms, they cannot judge them; if consequences are hidden, moral imagination fails; if ownership structures are layered, responsibility diffuses; if algorithms govern visibility, affected parties may not know why they succeed or fail. Transparency is a precondition for accountability, and moral economy requires that people be able to know what claims are being made upon them, who controls the relationship, how value is allocated, and what alternatives exist.

Externalisation is extraction's most common mechanism. To externalise is to shift costs onto others while retaining benefits for oneself. A factory pollutes a river without paying for restoration. A platform relies on unpaid moderation and on public infrastructure and open-source software without contributing proportionately. A fund extracts cash from a company and leaves future fragility to employees, creditors, or the public. Externalisation allows extraction to masquerade as efficiency: costs disappear from the balance sheet, harms are moved elsewhere, and market prices fail to reflect the full cost. The firm appears more productive than it is because someone else is paying the real bill. A business model that depends on hidden harm is, in a morally serious sense, misrepresenting its own profitability.

\section*{Extraction and Time}

Extraction often works by distorting time.

Productive activity usually requires patience. Soil fertility takes years, and so does trust, craft, research, open-source infrastructure, and mission fidelity. Extractive capital tends to accelerate claims on value. It seeks liquidity before maturity, pulls future cash flows into present valuation, and pressures companies to grow before they are stable; it discounts long-term harms and rewards exit over stewardship. It treats future resilience as underutilised capital, and prefers what can be monetised now to what must be cultivated slowly.

This is why the time horizon of capital is morally important. A patient investor supports a company through the slow work of becoming durable; an impatient investor forces growth that weakens the company before it can survive pressure. Extraction compresses time in favour of the claimant. It asks the future to pay for the present, asks workers to absorb pressure now so investors can receive returns sooner, and asks ecosystems to bear depletion so profits can be reported this quarter. A moral economy must recover patience as a financial virtue. Evergreen finance and repayment from durable surplus, mission-locked ownership, and commons governance are all attempts to repair distorted time, to say that not all goods should be forced into the timetable of liquidity.

\section*{The Category of Enough}

Extraction cannot recognise enough.

This may be its most Aristotelian feature. Natural acquisition is limited by the needs of life; chrematistics is unlimited because its end is accumulation. Modern extraction follows the same logic. If the governing aim is return, growth, liquidity, or market dominance, then every limit appears as inefficiency. A company has enough revenue to sustain itself, its workers, customers, and mission, but investors demand growth. A platform serves its users well, but must increase engagement. A landlord receives fair compensation, but market scarcity allows more. A fund has received a reasonable return, but uncapped upside permits continued claim. In each case, enough is overridden by more, because the structures in place reward expansion regardless of whether expansion serves anything.

This is why limits must be designed rather than left to sentiment. A person may intend restraint, but structures reward expansion. A founder may want mission, but capital demands exit. Repayment caps define when capital has received enough; asset locks prevent mission assets from being converted into private gain; steward ownership separates control from speculative ownership; commons licences define conditions under which shared goods may be used without enclosure; Sabbath places a limit on production and mastery; Jubilee interrupts debt's claim to perpetuity. Without such limits, extraction will rationalise itself endlessly. Enough is not anti-economic. It is the condition under which economic life remains ordered toward life rather than consuming it.

\section*{Extraction as Anti-Stewardship}

Extraction is the opposite of stewardship.

A steward receives a good in trust. The good is not merely possessed; it is entrusted for a purpose. The steward must preserve, cultivate, and use it rightly, remaining accountable to the giver, to the purpose of the good, to the community affected by it, and to those who come after. Extraction breaks this relationship entirely. It treats entrusted goods as opportunities for claim. It consumes inheritance and monetises dependency, liquidates mission, and converts common goods into private assets. It captures surplus without accepting corresponding responsibility, and it uses legal ownership to deny moral accountability.

Where stewardship accepts accountability and honours limits, extraction seeks insulation and treats limits as inefficiencies. Stewardship preserves mission; extraction converts mission into asset. Stewardship strengthens the life of the other; extraction feeds on it. This is why a recovery of stewardship must be structural. It is not enough to ask capital to behave nicely, as though extraction were merely a personality flaw to be addressed with better manners. Capital must be placed under forms of governance, ownership, accounting, and law that make stewardship durable and extraction difficult.

\section*{How to Recognise Extraction}

Extraction is not always obvious. It often appears under respectable names: efficiency and innovation, liquidity, growth, optimisation, fiduciary duty, market pricing, risk management, user engagement, asset utilisation, value creation. The vocabulary is familiar, polished, and usually wearing a lanyard.

A useful diagnostic begins not with categories but with questions about the actual structure of the relationship. What is the source of gain: real service, production, risk-bearing, coordination, or stewardship, or is it control over access, scarcity, dependency, distress, opacity, or legal privilege? Is the claim proportionate to contribution, or has it outgrown the service that justified it? Who bears downside risk? Are risks shared fairly, or shifted onto workers, borrowers, founders, communities, and ecosystems while capital captures upside? Are alternatives real, in the sense that the other party can refuse, exit, substitute, or bargain, or is consent formal but structurally constrained? What is invisible in the accounting? And can the system recognise enough?

These questions do not produce automatic answers, and they require judgment. But they shift analysis from surface form to moral structure. A loan is not extractive because it is a loan. A rent is not extractive because it is rent. A profit is not extractive because it is profit. The question is always what kind of relationship is created, and what kind of life it produces.

\section*{Why the Distinction Is Resisted}

The distinction between productive and extractive capital is not merely forgotten. It is resisted, because it threatens powerful interests.

If we can distinguish productive from extractive gain, then not all profit receives the same honour. Not all investment appears equally legitimate. Not all legal contracts appear just. Not all shareholder returns appear morally deserved. This is uncomfortable for those who benefit from ambiguity. It is useful for extractive actors to blur the distinction between enterprise and extraction, to appear as entrepreneurs and innovators, as investors and job creators, as efficiency experts or platform builders, when in fact they are engaged primarily in rent-seeking, enclosure, or disproportionate claim. The language of value creation is especially vulnerable to this blurring. A private equity firm may claim to create value when it increases financial returns through layoffs, underinvestment, and debt loading. A technology company may claim to create value when it encloses public knowledge in proprietary systems. The question that cuts through the language is always: value for whom, created how, and at whose expense?

A moral economy must reclaim the language of value creation from financial abstraction. Real value creation strengthens life. Extraction may increase financial value while reducing the goods that actually matter.

\section*{Extraction in the Firm}

The firm is one of the main sites where productive and extractive logics collide.

A firm can be a community of production, gathering people around a useful purpose, coordinating labour, knowledge, capital, and trust across time, serving customers, sustaining workers, and contributing to community life. But a firm can also become a vehicle for extraction. This happens when the firm's productive life is subordinated to financial claim. Workers become costs to minimise, customers become revenue streams to lock in, suppliers become pressure points, mission becomes marketing, and assets become collateral. The company still produces. It may even produce more efficiently for a time. But its purpose has shifted from serving life to feeding on it.

Extraction in the firm often begins with a change in the answer to one question: what is the company for? If the answer is to serve a real mission through useful work, then capital must be ordered around that purpose. If the answer is to maximise financial return to owners, then mission survives only when compatible with return. This is why governance matters, and why the purpose of the firm cannot be protected by rhetoric alone. It must be embedded in ownership rights and board duties, in investor agreements, in exit restrictions, and in surplus allocation. The chapters that follow examine what those structures look like when they are designed to serve life rather than extract from it.

\chapter*{Chapter 11: Austrian Economics and Entrepreneurial Discovery}
\addcontentsline{toc}{chapter}{Chapter 11: Austrian Economics and Entrepreneurial Discovery}

If the previous chapter named extraction, this one must do something equally important: protect moral economy from becoming anti-market.

A critique of extraction can easily slide into a critique of markets as such. If profit can be extractive, perhaps profit is inherently suspect; if price can conceal power, perhaps price is merely ideology; if capital can dominate, perhaps decentralised exchange is only a mask for domination. That temptation must be resisted. Markets are not morally self-sufficient, but they are not morally meaningless. They are one of the ways human beings coordinate knowledge and labour, need and risk, creativity, and time without a single centre of command. They can reveal things no planner knows, and they can allow ordinary people to respond to one another's needs through enterprise, trade, and adaptation. At their best, markets are a vast and unruly conversation conducted through prices and experiments, through failures and adjustments, through substitutions, discoveries, and the occasional alarming spreadsheet.

The Austrian tradition helps us see this. Carl Menger shows us that value is not simply embedded in objects but emerges from human need, judgment, and context. Friedrich Hayek shows us that knowledge is dispersed, tacit, local, and often impossible to centralise. Austrian accounts of entrepreneurship reveal the market not merely as a static allocation mechanism but as a process of discovery under uncertainty. These insights matter because a moral economy must not confuse moral seriousness with central control. But discovery must be ordered. The fact that markets can discover opportunities does not tell us whether those opportunities are good. The fact that entrepreneurial profit signals value does not tell us whether that value serves life. This is the beauty and the limit of Austrian economics: it teaches us that markets are epistemically powerful. It does not, by itself, teach us what markets are for.

\section*{Menger and Subjective Value}

Carl Menger's account of value begins from human need.

Goods have value because they are capable of satisfying human wants, and that value is not a property sitting inside the object like jam in a doughnut. It arises from the relationship between a person, a need, and a good capable of meeting that need. The same object may have radically different value in different circumstances: a cup of water matters little beside a clean river and enormously in a desert. Menger's marginalism refined this further: the value of a good depends not on the total usefulness of all goods of that type, but on the significance of the particular need that would go unmet if one unit were lost. This helps explain why water, necessary for life, may command lower prices than diamonds. Price and existential importance are not the same.

This was a genuine advance, and moral economy should welcome it. Human beings are not interchangeable units with identical needs, and value is not fully knowable from outside. A centralised authority may not know what a particular tool means to a craftsperson, what a loan means to a household in distress, or what a specific technology makes possible for an entrepreneur at a critical moment. Subjective value protects humility. It makes room for decentralised judgment and reminds us that people closest to a situation often know something outsiders do not.

But subjective value is not sufficient. The fact that a person values something does not settle whether the desire is well ordered. Human beings can value what harms them. Addiction creates value in the economic sense; manipulation creates demand; artificial scarcity increases willingness to pay; desperation raises the subjective value of access. A person may value relief from distress so highly that they accept terms that deepen their dependence. Subjective value explains price. It does not justify every price. Moral economy must supplement this insight by asking what kind of relationship forms the value in question. Menger helps us see that value is relational. Moral economy must then ask whether the relation is one of genuine service or of manufactured need.

\section*{Prices as Knowledge}

One of the great Austrian insights is that prices communicate knowledge.

A price is not merely a number. It condenses dispersed information about scarcity and demand, opportunity, cost, urgency, substitution, and expectation. No single mind needs to understand why tin has become more expensive for users of tin to start economising its use. The price itself communicates a change in relative scarcity, and people adjust. Producers search for alternatives, consumers substitute, and entrepreneurs notice opportunities. Vast coordination occurs without central command. Hayek made this central. The economic problem, for him, is not merely how to allocate given resources to given ends. The deeper problem is how to use knowledge that is dispersed among many people, much of it tacit, local, time-bound, and often discovered only through action. A farmer knows the soil; an engineer knows the failure mode; a founder knows the user pain. Markets coordinate this dispersed knowledge through prices and profit, through loss and entrepreneurial action, in ways that no planning committee can replicate, however large the committee, however confident the chair, and however heroic the catering budget.

But prices are not omniscient. They communicate certain kinds of knowledge and omit others. They can reveal scarcity but not always justice; they can signal demand without telling us whether demand has been manipulated; they can reflect costs borne by the seller but not costs externalised onto workers, ecosystems, public institutions, or future generations. If a river can be polluted without cost, the price of the product lies. If open-source maintainers are unpaid, the price of proprietary software built on their work lies. If tenants have no alternative, rent reveals dependency and power rather than justice. Prices are signals, not sacraments. Moral economy must read them critically, taking them seriously as real information while asking what conditions produced them and whose experience they leave out.

\section*{Entrepreneurship as Discovery}

Austrian economics gives entrepreneurship a special and important role.

The entrepreneur does not merely manage known resources. She discovers possibilities, seeing mismatches between present arrangements and future opportunity, noticing unmet needs and emerging technologies, new combinations, or hidden demand, and acting under uncertainty before others see what she sees. This connects directly to the Franciscan discussion of foresight. Olivi recognised that commercial foresight can be a real contribution, and Austrian economics develops this into a broader account of market discovery. Profit is not merely reward for possession. It can be reward for seeing and acting before others do. A moral economy should not crush discovery under suspicion. Many life-serving goods emerge because someone noticed what others missed: a new tool, a better process, a more humane service, a way to coordinate unused resources. Entrepreneurship can serve life by revealing possibilities hidden inside the present.

But discovery is morally ambivalent. An entrepreneur can discover a way to provide clean water more cheaply, or to addict children to a platform. She can discover an underserved community and build dignified services, or discover a vulnerable population and monetise desperation. The market may reward both, depending on the structure it operates within. Discovery is therefore not the same as service. The entrepreneur as discoverer must become the entrepreneur as steward, asking not only what opportunity presents itself, but what good it serves, what dependencies it creates, and whether it increases human agency or captures it. This matters especially for startup culture. The language of innovation often praises discovery without enough moral interpretation. Founders are rewarded for finding scalable opportunities, but not all scalable opportunities are good. Some scale because they serve genuine need; others scale because they exploit psychological vulnerability, regulatory gaps, labour precarity, or network lock-in. A moral economy does not ask entrepreneurs to stop discovering. It asks them to discern what they have discovered.

\section*{Profit, Loss, and Moral Discipline}

Austrian economics emphasises profit and loss as mechanisms of learning.

Profit signals that resources may have been combined in a way that creates value for others. Loss signals that resources may be better used elsewhere. Markets discipline error. Entrepreneurs test hypotheses, consumers respond, prices move; some ventures survive, others fail, and knowledge emerges through feedback. A moral economy that ignores this will waste resources and harm the people it hopes to serve. An institution may have noble intentions but produce goods no one needs, services no one uses, or systems that cannot sustain themselves. The Monti di Pietà already taught us this: a charitable fund that depletes itself cannot serve the poor. Many bad ideas have survived far too long because they learned to speak morally while refusing feedback. Virtue is not improved by ignoring reality.

But profit and loss are not morally complete signals. Profit may result from genuine service, but also from extraction. Loss may indicate waste, but also mission fidelity in a context where the market fails to reward real goods. A company protecting workers, stewarding land, or serving poor communities may be less profitable than one that externalises costs. A platform that manipulates users may be more profitable than one that respects attention. Profit is feedback, not absolution. Loss is feedback, not always condemnation. A repayment formula tied to free cash flow must ask whether that cash flow is morally healthy: whether it comes from genuine service or from underpaid workers, locked-in customers, and neglected maintenance. The question is not simply whether the market rewarded this, but what the market rewarded.

\section*{The Beauty and Limits of Spontaneous Order}

There is genuine beauty in spontaneous order.

A market square, a language, an open-source ecosystem, a craft tradition, a scientific community: each can emerge from countless actions that no single person designed. Order appears without command. People adjust to one another, patterns form, knowledge circulates, innovation occurs at the edges, and participants discover uses and possibilities that central designers would never have imagined. This beauty should not be dismissed by moral economy. There is humility in allowing order to emerge, and Hayek's knowledge problem gives an economic reason for a principle that Christian social thought also values: human action is situated, local, and plural, and higher levels of authority should not absorb the proper agency of lower ones.

But beauty is not innocence. Spontaneous orders can exclude. They can reproduce hidden hierarchies, be captured by better-resourced actors, neglect maintenance, fail to protect the vulnerable, and be enclosed by those who understand the system well enough to monetise it. Harmful orders can also emerge without central design. Segregation can be spontaneous, and so can financial bubbles, extractive platform dependency, and addiction markets. The fact that an order is not centrally planned does not make it morally good.

The limits of spontaneous order become especially visible wherever the background conditions of just coordination are absent. Markets require law and trust, property rules and contract enforcement, limits on coercion, and access to information. When land was acquired through dispossession, market exchange in land begins from injustice. When workers lack real alternatives, labour markets reflect necessity as much as freedom. When ecological harms are externalised, prices mislead. When capital markets require exit, startup choices are shaped before founders appear to choose. A market is not just a process. It is an institutionally constituted process, and the rules of property and contract, of liability and money and corporate governance, and of public investment shape what market actors can do. Appeals to ``the market'' often conceal prior design. The question is not whether markets should be shaped by rules, because they already are. The question is whether they are shaped to serve life or to enable extraction.

\section*{The Problem of Power in Austrian Thought}

The Austrian tradition is powerful on knowledge but often weaker on power.

It understands that central planners lack knowledge, that markets can coordinate dispersed information, that entrepreneurs discover opportunities, and that prices transmit signals. But it can be less attentive to the ways power shapes which signals exist, whose knowledge counts, who can act on opportunity, and who bears the consequences of failure. A worker may know that a process is unsafe but lack authority to change it. A founder may know what the company needs but find that investors control the board. Dispersed knowledge does not automatically become effective action, and the gap between knowing and acting is often filled by whoever has more power.

Markets transmit knowledge through willingness and ability to pay, but those with more purchasing power speak louder. Those with urgent needs but little money may be nearly silent. Future generations do not bid. Ecosystems do not bid. Open-source maintainers may create enormous value without capturing any of it through price. The Austrian and liberal emphasis on exit as market discipline (if a platform mistreats users they can leave, if an employer mistreats workers they can move) also depends on conditions that may not hold. Network effects and switching costs, debt, geography, professional specialisation, medical need, housing scarcity, monopoly, and data lock-in can make exit formally possible but practically unavailable. When exit is constrained, the market discipline of competition weakens, and dependency becomes extractable.

A moral economy must therefore join Hayek's knowledge problem with a power problem: who knows, who can act on what they know, who controls the rules, who can leave and who cannot, who bears loss and who captures gain. Without these questions, spontaneous order can become a language that protects existing power while appearing to celebrate freedom.

\section*{Discovery, the Commons, and Moral Discernment}

Austrian economics often focuses on private property and price signals. But some of the most important forms of discovery happen in commons.

Scientific communities discover through shared inquiry, open-source developers through collaborative experimentation, and cultural traditions through imitation and adaptation. These are not always markets, yet they are decentralised knowledge systems, demonstrating that Hayek's insight about dispersed knowledge is broader than market price. Prices are one way of coordinating knowledge. Reputation and reciprocity, vocation, professional norms, gift exchange, peer review, and commons governance can also coordinate knowledge, and sometimes do so more effectively for the goods in question.

Open-source software is the clearest contemporary example. Much of the digital world depends on code created, maintained, reviewed, and improved through non-market or mixed-motive collaboration. The value created is immense, but not always captured by those who create it. Companies build proprietary services on top of open code; cloud providers monetise shared infrastructure; AI companies train on public knowledge; investors seek scalable returns from shared foundations. Sometimes this supports the commons. Often it extracts from it. A moral economy informed by Austrian insight should recognise the generativity of decentralised discovery without assuming that private capture is the natural reward for it. The commons needs stewardship precisely because it creates value that markets, left to themselves, may not sustain.

Entrepreneurial discovery also needs moral discernment, because opportunities are not morally equal. The same market that rewards innovation in clean water technology also rewards innovation in addictive interface design. The same venture capital structures that support genuinely useful infrastructure also support platforms that monetise behavioural manipulation. The entrepreneur must ask not only what opportunity she sees, but what good it serves, whether the discovery increases human agency or captures it, and whether it solves a problem or creates a profitable one. Entrepreneurship can be a genuine vocation when it serves real need through responsible risk, creativity, and stewardship. It becomes chrematistic when discovery is directed toward accumulation without regard for the good being created or destroyed. The goal is not fewer entrepreneurs. It is entrepreneurs formed by virtues and operating within structures that support those virtues rather than punish them.

\section*{Why Markets Still Require Moral Culture}

Austrian economics is right that knowledge is dispersed, that central planning is dangerous, that prices carry information, that entrepreneurs discover, and that spontaneous order can be beautiful. A moral economy that ignores incentives, scarcity, and unintended consequences will fail. But markets still require moral culture to function well.

Markets require trust: contracts cannot specify everything, and if people lie whenever it is profitable, exchange becomes fragile. They require restraint, because not every vulnerability should be monetised and not every legal opportunity should be taken. They require truthful valuation, since prices need background conditions that make them meaningful, including competition, transparency, and internalised costs. They require limits on power, because without antitrust, labour protections, data rights, and governance constraints, market actors can become powerful enough to deform the market itself. And they require a vision of the good. Markets can coordinate means, but they cannot tell us which ends are worthy; they can reveal what people are willing to pay for, but not what they ought to love.

This is the core argument against market idolatry. Markets are good servants and dangerous masters. Hayek himself understood this in part: he valued evolved tradition and moral culture precisely because he knew markets could not sustain themselves on calculation alone. People must be willing not to exploit every loophole, abandon every obligation, or liquidate every inheritance. They must keep promises, practise truthfulness, and retain some concern for reputation in the deeper sense, for being worthy of trust, not merely for appearing profitable. When moral culture collapses, markets become more dependent on regulation. When regulation is captured, markets become more extractive. When extraction becomes normal, trust collapses further. A free market without moral culture does not stay free. Moral culture is not decoration. It is market infrastructure.

\section*{Austrian Economics and Non-Extractive Capital}

What, then, can non-extractive capital learn from Austrian economics, and what must Austrian economics learn in return?

From Austrian economics, non-extractive capital can learn humility. A fund cannot know in advance which enterprises will flourish, which founders will discover real opportunities, or which technologies will matter. It must allow experimentation, and governance should protect mission without smothering entrepreneurial discovery. It can learn the importance of local knowledge: founders often know their customers, technologies, and communities better than investors do, and capital should support their judgment rather than dominate it. It can learn that prices and cash flows contain real information. A company's ability to generate durable surplus matters, and revenue, margin, and repayment capacity can reveal whether the enterprise is actually serving genuine need. And it can learn that loss matters. Not every mission-driven enterprise should continue forever, and stewardship includes recognising failure and reallocating resources rather than sentimentally depleting the fund.

But Austrian economics must also learn from the moral economy tradition. It must reckon with the fact that capital structure shapes discovery. If the only available capital demands exit, the market will discover exit-shaped companies. If finance rewards dependency, entrepreneurs will discover dependency. If investors seek uncapped upside, firms will be pushed toward monopoly and capture. The market process is not independent of the capital architecture that funds it. Austrian economics must also reckon with the ways prices can reflect power rather than service, voluntary exchange can conceal structural necessity, and commons can generate value without being reducible to private property. Market freedom is not the absence of structure. It is the presence of structures that protect persons and communities from domination. A non-extractive market economy would preserve entrepreneurial discovery while changing the terms on which capital participates: funding discovery without owning the future without limit, allowing founders to build without forcing exit, and rewarding useful service while capping disproportionate claim. This is not anti-market. It is market retrieval: recovering the discovery function of markets from the extractive pressures that increasingly distort it.

\section*{The Transition to the Corporation}

Austrian economics helps us see the market as a process of discovery. But modern economic life is not made only of markets. It is made of firms.

The corporation is not simply a participant in the market. It is a legal and governance structure that shapes how capital, labour, ownership, risk, and purpose are organised. It can enable productive enterprise at scale, coordinate complex work, accumulate knowledge, steward assets, and serve customers across time. But it can also become a sovereign institution. When corporate law, shareholder primacy, limited liability, capital markets, and managerial incentives all align around financial return, the firm can become less a community of production and more a vehicle for capital extraction. A platform corporation may coordinate millions of market participants while itself functioning as a private regulator. A financialised corporation may reorganise productive life around shareholder return in ways that no market participant can resist.

The Austrian tradition rightly warns us that centralised power is dangerous. The next chapter will show that centralised power is not only a danger of the state.

\chapter*{Chapter 12: The Corporation Becomes Sovereign}
\addcontentsline{toc}{chapter}{Chapter 12: The Corporation Becomes Sovereign}

The market is not the only institution of modern capitalism. In fact, much of modern economic life does not happen in markets at all.

It happens inside firms, through hierarchy and management, through employment contracts and ownership rights, through boards and financing agreements, through intellectual property, supply chains, platforms, and governance structures. A market and a corporation are not the same kind of thing. A market coordinates many actors through exchange. A corporation coordinates people and assets, capital and knowledge, contracts and strategy under a legal form. It can act across time, own property, sue and be sued, borrow, employ, acquire, lobby, and survive the death of its founders.

The corporation is one of the most powerful institutional inventions in history. It can gather resources for productive purposes that no household or small partnership could undertake alone, building infrastructure and coordinating research, employing thousands and producing complex goods, and preserving institutional knowledge across generations. But it can also become a vehicle through which capital escapes ordinary human limits. It can be a community of production ordered toward a useful purpose, or it can become an asset owned by capital, governed for capital, and disciplined by capital.

The question is not whether corporations are good or bad in themselves. The question is what they are ordered toward, who controls them, whose claims govern them, and whether their legal architecture makes stewardship durable or fragile. In the corporation, moral economy becomes legal architecture.

\section*{The Firm as Moral Architecture}

A firm is not merely a bundle of contracts.

People give years of their lives to it. Knowledge accumulates inside it. Relationships form around it. Customers depend on it. Workers develop skills, identities, habits, and expectations through it. Its decisions can affect land and housing, public infrastructure, political life, and technological dependency. A firm has a purpose, whether acknowledged or hidden.

This is why corporate governance is morally serious. Governance determines who has authority to decide what the firm is for. It determines who can redirect surplus and force sale, who can change strategy or replace management, issue shares, take on debt, sell assets, alter mission, merge, liquidate, or distribute profits. It determines who speaks when purpose and profit conflict. Shareholder primacy is not merely a financial theory. It is a doctrine of institutional purpose, a teleology. It tells managers and boards, investors, analysts, lawyers, and employees what the company is finally for. It says that other goods matter insofar as they contribute to shareholder value. Workers may matter, customers may matter, communities may matter. But if shareholder value is the governing end, all other goods remain conditional. The corporation becomes sovereign when this teleology is hidden under the language of neutrality, efficiency, fiduciary duty, or market discipline. It is not that the firm has no purpose. It is that its purpose has been smuggled in wearing a suit and calling itself procedure.

\section*{The Promise and Danger of the Corporate Form}

Before criticising the corporation, we should understand why it became so powerful.

The corporate form solves real problems. It allows enterprise to outlive individual owners, pools capital from many investors, enables complex projects requiring large upfront investment, and allows risk to be organised, delegated, and limited. It creates a recognisable legal person that can hold assets, enter contracts, borrow, and maintain continuity across generations. Without corporate forms, many essential activities would be difficult or impossible: railways and hospitals, energy systems, telecommunications networks, pharmaceutical research, large-scale manufacturing, and complex software infrastructure all require coordination beyond the scale of ordinary partnership. One can admire the dignity of the small workshop and still recognise that it is not going to build a national electricity grid by Thursday.

The corporation also separates the enterprise from the personal lives of its participants. A founder's death need not destroy the company; investors can support productive work without managing operations; workers can join a stable institution without personally owning its assets. This continuity can serve life. A good corporation can preserve knowledge, provide stable work, invest for the long term, and serve customers over decades. Some companies become institutions of trust, holding craft, memory, and productive capacity that would be lost if every enterprise dissolved with its founders.

The moral problem is not corporate continuity. The moral problem is continuity without accountability to purpose. A corporation can outlive people, but what exactly lives on? A mission, a community of production, or merely a legal shell and a claim on cash flows? The corporate form gives enterprise a kind of institutional body. Moral economy asks what soul animates it.

\section*{Limited Liability and Moral Distance}

Limited liability means that shareholders can invest in a company without being personally liable for all its debts and harms beyond the amount they invested. This encourages capital formation, enables risk-taking, and allows large-scale pooling of resources. Without it, many useful enterprises would never be financed. Limited liability can serve life.

But it also creates moral distance. It allows investors to receive upside while limiting downside, separates control, benefit, and responsibility, and can encourage risk-taking whose costs fall on workers, creditors, communities, ecosystems, or the public. A shareholder may benefit from cost-cutting that makes a company more fragile, while employees bear job loss when it fails. A private equity sponsor may benefit from debt-financed distributions while creditors and workers bear the collapse risk. A platform may benefit from independent contractors while avoiding employer obligations. In each case, the architecture allows upside and responsibility to separate, and that separation is one of extraction's favourite conditions. The answer is not to abolish limited liability but to embed it within stewardship: paired with duties and transparency, with worker protections and environmental responsibility, with governance checks, and with limits on distributions that weaken the enterprise. A legal shield can be useful. It should not become moral invisibility.

\section*{From Chartered Purpose to Shareholder Primacy}

Historically, many corporations were created by charter for specific purposes: to build infrastructure, to conduct trade, to administer a public function, or to manage a university. Corporate privilege was not always treated as a general private right. It was granted under conditions, with some claim that the entity served a public interest.

This older world had many injustices. Chartered companies could be monopolistic and colonial, violent, and corrupt. They did not always serve the public so much as arrive bearing a charter and a cannon. But the idea of chartered purpose recognised something important: the corporation is a legal creation. It receives privileges from the political community and therefore should answer to some purpose beyond private gain.

General incorporation changed this. As incorporation became easier and more standardised, the corporation became a normal vehicle for private enterprise. This democratised corporate formation in some ways and reduced arbitrary state favouritism. But as incorporation became ordinary, the public question of purpose weakened. The corporation increasingly appeared as a neutral legal container whose purpose could be whatever its owners wanted within the boundaries of law. When a legal form created by public authority is treated as a private instrument for capital accumulation, the question ``what does society receive in exchange for granting this form?'' becomes harder to ask. The corporation gains legal personality and asset ownership, limited liability and perpetual life, entity shielding and access to courts, but its obligation to the common good becomes diffuse. Moral economy must recover the public nature of corporate privilege. A corporation is not a natural person exercising inherent freedom. It is an institution authorised by law. Its privileges are designed. Therefore its responsibilities can also be designed.

\section*{Colonial Corporations and the Fusion of Commerce and Sovereignty}

The corporation's sovereignty is not merely metaphorical. Some of the most powerful early corporations (the Dutch and British East India Companies, the Virginia Company, the Hudson's Bay Company) exercised political and military functions. They governed territories, negotiated treaties, raised armies, extracted resources, and controlled trade. Commerce and sovereignty were fused.

This history matters because it reveals a deep temptation in the corporate form: the ability of private capital to acquire public power. Colonial corporations were not simply businesses operating in markets. They were instruments of empire that combined enclosure and monopoly, violence and political privilege, racial domination, and financial return. They operated in the blurred space between state and market, between public authority and private profit. Responsibility disappeared into distance: investors in one place benefited from domination elsewhere.

The modern corporation is not identical with colonial chartered companies, and it would be careless to treat every firm as a continuation of empire. But the history warns that corporate power can become political power, and political power can be used to secure corporate gain. Today, large corporations shape law and infrastructure, speech and public knowledge, labour markets and tax systems, urban development, digital dependency, and geopolitical strategy. When a corporation controls infrastructure on which public life depends, it exercises quasi-sovereign power. The colonial corporation is a warning from history: commerce without moral and political limits can become sovereignty for profit.

\section*{The Firm as Asset}

When shareholder primacy dominates, the firm is increasingly seen as an asset: a bundle of cash flows and claimable property that can be bought, sold, leveraged, restructured, optimised, split, merged, or liquidated. This perspective can reveal real inefficiencies and sometimes preserve enterprises that might otherwise fail. But it becomes dangerous when it displaces the institutional view.

A firm is not just an asset. It is a living organisation of people and purposes, of relationships, knowledge, obligations, and trust. Treating it primarily as an asset destroys forms of value that do not appear on the balance sheet: worker loyalty and craft knowledge, customer trust and community rootedness, mission integrity, supplier resilience, and long-term capacity. Financial engineering often profits by making these hidden goods visible only after they are gone. A company can cut experienced workers and show improved margins, until quality declines; it can reduce maintenance and show higher cash flow, until systems fail; it can load debt and distribute cash, until downturns become catastrophic. The asset view sees liquidity. Stewardship sees inheritance.

This is why mission-locked governance matters. It protects the firm from being reduced to a financial object, saying that some aspects of the enterprise are not for sale simply because a price can be offered. Capital that serves life can invest in firms. Capital that feeds on life turns firms into prey.

\section*{Financial Engineering and Productive Enterprise}

Financial engineering is not always bad. Capital structure matters, and debt and equity can be arranged in ways that make enterprise more resilient, smooth cash flow, fund equipment, enable transition, and structure ownership in alignment with mission. The question is always what financial engineering serves.

When it serves productive enterprise, it helps the company do its work better. When it extracts, it treats the enterprise as material for return. The difference is visible in how financial tools affect the productive core. Does debt fund productive investment, or finance distributions to owners? Does restructuring strengthen the company, or prepare it for sale? Does an acquisition preserve mission, or harvest assets?

This becomes especially important in private equity and highly leveraged acquisition structures. A firm may be purchased using significant debt placed on the acquired company, redirecting cash flows that could have supported workers, innovation, resilience, or long-term investment toward debt service, fees, and distributions. If the company succeeds, owners may receive high returns. If it fails, workers, creditors, suppliers, and communities bear much of the damage. This reveals the moral significance of leverage: it can amplify productive investment, but it can also amplify extraction, allowing owners to control assets with less equity while the burden of repayment falls on the enterprise rather than on the capital provider. A moral economy does not reject financial design. It demands that financial design be accountable to productive purpose.

\section*{The Corporation and the Worker}

The corporation shapes the moral meaning of work.

In a small workshop, a cooperative, or a craft tradition, the relationship between work, product, community, and purpose may be visible. In a large corporation, that relationship can become abstract. Workers may perform specialised tasks within systems they do not control, for purposes they did not choose, under metrics they did not design. Large organisations can offer stable employment, training, benefits, and the chance to contribute to meaningful work at scale. But the corporate form can also distance workers from the purposes of their labour. An engineer may optimise engagement without seeing addiction. A content moderator may absorb psychic harm so a platform can remain usable. The question is not only whether the worker is paid. It is whether the worker's humanity is honoured in the organisation of work.

A corporation governed primarily by financial return will tend to see labour as cost. This shapes imagination: workers become something to optimise and reduce, to monitor, outsource, casualise, or replace when possible. A moral economy sees workers as among the persons through whom the firm exists, people whose knowledge, judgment, care, and dignity are part of the enterprise's moral reality. If workers are central to the firm's life, their voice matters. If they bear risk, they should share in security and upside. Worker cooperatives and codetermination, employee ownership and profit sharing, steward ownership, and strong labour standards are different ways of addressing the same moral reality: labour is not merely a cost to capital. A corporation becomes extractive when it treats workers as instruments of financial claim rather than participants in productive purpose.

\section*{The Corporation, the Community, and the Commons}

Corporations do not exist in empty space. They rely on communities: on roads and schools, on legal systems and public safety, on educated workers and social trust, on ecological systems, housing, and culture. Even the most ``private'' enterprise depends on a world it did not create. A company that benefits from a community cannot treat that community merely as a location, a labour pool, a market, or a tax jurisdiction. It owes something to the conditions that make its existence possible.

The extractive corporation denies this dependence. It locates where subsidies are high and leaves when incentives expire, uses local labour while suppressing wages, benefits from public infrastructure while avoiding taxes, and draws on community trust while externalising risk. When capital can leave but workers, communities, and ecosystems cannot, the corporation exercises power beyond ordinary market exchange. Power creates obligation, and a corporation that shapes community life must be accountable to more than shareholder return.

Corporations can also help build commons: contributing to open-source software, funding research, maintaining public infrastructure, publishing knowledge. But corporations often face incentives to capture commons once value becomes visible. Open protocols become proprietary ecosystems, public research becomes patented monopoly, user networks become platform property, and maintainer labour becomes corporate dependency. The corporation is particularly effective at enclosure because it has legal personality and capital, lawyers and patents, contracts, and strategic continuity. A loose community may create value, but a corporation may be better positioned to capture it. A moral economy must therefore ask how corporations can participate in commons without enclosing them: through open licences, data trusts, interoperability requirements, contribution obligations, and community governance.

\section*{Platform Corporations as Private Governors}

The platform corporation is one of the clearest examples of corporate sovereignty in our time. A platform does not merely sell a product. It creates and governs a field in which others interact, setting rules and ranking visibility, controlling access, collecting data, resolving disputes, changing terms, designing incentives, and determining who can reach whom. Marketplaces govern sellers, app stores govern developers, social networks govern speech, cloud providers govern infrastructure access, and payment processors govern transaction possibility. These corporations do not simply participate in markets. They make markets.

This creates a new form of sovereignty. The platform's rules may matter more to participants than public law does in daily practice. A small business may survive or fail based on search ranking, marketplace fees, account suspension, or algorithmic visibility. When platform power becomes infrastructural, exit becomes difficult, and the corporation can extract from dependency by raising fees and privileging its own services, by altering algorithms and demanding data, by imposing terms or restricting interoperability. Participants may formally consent, but the field of choice has been shaped by the platform's control over access.

Platform corporations reveal why the old market-state binary is insufficient. They are private entities exercising rule-like power over economic and social life, neither ordinary market participants nor public governments. This hybridity allows power to escape older forms of accountability. Should platform governance include affected stakeholders? Should essential digital infrastructure be treated as commons or public utility? Should algorithms that allocate livelihood be accountable to democratic oversight? These are not merely technology policy questions. They are questions of moral economy. The platform corporation shows what happens when capital, code, and governance merge.

\section*{The Myth of Neutral Governance}

Corporate governance often presents itself as technical. Boards oversee management, shareholders elect directors, managers pursue strategy, financial reporting ensures accountability, and markets discipline poor performance. This language makes governance appear neutral, procedural, and rational. It has all the soothing authority of a flowchart.

But governance is never neutral. It allocates power and defines purpose, determines whose interests count, decides which risks matter, gives some parties voice and others silence, and makes some forms of value visible and others invisible. A board that represents only capital is not neutral; a fiduciary culture centred on shareholder value is not neutral; a reporting system that counts financial return but not mission drift is not neutral. These are moral choices. The myth of neutral governance helps extraction because it hides the fact that the corporation has been ordered toward specific ends. When those ends produce harm, defenders can say: we followed governance, we obeyed fiduciary duty, we responded to market signals. But the rules themselves are the moral architecture.

This is why alternative governance forms matter; they are not decorative. A cooperative, a foundation-owned company, a steward-owned firm, a benefit corporation, or a mission-locked enterprise each modifies the answer to who governs, for whose benefit, and under what limits. Legal form alone is never enough. Forms can be captured and missions can drift. But form matters because governance determines which moral commitments can survive pressure. If mission depends only on founder personality, it is fragile. If mission is embedded in governance, it becomes more durable. Moral economy must treat governance design as one of its central disciplines.

\section*{When the Corporation Becomes Sovereign}

The corporation becomes sovereign when it no longer functions as a tool for human purposes but as an authority to which human purposes must conform.

This happens when communities reshape themselves to attract and retain corporate investment at the cost of labour standards, housing stability, or ecological health; when governments depend on corporations for infrastructure they no longer know how to govern; when workers organise their lives around corporate scheduling, metrics, surveillance, and precarity; when founders must shape companies for investor exit rather than mission; when users must accept platform terms because social, professional, or civic life depends on access; when law bends around the expectations of capital markets. This is not sovereignty in the formal sense of the state, but it is functional sovereignty: the power to set conditions under which others must live, work, speak, trade, build, or access necessary goods.

Every tradition in this book would recognise the danger. Aristotle's fear of means becoming ends, Christianity's condemnation of avarice and domination, the Scholastics' commutative justice, the Franciscans' concern for capital feeding on dependency, Smith's critique of rent-seeking, the Protestant warning about stewardship lost: all of these converge in the sovereign corporation. The corporation becomes sovereign when no one can effectively ask what it is for.

\section*{Reclaiming the Corporation for Stewardship}

The answer is not to abolish the corporation. The corporate form is too useful, too embedded, and too capable of serving real goods. The question is how corporations should be governed.

A corporation can be reclaimed for stewardship when its legal and financial architecture is ordered toward purpose rather than extraction. Purpose must be defined beyond financial return and protected when return conflicts with it. Capital must be limited: investors may deserve return, but not unlimited authority over mission, workers, community, and future. Control must be entrusted to stewards bound to the company's purpose rather than to tradable financial interests. Exit must be disciplined, because if a company can be sold to the highest bidder regardless of mission, stewardship is fragile. And accounting must expand: the firm must be able to see mission drift and worker depletion, ecological harm and dependency creation and commons extraction, not only revenue and profit.

These shifts are practical, not utopian. They already appear in various forms: cooperatives and mutuals, industrial foundations, steward-owned firms, community interest companies, benefit corporations, mission locks, perpetual purpose trusts, employee ownership, and foundation-controlled enterprises. None is perfect, and all can drift. But each challenges the assumption that the corporation must be governed as capital's instrument. The firm can be a site of stewardship. The architecture exists. The question is whether it is chosen, protected, and replicated at scale.

\section*{The Transition to Financialised Capitalism}

The corporation becomes especially dangerous when joined to financialised capitalism. A corporation ordered toward productive purpose can serve life. A corporation ordered toward shareholder value can become extractive. But when the corporation itself becomes a financial asset in markets dominated by liquidity, leverage, short-term metrics, and portfolio logic, extraction intensifies.

The firm is no longer merely expected to produce profit. It is expected to produce financial performance according to the demands of capital markets. Its future cash flows become tradable expectations, its assets become collateral, its workers become adjustable costs, and its mission becomes relevant only insofar as it supports valuation. Governance becomes accountable to investors who may have no continuing relationship with the enterprise beyond return. In that world, extraction is not an occasional abuse. It becomes a system. The pressure for return moves through funds and boards, through executives and debt markets, through analysts and ratings agencies, through compensation structures and investor expectations. Even well-intentioned managers may find themselves governing under an architecture that punishes stewardship.

The corporation gave capital a body. Financialisation gives capital a nervous system. Together, they can make extraction appear not only normal, but necessary. The next chapter turns to that crisis.

\chapter*{Chapter 13: The Crisis of Financialised Capitalism}
\addcontentsline{toc}{chapter}{Chapter 13: The Crisis of Financialised Capitalism}

The corporation gave capital a body. Financialisation gave it a nervous system.

Through the modern corporation, capital could enter enterprise, hold rights, govern assets, claim surplus, and persist across generations. Through financialisation, those claims became liquid and tradable. They could be leveraged, modelled, securitised, and increasingly detached from the concrete lives of the institutions they governed. Financialisation is not merely the growth of finance as a sector. It is the spread of financial logic into the whole economy, until a company is valued as a portfolio of cash flows, a home as an asset class, infrastructure as investable yield. Education is reconceived as human capital investment, healthcare as billing architecture, human attention as inventory. Finance no longer simply supports production. It increasingly defines what production is for.

A financial system can serve life. It moves resources across time, pools risk, funds enterprise, and protects households. The problem is not finance as such. The problem is finance becoming sovereign. When financial logic governs the firm, the household, the city, and the state, everything is asked to justify itself as a stream of future returns. The question ``what is this for?'' is replaced by ``how can it be capitalised?'' The grammar of stewardship gives way to the grammar of asset management. The category of enough disappears. The future is not inherited; it is discounted. Financialised capitalism is what happens when capital not only seeks return, but reorganises the world so that return becomes the measure of reality.

\section*{From Productive Finance to Financialised Control}

Finance becomes morally dangerous when it ceases to be a servant of productive life and becomes the ruling logic by which productive life is judged.

Productive finance asks how capital can support real work across time. It funds the things real work requires (tools, equipment, housing, infrastructure) and bridges the gap between present need and future surplus, sharing risk and making patience possible. Financialised control asks a different question. How can claims on future value be created, priced, leveraged, and extracted? This shift changes everything. A company is no longer primarily a community of production but an asset whose value lies in expected cash flows. A home is no longer primarily shelter but a financial instrument valued for appreciation and rental yield. A hospital is no longer primarily a place of healing but a billing machine with assets, liabilities, and revenue optimisation.

Financialisation does not necessarily deny the original purpose of these things. It subordinates them. The company may still produce, the home may still shelter, the hospital may still heal, but these goods increasingly survive only insofar as they can be reconciled with financial claim. Finance exists to serve life by moving resources across time. In financialised capitalism, life is reorganised to serve finance by generating claims across time. That inversion is the modern form of chrematistics.

\section*{Private Equity and the Extraction of Enterprise}

Private equity is one of the clearest sites where financialised capitalism must be morally examined. At its best, it can provide capital, discipline, and succession solutions to companies that need restructuring or strategic support. Some firms are genuinely poorly managed. Some family businesses lack successors. Not every private equity transaction is extractive. But the structure is morally dangerous.

Private equity often acquires companies using significant debt, then seeks to increase the value of the investment over a limited time horizon before sale. The acquired company may carry the debt used to purchase it. Cash flows may be directed toward debt service, fees, dividends, and eventual exit. Operational changes may be evaluated primarily through their effect on financial performance. This creates a powerful temptation to treat the company as a financial object before treating it as a productive community. Workers become adjustable costs. Long-term maintenance becomes a deferrable expense. Supplier relationships become margin pressure. Customer trust becomes monetisable pricing power. The company's future becomes collateral for the acquisition of itself, which is a sentence that should probably make everyone sit quietly for a moment.

The central moral question is not whether private equity can increase financial value. It often can. The question is how that value is produced. Through better products, wiser operations, long-term investment? Or through leverage, fees, layoffs, and asset sales? A company can become more valuable to investors while becoming less capable of serving the people who depend on it. The most morally revealing question is simple: what remains after capital exits? Is the enterprise stronger, more resilient, better governed? Or has value been removed while risk and fragility remain? Exit reveals stewardship. A steward leaves the entrusted good more capable of life. An extractor leaves with the surplus.

\section*{Debt as Domination}

Debt is one of the most ancient and morally charged economic relationships. It can serve life. It helps a household survive hardship. It allows a farmer to plant before harvest. It finances a home, funds a business, builds infrastructure. Lending can be an act of trust, solidarity, and productive participation. But debt can also dominate.

Debt creates an asymmetry across time. The borrower receives resources now and becomes obligated later. This can be just when the obligation is proportionate, transparent, repayable, and connected to real capacity. It becomes dangerous when repayment claims exceed the borrower's ability to pay without harm, when default transfers control, or when debt is used strategically to discipline behaviour. Before surplus can be shared, debt must often be serviced. Before mission can breathe, creditors may demand compliance. Debt can become governance by another name.

A household under heavy debt loses real freedom even if no one is visibly coercing it. Choices narrow, work becomes compulsory, risk becomes terrifying, and the future is pre-claimed. A company under heavy debt is forced to prioritise cash extraction over long-term investment, its future pledged to creditors. A state under debt pressure may be forced into austerity or privatisation that weakens public life. Debt turns time into obligation. When the obligation becomes disproportionate, debt becomes a form of rule. This is why usury remains a living moral category. The older prohibition against profiting from need was not primitive financial ignorance. It was an intuition that debt can convert vulnerability into claim. Debt should be servant to life across time. When it becomes master, the future is colonised.

\section*{Housing and Speculation}

Housing reveals financialisation with particular intimacy because housing is simultaneously asset and home. A house can store savings, secure a loan, provide inheritance, and form part of a household's financial stability. There is nothing inherently wrong with housing having economic value. But housing is not only an asset. It is shelter. It is the place where people sleep, where families form, where memory accumulates, where belonging takes root. When housing is governed primarily as an asset class, the logic of shelter is subordinated to the logic of return.

Investors seek yield. Land values rise. Homes are purchased as stores of wealth. Rents rise to whatever the market will tolerate. Short-term rentals compete with residents. Luxury development coexists with homelessness. Young households are priced out of the places where they work. A housing market can appear healthy from the standpoint of asset owners while becoming morally disordered from the standpoint of human flourishing. Rising prices celebrate wealth creation for those who own property; for those without, they represent exclusion, dependency, and eventually displacement. Housing teaches us that financialisation is not abstract. It enters the place where people sleep.

Speculation intensifies the problem because it profits from expected future scarcity or appreciation without improving the home, serving the tenant, or strengthening the community. This is rentier logic: wealth from control of access. A moral economy must distinguish the landlord who maintains housing well, charges fair rent, and treats tenants with dignity from the investor who profits chiefly from scarcity, displacement, and leverage. The home becomes chrematistic when its exchange value overwhelms its use. Recovering housing as shelter before asset requires institutional forms that remove at least some housing from speculative pressure. Community land trusts, cooperatives, mission-locked ownership, and financing structures that prioritise long-term affordability over maximum appreciation all point in this direction.

\section*{Growth Without Shared Flourishing}

Modern capitalism often defends itself through growth, and the defence cannot simply be dismissed. Growth has, in many contexts, improved material conditions dramatically. It has helped reduce poverty, expanded access to goods, funded healthcare, and supported education. A moral economy should not romanticise stagnation. Scarcity can crush human possibility, and economies need productive capacity. Empty shelves are not a moral achievement.

But growth can also conceal extraction. An economy can grow by expanding debt and inflating asset values. It can grow by monetising what was once freely given, by making housing more expensive, by producing addiction, by depleting ecosystems, by externalising costs onto the future. The fact of growth tells us nothing until we know the form of life. A cancer grows; so does a forest. Financialised capitalism often produces growth without shared flourishing because it excels at converting life into financial value. Attention becomes advertising revenue. Illness becomes billing. Loneliness becomes engagement. Each conversion may increase measured economic activity. Not all increase human flourishing. A moral economy must therefore distinguish expansion from health, asking not merely whether the economy is growing but what is growing, for whom, and at what cost to the goods that actually sustain human life.

\section*{Why Growth Alone No Longer Saves Us}

There was a time when the promise of growth could answer many criticisms. Even if capitalism was unequal, growth would raise all boats. Even if markets were harsh, growth would create opportunity. Even if workers suffered now, growth would create better jobs later. This promise has weakened.

In many advanced economies, growth no longer reliably produces broad security, affordable housing, stable work, or ecological health. Productivity gains do not automatically translate into shared prosperity. Asset owners may benefit disproportionately while workers experience precarity despite economic expansion. Young people inherit higher costs, lower security, and ecological risk. The issue is not simply that growth is too low. The issue is that the architecture beneath the growth determines what growth does. Growth flowing through extractive structures deepens extraction. When companies are governed for shareholder value, growth enriches shareholders while workers remain insecure. When housing is financialised, growth raises land values while excluding residents. Growth cannot save an economy whose institutions are ordered toward extraction. More activity inside the wrong architecture may simply strengthen the wrong architecture.

The question is not growth or no growth. The question is whether growth is governed by stewardship. Capital that serves life may grow, but it grows as a tree grows: rooted, bounded by form, oriented toward fruit and continuity. Capital that feeds on life grows as fire grows: by consuming what surrounds it.

\section*{The Rentier Economy}

A rentier economy is one in which wealth increasingly comes from ownership and control of assets rather than productive contribution. Rentiers gain because they control something others need. It may be land or housing, intellectual property, infrastructure, data, a platform, or simply access to a network that has become indispensable. They do not necessarily produce more. They position themselves to charge for access.

The rentier asks: what must others pass through? Build the platform, own the data, hold the patent, become the intermediary; then charge for passage. Some of this may begin productively. A platform may solve a real coordination problem, a patent may reward real invention, land ownership may support real stewardship. But once the chokepoint is established, the logic can shift from service to extraction. The rentier economy is especially dangerous because it can appear peaceful and efficient. No one is visibly robbed. Contracts are signed, fees are paid, users consent, tenants rent, developers accept platform terms. But underneath the formal exchange lies a deeper dependency: access to necessary goods is controlled by those whose primary contribution may be ownership itself.

This stratifies society between those who own claims and those who must pay to access life. It rewards asset appreciation over work. It makes younger generations dependent on older asset holders, turns public investment into private gain, and increases the power of those who can wait and own over those who must work to live. Every tradition in this book has named this pattern. Aristotle feared wealth detached from use; the Church condemned usury and avarice; Smith warned against monopolists and rent-seekers. The rentier economy is their shared nightmare in modern form.

\section*{Speculation and the Detachment from Reality}

Speculation is not always wrong. A degree of it is inherent in economic life. A farmer speculates when she plants, an entrepreneur when he builds, an investor when she funds an uncertain project. Human beings act under uncertainty, and foresight can be a real contribution. Life-serving speculation bears risk in order to bring about or preserve a real good.

Extractive speculation profits from price movement without contributing proportionately to the good whose price moves. It may amplify volatility, inflate bubbles, hoard scarce assets, or draw capital away from productive use. It treats food, energy, land, or technology as chips in a game whose consequences are borne by others. Financialised capitalism encourages this detachment because claims can be traded far from underlying goods. Securitised mortgages and complex credit instruments create layers of claims so abstract that the underlying human reality disappears. The trader does not see the family losing a home. The investor does not see the worker in the acquired company. Abstraction is not always bad, but moral economy must continually reconnect claims to realities. What good underlies the claim? What service justifies the return? What risk is being borne, and by whom? Speculation becomes extraction when it profits from reality while refusing responsibility for reality.

\section*{The Financialisation of the Self}

Financialisation does not stop at institutions. It enters the person.

Human beings are increasingly encouraged to understand themselves as portfolios. The self becomes a bundle of assets and skills, of credentials and habits and future earning potential. Education becomes investment in human capital. Health becomes productivity optimisation. Friendship becomes networking. Rest becomes recovery for performance. Some of this language can be useful. People do need to think prudently about education, work, health, and money. But when the whole self is financialised, the person becomes both investor and asset: a tiny private equity fund with a nervous system.

This is spiritually exhausting. One must constantly optimise. Build the profile, increase value, manage risk, convert passion into income, convert suffering into content, convert relationships into opportunity, convert rest into future output, convert the soul into brand. The old Protestant danger returns in secular form: vocation without grace, discipline without Sabbath, work without rest, self-examination without mercy. Financialised capitalism does not merely extract from institutions. It trains people to extract from themselves. An economy ordered toward life must therefore protect what cannot be monetised. It must honour people who are not maximising themselves, and refuse to measure human worth by earning potential, visibility, or efficiency. Sabbath is resistance to self-financialisation. Craft is resistance to pure scalability. Care is resistance to productivity as worth.

\section*{The Fragility Hidden by Financial Performance}

Financialised capitalism often appears strong just before it reveals fragility. A company can report rising earnings while neglecting maintenance. A household can appear prosperous through debt. A bank can appear profitable while accumulating systemic risk. A housing market can appear healthy while becoming unaffordable. Financial metrics often reward the removal of redundancy, but life requires redundancy.

Households need savings. Companies need reserves. Hospitals need spare capacity. Supply chains need buffers. Democracies need public institutions that are not optimised to the edge of collapse. Financialisation often treats these reserves as inefficiency, seeking higher returns by reducing slack, accelerating turnover, increasing leverage, and extracting idle capacity. In stable times, this appears rational. In crisis, the hidden fragility becomes catastrophic. The moral economy tradition has a better word for some forms of redundancy: prudence. Reserves are not always inefficiency. Slowness is not always failure. Maintenance is not lack of innovation. Rest is not idleness. A financial system that cannot distinguish waste from resilience will eventually destroy resilience in the name of efficiency. Extraction often appears as performance because it consumes resilience before anyone accounts for it.

\section*{The Crisis of Trust}

Financialised capitalism corrodes trust. It teaches people that every relationship may conceal a business model. The platform wants attention. The lender wants fees. The landlord wants yield. The hospital wants billing codes. The app wants data. The investor wants exit. Not every actor is cynical; many people inside these systems are sincere, generous, and competent. But the structure trains suspicion. The customer suspects manipulation, the worker suspects extraction, the user suspects surveillance, and these suspicions are often correct.

Trust collapses when people believe, often correctly, that the other party's gain depends on their dependency, ignorance, or lack of alternatives. And trust matters deeply. Markets require it. So does democracy. So does the patient cooperation that produces commons, science, and craft. When trust collapses, every transaction becomes more defensive, more legalistic, more surveilled, more costly, and more brittle. Extraction is therefore self-undermining. It may generate returns in the short term, but it destroys the moral infrastructure on which productive cooperation depends. Once people expect extraction, they either withdraw, retaliate, comply cynically, or seek their own opportunities to extract. The system becomes a school for mistrust.

A moral economy must rebuild trust by making trustworthiness structural. Capped returns, mission locks, transparent accounting, worker voice, fair contracts, exit restrictions, legal duties that match stated purpose: these are not bureaucratic embellishments but the materials from which trust is rebuilt under pressure. Trust grows when people can see that the other party cannot easily exploit them even under pressure. Institutional design is therefore pastoral in the broad sense. It cares for the conditions under which people can act without constant fear of being used.

\section*{The Political Crisis of Financialisation}

Financialised capitalism produces political crisis. When people experience the economy as extractive, they lose faith in institutions. Work no longer leads to security. Housing becomes unattainable. Healthcare becomes frightening. Education requires debt. Public services decline. Wealth accumulates among asset holders. Political leaders speak of growth while ordinary life feels more fragile. This gap breeds anger, some of it directed toward legitimate critique of monopoly and institutional capture, some redirected toward scapegoats. Extractive systems often survive by allowing the wounded to fight one another rather than the structures that wound them.

Financialisation hollows out democracy because democratic institutions appear less able to govern capital. Governments fear investor flight. Cities compete for corporate headquarters. Public agencies depend on private vendors. Infrastructure requires private financing. Digital platforms govern speech beyond national borders. When capital becomes mobile and communities remain rooted, democracy weakens. People may still vote, but many of the conditions shaping their lives seem decided elsewhere: in boardrooms, in funds, on platforms, in global supply chains. This produces the sense that politics is both loud and powerless.

A moral economy must restore democratic and communal agency over the conditions of life. This does not mean centralised state control of everything. That would ignore the lessons of Hayek and the importance of plural institutions. It means capital must be re-embedded within legal, civic, local, and moral accountability. It must not be free to govern without being governed. Subsidiarity and solidarity become crucial. Financialisation is political because extraction is power.

\section*{Why Charity After Extraction Is Not Enough}

Financialised capitalism often tries to repair its moral damage through philanthropy, corporate social responsibility, impact investing, or ESG. Some of this is good. Charity matters. Philanthropy can fund real goods. Corporate responsibility can improve practices. ESG can make some risks visible. We should not dismiss every effort simply because it occurs within imperfect systems. The world is too damaged, and the good too urgently needed, for that kind of purity.

But charity after extraction is not moral economy. If wealth is accumulated through extractive structures and then partially redistributed through philanthropy, the original disorder remains. If a company harms workers, communities, or ecosystems and then funds social programmes, the giving may help but does not justify the harm. If a platform encloses a commons and then donates to digital inclusion, the deeper governance problem is not solved. The question is not only what capital does after it has gained, but how capital gains in the first place.

Older moral economy understood this. Almsgiving did not automatically purify usury. Generosity did not erase exploitation. The rich were called not only to give but to practise justice. Modern moral economy must recover that order. Justice precedes charity. Structure precedes compensation. Stewardship precedes philanthropy. This is not an argument against giving. It is an argument against using giving to avoid conversion. A fund should not extract and then donate. A company should not exploit and then sponsor. Capital that serves life must be structured to serve before surplus is distributed.

\section*{The Search for Post-Extractive Finance}

If financialised capitalism is a crisis of architecture, then the response must be architectural. We need financial forms that can fund useful enterprise without forcing it toward extraction. That means investors who can receive fair return without demanding infinite upside. It means companies that can access capital without surrendering mission. It means repayment structures based on durable surplus rather than exit value, governance that protects purpose beyond founder intention, accounting that distinguishes real surplus from hidden depletion, ownership forms that prevent mission assets from being liquidated, and infrastructure finance that preserves public accountability.

This is the search for post-extractive finance. It does not mean finance disappears. It means finance is converted back into service. The practical questions it must answer are demanding. How much return is enough? When has capital been fairly repaid? Who controls the enterprise after investment? How is stewardship preserved across generations? How can capital support growth without demanding exit? But these are also ancient questions in modern form. When is gain legitimate? When does lending become usury? When does ownership carry obligation? When does accumulation become unnatural? The search for post-extractive finance is not a novelty. It is a retrieval.

\section*{The End of Diagnosis}

This chapter concludes the diagnostic part of the book. The traditions traced in earlier chapters named the danger of wealth detached from life. The analysis of markets, corporations, and financialised capitalism showed how that danger operates in modern form. We now have a working definition of extraction, an understanding of how it colonises firms and finance, and a vocabulary for distinguishing it from productive capital.

Now the book turns toward recovery. Critique without construction becomes despair, performance, or purity. A moral economy must build. It must recover traditions capable of guiding institutional design and translate them into forms of ownership, finance, governance, and technology. The remaining chapters attempt exactly that.

\chapter*{Chapter 14: Subsidiarity, Solidarity, and the Common Good}
\addcontentsline{toc}{chapter}{Chapter 14: Subsidiarity, Solidarity, and the Common Good}

If Part IV diagnosed extraction and financialised capitalism, Part V begins the recovery.

The question now is not merely what is wrong, but how institutions, markets, and finance can be ordered toward human flourishing. This requires construction as well as critique: traditions capable of guiding institutional design, and principles that can be translated into forms of ownership, finance, governance, and technology. It is one thing to identify the disease. It is another to stop admiring the X-ray and begin treating the patient.

Catholic Social Teaching offers a rich resource for this work. Far from being purely devotional, it provides practical principles for governance, labour, and social organisation that have been refined over more than a century of engagement with the actual conditions of modern life. It is not a perfect tradition. It has had its own internal tensions, and it has sometimes failed to apply its principles as broadly as they demand. But it has preserved, against considerable pressure, something essential: the insistence that economic life must be judged by whether it serves human flourishing, and that institutions which fail this test must be reformed. This chapter focuses on three of its core principles. Subsidiarity, solidarity, and the common good are not abstract ideals but operational guides for institutional design, financial architecture, and corporate governance, tools for preventing extraction while enabling productive enterprise.

\section*{A Tradition of Critique: From Rerum Novarum to Laudato Si'}

Modern Catholic Social Teaching begins with Pope Leo XIII's Rerum Novarum, published in 1891 at the height of industrial capitalism's most brutal period. Leo XIII was responding to conditions the older moral vocabulary had not fully anticipated: mass urbanisation, the rise of factory labour, the destruction of older household economies, and the concentration of industrial capital on a scale that made earlier concerns about individual usurers seem almost quaint. A lone moneylender can do damage. An industrial order can reorganise a civilisation.

The document was not a rejection of markets or private property. Leo XIII explicitly affirmed both, against socialist proposals that he found dangerous and anthropologically false. But he insisted that ownership was never absolute. It was always a stewardship toward family, community, and the common good. Workers had rights that employers were obligated to respect, the state had duties to protect those who could not protect themselves, and wealth carried obligations that its legal possession could not dissolve. The Church was not indifferent to economic life. Economic life was a moral field, and the Church had a responsibility to say so.

Subsequent documents expanded and sharpened the critique. Quadragesimo Anno, in 1931, addressed the rise of financialised industrial capitalism, warning against the excessive concentration of economic power and the way financial interests could come to dominate productive enterprise. Centesimus Annus, written in 1991 at the end of the Cold War, reflected on market mechanisms, globalisation, and the enduring claims of social justice. Caritas in Veritate, in 2009, addressed the global economy after the financial crisis, insisting that charity and justice could not be separated, and that economic life required institutions ordered toward authentic human development. Laudato Si', in 2015, brought ecological stewardship explicitly into the framework, arguing that the same logic of extraction that impoverishes communities also degrades the natural world, and that the cry of the poor and the cry of the earth cannot be heard separately.

Across this tradition, running from Leo XIII through Francis, the Church has consistently argued that markets, corporations, and finance are instruments. They are powerful, useful, and necessary, but always in need of moral ordering. Legality is insufficient, since an institution may comply fully with the law while undermining human flourishing. Growth and efficiency are real goods, but only when subordinate to a moral order that protects life, community, dignity, and the natural world. The tradition names what modern economics often cannot: structural injustice, the dignity of work, ecological responsibility, and the duties of capital alongside the rights of ownership.

\section*{Subsidiarity: Governance at the Right Scale}

Subsidiarity is one of the most practically important principles in Catholic Social Teaching, and one of the most widely misunderstood. It is sometimes read simply as a preference for small government or local control. That reading misses its deeper structure. Subsidiarity is not a bumper sticker for administrative decentralisation. It is a principle about the right location of authority.

Subsidiarity holds that decisions affecting human beings should be made at the lowest level of social organisation genuinely competent to make them well. Higher levels of authority, whether state or capital, regulator or investor, should not absorb what lower levels can handle responsibly. But the principle has two sides. Higher authority should not dominate where lower authority is adequate, and it must act where lower authority is overwhelmed, captured, or failing. Subsidiarity is therefore not merely a defence of local autonomy. It is a principle about the right ordering of authority at every level.

Applied to economic life, subsidiarity carries significant implications. Families, workers, communities, and enterprises should retain real decision-making power over their own affairs, rather than having decisions made for them by distant capital, managers, or policy-makers who lack the local knowledge and bear none of the consequences. This connects directly to Hayek's insight about dispersed knowledge: those closest to a situation often know something centralised authority cannot. But subsidiarity goes further. It insists that this is not merely an efficiency point about knowledge; it is a moral point about dignity. People have a right to participate in the decisions that shape their lives.

In practice, subsidiarity challenges the typical structure of financialised corporate governance, in which distant shareholders or fund managers make decisions that determine the lives of workers, communities, and the ecosystems they depend on. It supports worker participation in governance, not only because workers may know more about operations than distant owners, but because they are the human beings whose lives are most directly shaped by what the enterprise does. It supports community voice in corporate decisions that affect local conditions. And it supports investor relationships in which capital participates in enterprise without dominating it, where founders, managers, workers, and communities retain genuine authority over purpose and direction.

Subsidiarity is not a licence for local tyranny, nor a denial that higher authority sometimes must act. When local actors are themselves extractive, when communities discriminate, when firms exploit workers who lack alternatives, when corruption shields the powerful from accountability, higher authority must intervene. The principle is not that local is always right. Anyone who has attended a local meeting knows that proximity alone is not sanctification. The principle is that authority should be located where knowledge, responsibility, and accountability are greatest, and that higher authority should support rather than supplant what lower authority can do well.

\section*{Solidarity: The Structural Alternative to Extraction}

If subsidiarity governs the distribution of authority, solidarity governs the distribution of care and obligation. It is one of the most demanding concepts in Catholic Social Teaching, and one of the most easily sentimentalised. Solidarity is often treated as a warm feeling toward humanity, ideally accompanied by a photograph of linked hands and a tasteful font. But in the tradition, solidarity is much more demanding than sentiment.

Solidarity is not charity. Charity is the personal response to another's need; it is indispensable and must not be despised. But solidarity is something more structural. It is the recognition that human beings share a common fate, that the flourishing of each is genuinely bound up with the flourishing of all, and that this interdependence creates obligations that are not merely personal but institutional. Solidarity asks not only what I am willing to give to someone in need, but what kinds of institutions, governance structures, and legal frameworks are required to protect the dignity and agency of all who participate in economic life.

Pope John Paul II used the language of solidarity to describe a commitment that goes beyond sentiment to institutional form. A society marked by solidarity designs its institutions so that no one suffers harm simply because of unequal power or dependency. A firm marked by solidarity does not merely offer charitable programmes alongside extractive operations; it governs itself in ways that protect workers, customers, suppliers, and communities from being dominated. A financial system marked by solidarity designs repayment around the borrower's real capacity and resilience, not merely around what can be legally extracted.

The contrast with extraction runs through everything. Where extraction profits from dependency, solidarity refuses to treat dependency as a revenue opportunity. Where extraction shifts risk onto those least able to bear it, solidarity distributes risk in proportion to power. Where extraction externalises costs, solidarity insists that those who create costs must bear them. Solidarity is therefore not a supplement to productive enterprise but a condition for enterprise that genuinely serves life rather than feeding on it. In a financialised economy, this means that repayment terms must account for the borrower's capacity and resilience, that shared infrastructure and knowledge must remain accessible rather than enclosed for private advantage, and that labour must be treated as a participant in enterprise rather than a cost to minimise. These are not optional expressions of generosity. They are obligations encoded in the structure of institutions that take seriously the common human dignity of all who participate.

\section*{Structural Sin and Economic Systems}

Catholic Social Teaching has developed the concept of structural sin to describe injustices embedded in social, economic, and political systems rather than arising from isolated individual wrongdoing. This is one of the tradition's most analytically useful contributions, because it names something that purely personal accounts of morality cannot adequately address.

Structural sin occurs when the arrangements of society systematically disadvantage some while advantaging others, in ways that persist even when the individuals inside those arrangements are personally decent. The laws and institutions, the markets and governance forms and cultural assumptions, can be tilted long before any individual chooses to tilt them. A manager who genuinely cares about workers may still govern a firm whose legal architecture requires prioritising shareholder return over their welfare. An investor who personally wishes to be patient may manage a fund whose structure demands liquidity and exit. A platform that began with genuine community intentions may be governed by incentives that eventually require monetising dependency. The harm is systemic, not merely personal. Changing the individuals does not change the structure. Replacing one well-meaning manager with another does not fix a governance system that rewards extraction, any more than changing the driver fixes a car designed without brakes.

Financialised capitalism illustrates structural sin with unusual clarity. Debt regimes can trap households while enriching lenders through entirely legal arrangements. Housing markets can privilege investors while excluding residents without any single actor being obviously villainous. Corporate governance can prioritise shareholder value at the expense of workers, communities, and the commons while every individual board member acts within their formal duties. Platform monopolies can enclose shared goods and extract from dependency while pointing to user agreements that were technically accepted. The concept of structural sin therefore demands structural repentance, not only the conversion of individuals but the reform of institutions. It insists that legality does not guarantee justice, and that moral responsibility extends to the systems we design, perpetuate, and benefit from, not only to the direct acts we personally perform. This is why mission-locked corporate structures, repayment terms tied to durable surplus, commons stewardship that prevents enclosure, and accounting that makes externalised costs visible are not add-ons to otherwise neutral institutions. They are responses to structural sin, and they must be embedded in the architecture of economic life.

\section*{The Common Good}

The common good is the organising horizon of Catholic Social Teaching's economic ethics. It refers to the sum of social conditions that allow individuals and communities to flourish. It is not the aggregate of private preferences, nor the welfare of a majority purchased at the expense of a minority, but the genuine set of conditions under which human beings can live well together.

The concept has deep roots in Aristotelian and Scholastic thought. The city, for Aristotle, exists for the sake of the good life. The common good in the Catholic tradition inherits and deepens this insight: society is not merely a mechanism for satisfying individual preferences but a community whose institutions and practices must be ordered toward genuine human flourishing. It encompasses access to work and housing, to healthcare and education, to a clean environment and meaningful opportunity. It includes the goods that cannot be reduced to private ownership, among them ecological health, public trust, open knowledge, and the resilience of communities.

The common good as an economic standard produces a different set of evaluative questions than efficiency or growth. It asks whether markets enhance or undermine human dignity, whether corporations produce value for workers and communities or only for owners, whether finance enables productive activity or subordinates enterprise to exit pressure, whether governance protects shared infrastructure or permits the concentrations of power that enable extraction. It provides a horizon against which economic institutions can be measured, not merely by whether they produce more, but by whether they contribute to conditions in which human beings can actually live well.

The common good cannot be reduced to any single institution's jurisdiction. It is not determined by the state's decisions, revealed by market prices, expressed in shareholder preferences, or made possible by technology alone. It is a moral standard that stands over all of these, against which all of them can be judged. Its pursuit therefore requires plural institutions, each with its proper role: markets that coordinate knowledge and reward useful service, firms accountable to all their stakeholders, states that protect shared infrastructure and prevent domination, a civil society that forms people in the virtues necessary for moral life, and the religious and philosophical traditions that keep alive the question of what human beings are actually for.

\section*{The Practical Implications of the Triad}

Subsidiarity, solidarity, and the common good are not separate principles that happen to coexist in the same tradition. They form a coherent and mutually supporting structure, and they discipline one another.

Subsidiarity without solidarity can become a localism that protects the powerful against accountability at a smaller scale, since a community that governs itself may still exclude, discriminate, or exploit those who lack a voice. Solidarity without subsidiarity can become a centralism that absorbs the proper agency of local actors in the name of their welfare, acting for people rather than with them. Both principles require the common good as their orienting horizon, the standard against which both local authority and collective obligation must be judged.

Together, the three principles translate into concrete requirements for institutional design. Authority over economic decisions should be distributed to those with the most knowledge and the most at stake. Risk and obligation should be shared in ways that prevent the powerful from imposing fragility on the powerless. Institutions, markets, and finance should be evaluated not only by their financial performance but by their contribution to the conditions of human flourishing. These requirements are demanding, but they are not vague. They can be specified in ownership structures, governance documents, repayment formulas, accounting frameworks, labour standards, commons protections, and exit restrictions. The tradition provides the moral grammar. The remaining chapters attempt to give that grammar institutional form.

Catholic Social Teaching does not claim to offer a complete economic blueprint. It does not tell us the optimal interest rate, the ideal corporate governance model, or the right balance between markets and public provision in every context. What it offers is something more fundamental: a set of moral criteria against which economic arrangements can be judged, a vocabulary for naming structural injustice, and a tradition of reflection that refuses to treat economic life as morally autonomous. In a world where capital has become increasingly sovereign, that refusal is itself a form of resistance.

\chapter*{Chapter 15: Sphere Sovereignty, Stewardship, and Moral Markets}
\addcontentsline{toc}{chapter}{Chapter 15: Sphere Sovereignty, Stewardship, and Moral Markets}

Catholic Social Teaching, as the previous chapter traced, provides a powerful moral framework for economic life. But it is not the only tradition that has tried to give stewardship institutional form. Protestant thought, Neo-Calvinist political theory, German ordoliberalism, and the cooperative tradition each approach the same fundamental problem from a different angle, and each contributes something the others need.

This chapter draws on those traditions to argue that stewardship is not merely a personal virtue but an economic category, something that can be encoded into ownership structures, governance documents, financial instruments, and regulatory frameworks. The moral principles established in earlier chapters must become institutional if they are to survive the pressures that consistently push toward extraction. Good intentions, admirable as they are, are not a business model. Nor are they a governance structure, a financing instrument, a voting right, or a legal barrier against sale to the highest bidder. If stewardship is to last, it needs architecture.

\section*{Neo-Calvinist Sphere Sovereignty}

Abraham Kuyper, the Dutch theologian, journalist, and statesman who served as Prime Minister of the Netherlands at the turn of the twentieth century, developed one of the most architecturally sophisticated accounts of how a plural society should organise authority.

His principle of sphere sovereignty holds that human society is not a single unified structure commanded from one centre, but a plurality of distinct spheres, among them family and church, school, enterprise, and state, each with its own proper authority and integrity. No sphere is sovereign over the others in the absolute sense. The state may not absorb the family into itself, the church may not govern the economy, and the market may not colonise education, worship, or the natural world. Each sphere has a domain, and the proper ordering of society requires that each domain be respected.

When applied to markets and corporations, sphere sovereignty resists the tendency of financialised capitalism to reduce everything to a single logic. It insists that the firm is not only an economic entity. It is also a community of work, a participant in civic life, a user of ecological systems, and a structure that forms the people inside it. Its responsibilities do not end at the edges of its financial claims. The enterprise inhabits several spheres at once and must answer to each. Sphere sovereignty also provides a principled account of pluralism in ownership and governance. It recognises that cooperatives and mutuals, mission-locked firms and family businesses, foundation-owned enterprises and social ventures each represent a different ordering of authority appropriate to a different purpose. No single model should dominate all markets. Diversity in ownership forms preserves resilience, protects against the concentration of power, and honours the genuine differences in what different kinds of enterprise are for. A farm, a hospital, a software commons, a family business, a bank, and a monastery do not need the same governance form simply because all of them must somehow deal with money. They are different creatures, and different creatures require different habitats.

Kuyper's framework also illuminates why the sovereignty of capital over the firm is a specific form of disorder. When financial claims come to govern not only the enterprise's financial decisions but its mission, its workers' lives, its relationships with suppliers and communities, and its effects on the natural world, capital has overstepped its proper sphere. It has claimed authority that belongs elsewhere and subordinated other spheres to a single logic. The task of recovering moral economy is, in Kuyperian terms, the task of restoring the proper authority of each sphere over its own domain.

\section*{Ordoliberalism and the Moral Constitution of Markets}

A different but complementary approach came from the German ordoliberal tradition, associated with Walter Eucken and the Freiburg School, which deeply influenced the post-war West German social market economy.

Ordoliberalism begins from a premise that distinguishes it sharply from both laissez-faire liberalism and state-directed planning: markets are not natural. They do not arise spontaneously from human nature and maintain themselves without institutional support. They are constitutional achievements, created and sustained and shaped by the legal, institutional, and cultural frameworks in which they are embedded. There is no such thing as ``the market'' in the abstract, floating above history like an economic weather system. There are only markets constituted by rules, by institutions and cultures and enforcement mechanisms, and someone always helped design them, even if they now pretend otherwise.

Different rules of property, contract, competition, liability, corporate governance, labour law, and monetary policy produce different markets with different moral characters. The choice of rules is therefore a moral and political choice, not merely a technical one. Eucken and his colleagues were deeply suspicious of both the unregulated market power that had produced the cartels and monopolies of the Wilhelmine and Weimar periods, and the state planning that had produced the pathologies of fascism and Soviet-style socialism. Their answer was a third path: an economy whose framework rules were designed to make markets genuinely competitive, to prevent the concentrations of power that could corrupt both the market and the state, and to embed the social protections that allowed people to participate in economic life without being destroyed by its vicissitudes. The social market economy of postwar Germany, with its codetermination requirements, its competition law, its strong labour protections, and its institutional constraints on financial short-termism, was in significant part the institutional expression of ordoliberal thought.

For our purposes, ordoliberalism contributes a crucial insight: stewardship at the level of individual firms and actors is not enough if the institutional framework of markets rewards extraction and punishes stewardship. A founder who wants to build a durable, mission-driven enterprise may find that the capital available, the corporate law applicable, the governance norms expected, and the exit pressures imposed all push systematically toward the opposite. Personal virtue cannot survive indefinitely against institutional structure. A morally serious entrepreneur placed inside an extractive market architecture will become heroic, exhausted, compromised, or outcompeted, and heroism is a poor basis for an economy. Moral markets require moral architecture.

\section*{Cooperative and Mutualist Traditions}

While ordoliberalism worked at the level of market-wide framework conditions, the cooperative and mutualist traditions have worked at the level of the individual enterprise, asking what organisational forms can embed stewardship into the structure of business itself.

The cooperative tradition has a long and varied history, running through workers' cooperatives and agricultural cooperatives, consumer cooperatives, credit unions, and mutual insurers. Each represents a different application of the same core idea: that the people most directly affected by an enterprise should have meaningful control over it, that surplus should serve the enterprise and its members rather than flowing primarily to external investors, and that the governance of productive activity should be accountable to the people who do the work and use the services. The Mondragon cooperative in the Basque Country is perhaps the most studied large-scale example, a federation of worker-owned cooperatives that has maintained significant employment, competitive performance, and mission fidelity over decades. Rabobank in the Netherlands, originally a federation of agricultural credit cooperatives, shows what mutual finance can become at scale. The John Lewis Partnership in Britain represents employee ownership in retail. Building societies and credit unions across the English-speaking world have historically served their members with terms not available from profit-maximising banks.

None of these is perfect, and cooperative forms face genuine challenges: governance grows more complex as scale increases, capital is harder to access, and competitive pressure can erode mission over time. A cooperative is not magically immune to bureaucracy, bad incentives, or the ancient human capacity to make meetings unbearable. But these traditions demonstrate that member ownership and democratic governance can produce economically viable enterprises with significantly different moral characteristics than investor-owned firms. The mutualist tradition, related but distinct, adds the insight that vulnerability need not be monetised by commercial actors. It can be shared among those who face it together. When people pool their risks and govern their protective institutions themselves, the logic shifts from extraction to mutual service.

\section*{Christian Entrepreneurship and the Vocation of Enterprise}

The traditions gathered in this chapter share a common concern with what earlier chapters called vocation: the idea that economic activity can be genuinely meaningful when it is ordered toward purposes beyond private gain, and that this meaning must be institutionally protected or it will be progressively eroded.

Protestant and Neo-Calvinist thought contributed a specific account of entrepreneurship as vocation, not merely as risk-taking in pursuit of profit, but as the creative application of gifts and foresight in service of real human need. The entrepreneur, on this account, is not simply an optimiser, a pitch-deck generator, or a person who has discovered how to say ``scalable'' with unusual conviction. She is a steward of resources, talent, relationships, and opportunity, accountable to the people whose lives her enterprise affects: workers and customers, suppliers, communities, and future generations. Her success is not measured only by financial return, but by whether the enterprise genuinely serves the purposes that justified its creation.

This account of entrepreneurship is demanding. It requires not only personal virtue but institutional support. A founder who understands herself as a steward needs governance structures that protect her mission when investors arrive with different priorities, capital structures that do not require exit on a timeline incompatible with building something durable, and cultural norms that honour long-term stewardship rather than celebrating only liquidity events. Without these supports, the language of vocation can become private consolation while the structures of extraction operate undisturbed. The moral formation of entrepreneurs, though necessary, is not sufficient. Virtue formed in isolation from structure will be progressively worn down. The goal is not to produce morally superior individuals operating within extractive systems, but to build systems in which moral entrepreneurship is the economically viable path rather than the heroic exception.

\section*{Stewardship as an Economic Category}

Across all these traditions, from Catholic Social Teaching and sphere sovereignty to ordoliberalism, cooperative enterprise, and Protestant vocation, a single central insight emerges: stewardship can be institutionally encoded. It is not merely a personal disposition that depends on the character of whoever happens to be in charge. It is a structural feature of ownership, governance, financial design, and regulatory framework, something that can be built into the architecture of economic life.

This changes what moral economy is doing. It is not primarily a programme of exhortation, asking people to be kinder, more generous, or more spiritually attentive, though none of these should be discouraged. It is a programme of institutional design, asking what legal forms, governance structures, financial instruments, and market rules would make stewardship durable and extraction difficult. The question is not only what people should do, but what institutions should require, enable, and protect.

Governance structures can embed purpose through boards that represent multiple stakeholders, through fiduciary duties redefined to include workers and communities, and through voting structures that separate governance rights from financial claims. Ownership rules can prevent extraction through asset locks that protect mission from liquidation, through non-tradable shares that preserve governance stability, and through foundation ownership that removes the enterprise from speculative acquisition. Financial design can align capital with productive outcomes through repayment caps that define when capital has been fairly compensated, through risk-sharing agreements that distribute downside proportionately, and through financing tied to genuine capacity rather than rigid exit schedules. None of these mechanisms is self-sustaining. Every governance structure can be captured, every ownership form can drift, every financial instrument can be misused. The traditions examined in this chapter are not naive about this. But the lesson is not that good architecture is sufficient. The lesson is that without good architecture, personal virtue cannot sustain moral economy against systematic structural pressure.

\section*{Plural Institutions for a Moral Economy}

A moral economy is not a single institutional form. It is a plural ecology of enterprises, governance structures, ownership models, financial instruments, and regulatory frameworks, each appropriate to its context and purpose, together creating conditions in which stewardship is viable and extraction is constrained.

This pluralism is not moral relativism. Not every institutional form is equally aligned with human flourishing. Extractive structures cause real harm even when they are legally permitted and even when financial metrics call them a success. But different human goods require different forms of organisation, and no single model, however well-designed, can adequately serve all of them. The family is not the state. The cooperative is not the foundation-owned firm. The public utility is not the social enterprise. Each exists for its own reasons, and each requires its own form of governance, accountability, and protection.

What the traditions in this chapter share is not a single blueprint but a common orientation: that the economy exists to serve human beings and the communities they form, that capital is a useful servant and a dangerous master, that stewardship requires both personal commitment and institutional support, and that the design of economic institutions is a moral task requiring wisdom, humility, and sustained attention. The next chapter turns to the specific institutional forms through which these principles take concrete shape: steward ownership, mission locks, evergreen finance, and non-extractive capital.

\chapter*{Chapter 16: Non-Extractive Capital}
\addcontentsline{toc}{chapter}{Chapter 16: Non-Extractive Capital}

The previous chapters recovered a moral vocabulary. Now the question becomes practical. How can capital be structured so that it serves life?

Non-extractive capital is not ``ethical investment.'' It is not philanthropy with a revenue model. It is not ESG, impact branding, or a moral preference draped over conventional ownership like a nice linen napkin over a machine designed to eat furniture. It is an attempt to redesign the relationship between capital, enterprise, governance, surplus, and mission. The central premise is simple: capital may receive a fair return for real risk, service, and patience, but it must not gain unlimited power over the enterprise, its mission, its workers, its community, or its future. Non-extractive capital asks capital to do something difficult: to serve without becoming sovereign.

\section*{The Problem It Tries to Solve}

Many mission-driven companies face a painful choice. They need capital in order to build, to hire and develop technology, to reach customers and weather uncertainty, to invest before revenue is stable. Without it, they may remain fragile or fail before they have the chance to serve. But the capital available to them often comes with expectations that threaten the very mission they are trying to preserve.

Conventional venture capital seeks outsized returns through growth and exit. This model may work for companies capable of exponential scaling and winner-take-most dynamics. But it is poorly suited to the much larger number of enterprises that are useful, durable, cash-flowing, and oriented toward stewardship rather than domination. A company may be good but not venture-scalable: profitable but not exit-oriented, serving a real need but incapable of generating monopoly returns, requiring patience rather than hypergrowth, designed to preserve a commons rather than enclose one.

When such a company accepts capital built around exit logic, the company is gradually reshaped. Growth comes to matter more than resilience, valuation more than usefulness, investor rights more than mission, exit more than continuity. This is not always because investors are malicious. They may be sincere and even mission-aligned. But structure speaks louder than intention. A fund designed to return capital through large exits must seek large exits. A governance structure that gives capital the final word will eventually make mission conditional on capital's consent. The result is structural drift, and non-extractive capital tries to solve it by changing the underlying architecture.

\section*{Steward Ownership}

Steward ownership begins from a simple but radical idea: a company should be controlled by people who are responsible for its purpose, not by whoever can buy the financial rights.

In conventional ownership, control and economic upside are bundled together. Those who own shares can influence or control the company, and those shares may rise in value, be sold, or confer rights to future profits. This creates a powerful incentive, since control exists for the benefit of financial ownership. Steward ownership separates these elements. Control is held by stewards bound to the company's mission, economic rights are limited, and the company is not treated as a speculative object to be bought, sold, and liquidated for maximum gain. Profits can still exist. Capital can still be repaid. Workers can still be paid well. The company can still compete, grow, and occasionally discover that printers remain mankind's least redeemed technology. But the governance of the company is protected from becoming an instrument of extraction.

A steward-owned company exists to fulfil its purpose over time. It is not ownerless, not unmanaged, not allergic to profit. But profit is placed in service of mission, not mission in service of profit. The steward is not an absolute owner. The steward is entrusted with governance, accountable to purpose, to the people the enterprise serves, and to the future. In a steward-owned company this obligation is embedded in legal architecture rather than left to conscience: transfer restrictions and mission locks, golden shares, foundation ownership, perpetual purpose trusts, and asset locks can all be used to prevent sale, mission drift, or speculative control. Steward ownership does not guarantee virtue. Stewards can fail, can grow complacent or self-serving or ideologically rigid. But steward ownership changes the default. If capital cannot buy control, it cannot easily force exit. If mission is locked, it cannot easily be sold away. If governance is entrusted to stewards, the company can make long-term decisions that ordinary capital markets might punish.

\section*{Mission-Locked Governance}

Mission is fragile when it depends only on founder intention. Founders may begin with clarity and courage, but founders age, burn out, disagree, die, sell, or get outvoted. Investors enter. Boards change. Crises arrive. A mission that lives only in the founder's heart is vulnerable, however noble the heart and however determined the founder's LinkedIn posts.

Mission-locked governance turns purpose into structure. It asks what must be legally, financially, and institutionally true so that the company's mission can survive pressure. A mission lock may define the company's purpose in the governing documents, restrict sale to buyers who would violate that purpose, require the consent of a guardian entity before major changes, prevent the distribution of assets beyond defined limits, give veto rights to a foundation or purpose trust, or protect worker participation and commons commitments. The point is not to freeze the company forever. Good governance must allow learning and adaptation, and markets change. But adaptation is different from liquidation. A mission lock says that the company may change how it serves its purpose, but it may not abandon the purpose in order to maximise financial return.

This matters most when capital enters. Investors often ask for protections: information rights and repayment rights, covenants, board seats, liquidation preferences, drag-along rights, and exit provisions. Some protections are reasonable, and investors should not be deceived or recklessly disregarded. But investor protections can become mission threats when they allow capital to override purpose. Mission-locked governance seeks a difficult balance. Capital should be protected against fraud and negligence, but mission should be protected against capital's impatience, its appetite for control, and its hunger for liquidity. The company must be governable. The investor must be treated fairly. The mission must remain sovereign over both. This is governance as moral architecture.

\section*{Evergreen Finance}

Conventional venture capital relies on exit. Investors put money into a fund, the fund invests in companies, the companies are expected to grow rapidly, and a small number of large exits return the fund within a limited cycle. This model produces a particular kind of company, one shaped for rapid growth, scalable technology, defensible moats, and exit potential. It can support genuine innovation, but it also pressures companies toward monopoly dynamics, enclosure, aggressive scaling, and financial narratives that may not fit their real mission.

Evergreen finance follows a different logic. An evergreen fund is designed to continue over time rather than to liquidate all its investments within a short cycle. Returns can be recycled, capital can remain patient, and the fund can support companies that generate durable surplus rather than spectacular exits. This matters because the time horizon is itself moral architecture. A short horizon pushes toward extraction, asking how value can be increased and realised soon. A long horizon asks how the enterprise can become durable, resilient, and capable of producing surplus without harming its mission.

Evergreen finance is not automatically non-extractive. A permanent fund can still extract indefinitely if its return expectations are uncapped or its governance is not mission-aligned. Permanence can become another form of domination, just with better stationery. But when joined to repayment caps, mission locks, and steward ownership, evergreen finance can support patient enterprise. The Monti di Pietà faced a similar structural tension: the fund must survive in order to serve, but if the preservation of capital becomes the highest good, mission drifts. A non-extractive evergreen fund must preserve itself for mission, not preserve mission as a story that justifies the fund. This requires financial discipline as well as moral discipline.

\section*{Repayment Caps and Anti-Usury Logic}

A repayment cap defines the maximum return capital can receive. It is one of the simplest and most morally significant mechanisms in non-extractive finance. It says that capital may be compensated, but its claim is not infinite. Once the agreed return has been reached, the financial claim ends or changes form, and capital does not own the enterprise's future without limit.

This addresses the ancient problem of usury in contemporary form. The Scholastic and Franciscan traditions asked when return on money becomes morally excessive, especially when it profits from need or generates gain without real participation in production. They shared the conviction that financial return requires moral justification and must remain proportionate to contribution. A repayment cap gives that conviction structural form. It acknowledges that investors may deserve return, since they may have taken real risk, provided capital when the company needed it, waited, and made productive activity possible. A fair return is not extraction. But the cap says that the investor's claim cannot expand forever simply because the company succeeds. It cannot become a permanent tax on the enterprise. It cannot convert temporary risk into perpetual ownership. The cap must be set carefully: too low and investors may not participate, too high and it becomes symbolic, too rigid and the company may be harmed during fragile periods. The moral task is to design a cap that reflects risk, time, contribution, and mission. A repayment cap is therefore not merely a number. It is a statement about the moral limits of capital.

\section*{Durable Surplus Instead of Exit Value}

The most important shift in non-extractive capital is the shift from exit value to durable surplus. Conventional venture capital asks what the company will be worth to a future buyer or public market. Non-extractive capital asks whether the company can generate durable surplus while preserving its mission. These are profoundly different questions.

Exit value depends on what someone else will pay for the company, shaped by growth narratives, market comparables, revenue multiples, and future expectations. It can reward real value creation, but it can also reward monopoly potential, dependency, enclosure, and speculative hype. Durable surplus depends on the company's ability to generate more income than it needs for healthy operation, resilience, reinvestment, fair compensation, and mission integrity. It asks whether the company can live. Not whether it can be sold, not whether it can dominate, but whether it can produce enough recurring surplus from real service to repay capital without harming itself.

A company that grows steadily, pays its workers well, serves its customers honestly, maintains reserves, contributes to commons, and repays its investors over time may be highly successful, even if it never exits. Durable surplus is not mere profit. It must be what remains after the company has met its obligations to mission, workers, resilience, maintenance, customers, and future capacity. If repayment to investors comes only by underpaying workers, neglecting infrastructure, extracting from users, or externalising costs, then the surplus is not durable. It is disguised depletion. This is why moral accounting matters: a repayment formula must distinguish between cash that can rightly be distributed and cash needed for resilience.

\section*{Productive versus Extractive Capital in Practice}

The distinction between productive and extractive capital becomes practical in the investment process. A non-extractive investor must ask different questions from a conventional one. Not only how large the market is, but what real need is being served. Not only whether this can scale, but what happens if it does. Not only what the revenue model is, but whether that model depends on dependency, lock-in, addiction, opacity, or enclosure. Not only what the exit path is, but how capital can be repaid without exit. These questions change how a portfolio is built.

Non-extractive capital may favour companies with real revenue potential, durable customer relationships, moderate growth, useful products, open-source or infrastructure value, founder commitment to mission, and the capacity to repay from surplus. It may avoid companies whose economics depend on monopoly capture, user manipulation, labour arbitrage, speculative exit, regulatory evasion, or the enclosure of commons. A non-extractive company can still be ambitious. It can grow, innovate, hire, compete, and build at scale. But its ambition is ordered. It does not seek scale for its own sake or treat domination as success. The practical challenge is that non-extractive investing requires being financially literate and morally literate at the same time. If only morally literate, it will fund beautiful but unsustainable projects. If only financially literate, it will reproduce extraction with better language. The discipline lies in holding both together.

\section*{Governance as Moral Architecture}

Governance determines whether non-extractive intention can survive pressure. Many companies and funds begin with good intentions, but pressure reveals architecture. When revenue slows, when investors become impatient, when an acquisition offer arrives, or when reserves run low, governance decides what happens. Governance is where moral claims become operational authority.

A non-extractive company needs governance that can answer clearly. Who has the right to say no to an acquisition? Who decides whether surplus is distributed or reinvested? Who protects mission when founders change? These questions cannot be left vague, and the answers must be embedded in documents rather than trusted to personality. Non-extractive governance must define the mission in operational terms, specify which decisions require mission review, make investor repayment rights transparent, limit sale and asset transfer, and clarify how conflicts are resolved. Mission locks can become rigid, steward bodies can become self-perpetuating, and investor rights can be too weak to attract capital or too strong to protect mission. The goal is not maximum complexity. The goal is credible protection: governance that makes the right thing easier and the extractive thing harder.

\section*{Moral Accounting and the Repayment Formula}

A non-extractive fund needs accounting that can distinguish healthy repayment from disguised extraction. But this must be practical. If it is too complex, it will fail, becoming either bureaucracy or a very expensive fog machine. The goal is to make visible the few moral distinctions that matter most to the financing relationship.

The central question is whether a company can repay capital from durable surplus without weakening its mission, its resilience, its workers, its users, its commons, or its future capacity. Operating needs come first, then fair wages and essential obligations, then maintenance and security, then taxes, then necessary reinvestment, then resilience reserves, and only then should investor repayment come from whatever surplus remains. The principle is constant: capital should be repaid from surplus, not from depletion.

The key risk to guard against is false surplus, the surplus that appears only because the company is underpaying workers, ignoring maintenance, accumulating technical debt, exploiting customers, or depleting the founder's health. Such repayment is extractive even when the spreadsheet looks healthy. A non-extractive fund must refuse repayment that comes from hidden harm. This is the Monti problem in modern form: the fund must be preserved, but not by consuming the borrower. Moral accounting is the practice of keeping capital honest.

\section*{Workers, Users, and Communities}

If non-extractive capital only changes the relationship between investors and founders, it is incomplete. Capital can be non-extractive toward founders while companies remain extractive toward workers, users, suppliers, communities, or commons. A founder can retain mission control and still underpay workers. A steward-owned firm can still exploit its users. A repayment cap can still be serviced through customer lock-in.

Non-extraction must therefore reach beyond financing to ask how the enterprise treats all those from whom value is drawn. Workers should not be treated merely as costs that make repayment possible. Users should not be locked into dependency so that revenue becomes predictable. Open-source ecosystems should not be mined without contribution. The repayment formula must be linked to the broader health of the company, because a company that repays its investors while increasing worker precarity, enclosing a commons, or designing for addiction is not truly non-extractive. The practical answer is to identify each company's material extraction risks. For a software company these might be user lock-in, data extraction, open-source dependency, technical debt, and worker burnout; for a manufacturing firm, labour conditions, supply-chain pressure, and ecological impact. Non-extractive finance must be contextual. It should ask where this particular company is most tempted to extract, and how it will know if it does.

\section*{Open Source and Non-Extractive Finance}

Open source is one of the most important test cases for non-extractive capital. Open-source software creates enormous shared value, supporting public infrastructure, private companies, governments, universities, and daily digital life. Yet many maintainers are not adequately supported, and corporations often depend heavily on open-source infrastructure while contributing less than they extract. Conventional venture capital often worsens this by pushing open-source companies toward enclosure: proprietary features, hosted services, cloud lock-in, dual licensing, or eventual acquisition.

A project begins as a commons, attracts users, forms a company around itself, accepts capital, faces growth expectations, and discovers that defensibility seems to require monetising the very community that created the value. This is a profoundly modern tragedy: build a commons, fund it, then slowly discover that the funding model wants a tollbooth. Non-extractive capital offers another path. It can fund open-source companies in ways that do not require enclosure or exit, through repayment drawn from durable surplus generated by services, support, hosting, training, or enterprise contracts, through capped returns, through governance that protects the open-source mission, and through steward ownership that prevents sale to a proprietary acquirer. Shared digital infrastructure should not have to choose between starvation and enclosure. Non-extractive finance can help create a sustainable third path.

\section*{The Risk of Moral Branding}

``Non-extractive'' is a powerful phrase, and precisely because it is attractive, it can be diluted. A fund may call itself non-extractive while demanding high returns. A company may call itself steward-owned while leaving its mission weakly protected. An investor may accept a repayment cap set so high that it functions like uncapped extraction. A platform may call itself commons-oriented while quietly capturing user dependency. Extraction is skilled at wearing the clothes of reform.

This is why definitions, structures, and audits matter. Non-extractive cannot mean good intentions. It must mean something testable: a capped financial claim, control separated from speculative ownership, mission legally protected, repayment tied to durable surplus, transparency about fees and governance rights, and restrictions on sale against mission. The word must be disciplined by form. Not because form guarantees virtue, but because without form, moral language will be absorbed by the very system it critiques.

\section*{Failure Modes}

Non-extractive capital has its own failure modes, and they should be named honestly. The first is undercapitalisation: if return expectations are too low, liquidity too limited, or investor education too weak, the fund may not raise enough capital to matter. The second is naive underwriting: if the fund overestimates a company's repayment capacity, it may deplete itself or pressure companies during fragile periods. Non-extractive finance still requires rigorous financial analysis. Kindness is not an underwriting methodology.

Hidden extraction is a persistent risk, since portfolio companies may generate surplus through worker burnout, customer lock-in, undermaintenance, or commons free-riding that the fund fails to see. Governance rigidity can prevent necessary adaptation or trap companies in obsolete models. Founder capture, protecting companies from investors while inadvertently protecting founders from accountability, is a related danger. And moral complacency, the assumption that once a structure is labelled non-extractive the moral work is done, is perhaps the most insidious failure of all. These failure modes do not invalidate the project. They clarify the work. A serious non-extractive fund must be financially competent, legally careful, morally disciplined, and institutionally humble. The Franciscans taught us that institutions with good purposes can slowly become what they were designed to oppose. The lesson is not to avoid building institutions, but to build them with that vulnerability fully in view.

\section*{Non-Extractive Capital as Moral Economy in Practice}

Non-extractive capital is not the whole of moral economy. It will not solve housing financialisation by itself, reform all corporations, replace public policy, or eliminate the need for regulation, labour organising, taxation, antitrust, public investment, commons governance, or spiritual conversion. But it is one concrete answer to a central problem: how can mission-driven enterprises access capital without becoming instruments of extraction? In this sense, it is moral economy in practice.

The ancient concerns mapped in earlier chapters become, in non-extractive finance, things you can actually write into a contract. The concern with accumulation without limit becomes a repayment cap. Scholastic proportionality becomes return linked to contribution and risk. Franciscan productive capital becomes patient investment in real enterprise. Protestant stewardship becomes steward ownership and mission-locked governance. Smith's suspicion of rent-seeking becomes anti-extraction due diligence. Catholic Social Teaching becomes governance ordered toward dignity, subsidiarity, solidarity, and the common good. The question ``what is an economy for?'' eventually becomes a more pointed one: what is in the contract?

\section*{Capital That Serves Life}

Capital that serves life must be converted from master to servant. This conversion requires limits and purpose, governance and patience and moral accounting. It requires investors willing to accept a fair return rather than maximum extraction, founders willing to accept stewardship rather than personal sovereignty, and companies willing to generate surplus without hidden depletion. Non-extractive capital is one attempt to make this conversion real.

Its goal is not to end capital, but to redeem its role. Capital is stored capacity. It can help human beings build across time, fund what does not yet exist, preserve what would otherwise be lost, and support enterprise, infrastructure, and common life. It can serve, but only if it is ordered. Stewarded capital asks what it is for; unordered capital seeks more. Stewarded capital accepts limits and supports continuity; unordered capital demands control and seeks exit. One strengthens agency; the other extracts from dependency. This is the difference between capital that feeds on life and capital that serves it, ancient in structure, contemporary in form, and urgent in this moment.

The next chapter tests this framework in the defining terrain of our time: digital infrastructure, open source, artificial intelligence, and the commons, where every question this book has raised appears again, at higher speed, and with considerably less margin for error.

\chapter*{Chapter 17: The Digital Commons and the Future of Civilisation}
\addcontentsline{toc}{chapter}{Chapter 17: The Digital Commons and the Future of Civilisation}

The digital world is often described as new. New technologies and platforms, new forms of intelligence and labour, new powers of computation and prediction and automation. In many ways this is true. Software mediates daily life, data shapes decisions, platforms govern visibility, and algorithms sort attention. Artificial intelligence now enters domains once reserved for human judgment, craft, analysis, and care.

But beneath the novelty lies an ancient question: will the infrastructure of human life be stewarded as a common good, or enclosed as a source of extraction? Every tradition traced in this book feared some version of this, from Aristotle's concern about money detached from use, to the Franciscan worry about finance detached from need, to Smith's warnings about concentrated commercial power, to the lesson of financialised capitalism that capital can turn everything into future claim. The digital age gives this old drama new tools, new scale, and new stakes. The server racks are new. The temptation is ancient. The central question remains unchanged: what is this economy for?

\section*{What Makes Infrastructure a Commons}

Infrastructure is what other things depend on. Roads and ports, schools and hospitals, water systems, legal systems, languages, and scientific knowledge all make other forms of life possible. They are not merely products; they are conditions of participation. Digital infrastructure now plays this role. Operating systems and programming languages, encryption libraries, cloud services, communication protocols, payment rails, and AI models shape what can be built, who can participate, and who controls access.

A good becomes commons-like when many depend on it, many contribute to it, and its value exceeds any single owner's contribution. A commons is not merely an unowned resource, which would be a recipe less for flourishing than for everyone standing around wondering who is supposed to fix the roof. It is a shared field of use that requires governance, maintenance, access rules, contribution norms, and protection against depletion or capture. This is especially clear in open-source software. A small library maintained by volunteers may become a dependency for thousands of companies. A programming language may become the foundation for entire industries. A security tool may protect public institutions, hospitals, activists, and journalists. The value is distributed, the dependency is broad, and the maintenance burden may fall on very few. Such infrastructure cannot be understood well if we see it only as private property, volunteer hobby, or market input. It is a commons, and commons require stewardship.

\section*{Open Source as Stewardship in Practice}

Open source is one of the clearest contemporary examples of capital serving life, when it is well governed and sustained. At its best, open source allows knowledge to circulate freely, lets people inspect and modify and improve the tools they use, lowers barriers to entry, and supports education, security, interoperability, and public trust. It can be a form of digital hospitality: here is a tool, a library, a protocol; you may use it, learn from it, improve it, and build with it. The gift does not disappear when shared. In many cases, it becomes more valuable.

This is not romanticism. Open-source communities can be exclusionary, underfunded, badly governed, or dominated by powerful actors. Maintainers burn out. Security vulnerabilities persist. Corporate free-riding is common. But open source still reveals something important: not all value creation requires enclosure. Shared creation can generate enormous value, and collaboration, reputation, curiosity, public spirit, and vocation can all motivate productive work. The weakness is the weakness of a gift without an institution. A commons can be used without contribution. A maintainer can be praised but unpaid. A community can create value while corporations capture it. Open source needs funding models and governance bodies, maintenance institutions, procurement policies, and business models that do not simply enclose the commons they depend on. Otherwise it becomes an invisible subsidy to extractive capital. The moral problem is not that companies use open source, since use is part of the point. The problem is when use becomes extraction: corporations building profitable services on shared infrastructure while contributing little, or depending on maintainers without helping to sustain them. Open source is stewardship in practice only when the gift is protected from becoming raw material for extraction.

\section*{Digital Dependency and the Sovereignty Question}

Digital dependency is often invisible until something breaks. A public agency cannot migrate from a proprietary vendor. A hospital depends on software it cannot inspect. A small business loses access to its customers when a platform changes the rules. A city cannot operate critical services without foreign-owned systems. A company built on open-source software discovers that the maintainer has burned out. Digital infrastructure creates relationships of dependence, and dependence creates moral obligations.

This does not mean every dependency is bad. Human life is interdependent, technical systems require specialisation, and no institution can build everything itself unless it has a great deal of money and a suspicious amount of time. The question is whether dependency remains accountable: whether users can leave, whether data can move, whether systems can interoperate, whether public institutions can audit what they rely on, and whether access can be withdrawn arbitrarily once lock-in is complete. Dependency becomes extractive when exit is impossible, governance is distant, and access is monetised through lock-in. This is the digital form of rent. The platform or cloud provider, the app store or AI infrastructure company, becomes the gate through which others must pass, extracting recurring claim from dependency rather than from ongoing service. When a corporation controls infrastructure on which public and economic life depends, it exercises quasi-sovereign power, and the question of democratic accountability becomes urgent.

\section*{AI and the Concentration of Power}

Artificial intelligence intensifies the moral questions of digital infrastructure because it concentrates several forms of power at once. AI systems require data and compute, talent, capital, and distribution, resources that are not evenly spread. The most powerful systems tend to be built by institutions with enormous financial resources, specialised hardware, vast data pipelines, and platform reach. This creates a concentration dynamic that compounds itself: more data improves systems, better systems attract more users, more users generate more data and revenue and bargaining power, more capital buys more compute, and more compute raises the barrier for anyone else.

AI therefore risks becoming not merely a technology but an infrastructure of dependency. Shared human knowledge becomes training material, public language becomes model input, creative work becomes dataset, compute becomes chokepoint, and subscription becomes dependency. This does not mean AI is inherently extractive. AI can serve life. It can assist research and medicine, education and accessibility, administration and security and scientific discovery. It can help small teams do work that once required large institutions. But whether AI serves life depends entirely on how it is governed, owned, financed, deployed, and made accountable. An AI system trained on shared knowledge but governed as a private empire raises serious moral questions. An AI system that displaces workers while concentrating gains in capital raises questions of justice. The question is not whether AI is powerful. The question is what the power is for.

\section*{AI as a New Chrematistics}

Aristotle's chrematistics was the art of accumulation without limit. AI can become chrematistic when intelligence itself is reorganised around extraction, when computational capability is directed not toward wisdom, care, or service, but toward optimising prediction, capture, and monetisation on behalf of whoever owns the system. AI can optimise advertising and pricing, surveillance and labour discipline, addiction and persuasion, financial trading and regulatory evasion, just as readily as it can optimise healthcare, education, or public safety. The tool amplifies the telos of the institution deploying it.

The moral danger is not that machines think. The danger is that machines intensify the aims of the systems that deploy them. This is why the question of purpose is unavoidable. What is this AI system for? Who owns it, and who governs it? Who is accountable for its harms, and who benefits from its productivity gains? Whose labour trained it, and whose knowledge was absorbed? Without moral economy, AI will be governed by the default answers of financialised capitalism: scale and growth, engagement, efficiency, market share, and return on capital. Those answers are not enough. AI needs a richer moral architecture because it touches not only production but cognition, attention, trust, speech, and imagination. If capital that feeds on life gains control over artificial intelligence, extraction becomes faster, smarter, and harder to see.

\section*{Digital Sovereignty}

Digital sovereignty is often discussed in geopolitical terms: whether states can control their own digital infrastructure, avoid dependency on foreign cloud providers, protect data and democratic authority, and govern technology according to their own laws and values. These are important questions. But digital sovereignty is not only geopolitical. It is moral-economic.

A society is not digitally sovereign if it merely replaces foreign dependency with domestic extraction. It is not sovereign if its public institutions depend on private platforms whose governance they cannot influence. Sovereignty requires the stewardship of shared infrastructure: public institutions need technical competence, procurement models that support open standards, funding for open-source infrastructure, governance that prevents capture, and data rules that protect persons and communities. Digital sovereignty also requires institutional pluralism. Not everything should be state-run, corporate, or volunteer-maintained. A healthy digital ecosystem needs public institutions and mission-driven companies, cooperatives and open-source foundations, universities, standards bodies, and steward-owned firms. Each has a proper role, and no single sphere should dominate the others. The principles traced in earlier chapters, subsidiarity and solidarity, the common good, sphere sovereignty, and non-extractive capital, all apply to digital infrastructure. Digital sovereignty is not only a strategy. It is a moral economy of infrastructure.

\section*{The Governance of Shared Resources}

Commons do not govern themselves automatically. Romantic language about the commons can obscure the hard work of stewardship. Shared resources require rules and roles, maintenance and boundaries, conflict resolution, contribution systems, accountability, and renewal. Without governance, a commons can be depleted, captured, neglected, or dominated by informal elites. Digital commons are no different. An open-source project needs maintainers with the authority to accept, reject, review, and release code. Communities need norms, security vulnerabilities need a response, conflicts need resolution, and corporate contributors need boundaries.

A commons can fail by being too open to capture, but also by becoming closed, hostile, or unresponsive. It can unintentionally exclude newcomers. It can burn out its maintainers by expecting infinite unpaid labour. Stewardship requires form. Digital commons may need foundations and cooperatives, trusts and public funding, membership governance, contributor agreements, licences, security funds, and maintainer salaries. These are not bureaucratic distractions. They are the infrastructure of the infrastructure, which is admittedly not the sort of phrase one says at dinner if one wishes to remain popular, but it is true. The principles are clear regardless of the governance model: preserve access, fund maintenance, prevent enclosure, make power accountable, protect contributors, and ensure that those who profit from the commons contribute to its ongoing life.

\section*{Non-Extractive Business Models for Digital Infrastructure}

One of the hardest questions is how digital commons can be financially sustained. Pure volunteerism is fragile, philanthropy inconsistent, government funding slow, corporate sponsorship distorting, and conventional venture capital likely to push toward enclosure. There is no perfect model, which is unfortunate but at least keeps consultants employed. But some models are more aligned with stewardship than others.

A steward-owned open-source company can provide services, support, hosting, security, training, or enterprise features while protecting the core commons from sale or enclosure. Repayment-capped finance can fund growth without requiring acquisition by a proprietary platform. Public procurement can create reliable demand for open infrastructure. Foundations can hold trademarks, governance rights, or core assets. Cooperatives can let users or maintainers share in governance. The key is to align revenue with service rather than dependency. A life-serving digital business model asks how it can be paid to maintain, improve, support, and responsibly steward infrastructure. An extractive one asks how it can control access, capture data, create lock-in, and monetise dependency. A digital infrastructure company claiming to be non-extractive should be able to say clearly what remains permanently open, how its maintainers are paid, whether customers can leave with their data and workflows, whether investors are capped, and how community voice is represented. Without credible answers, ``open'' becomes a marketing word.

\section*{AI, Labour, and the Distribution of Surplus}

AI raises a profound question about surplus. If AI systems increase productivity, who receives the gains? The answer will not be determined by technology alone but by ownership and governance, labour law, bargaining power, public policy, intellectual property, data rights, and capital structure.

If AI is governed by extractive capital, the productivity gains will likely flow upward. Workers will face displacement, surveillance, or wage pressure. Creators will see their work absorbed into models without meaningful compensation. Public institutions will pay recurring fees for systems built partly on public knowledge, and the surplus generated by collective human culture will be captured by a small number of owners of compute, models, and platforms. If AI is governed by stewardship, the gains could be distributed differently, reducing drudgery while preserving worker dignity, improving public services without privatising public capacity, strengthening commons rather than enclosing them, and creating surplus that could fund shorter work time, better care, and shared infrastructure. The difference is not the model alone. It is the political economy around the model. Who owns the AI company, and what return does capital demand? Are workers represented? Are data sources acknowledged? Are the models open or auditable? AI makes the question of capital structure urgent because the technology magnifies whatever logic governs it. The future of AI is therefore not only technical. It is institutional.

\section*{Human Judgment in the Age of Machines}

Moral economy cannot delegate judgment to machines. AI systems can classify and predict, generate and summarise, optimise and recommend. They can assist judgment and reveal patterns. But they cannot answer the foundational question of this book: what is an economy for? That question requires discernment about ends, wisdom about goods, and attention to persons, communities, history, vulnerability, and dignity. It requires responsibility. AI can help calculate means. It cannot define ultimate ends.

This matters because institutions often adopt technology in order to avoid responsibility. An algorithm recommends. A model scores. A decision becomes ``data-driven.'' The human actor then claims to have merely followed the system: a bank denying credit because the model says so, a platform suspending a user because the system says so, a company firing workers because the optimisation says so. Tools can inform decisions, but responsibility cannot be transferred to the tool. A moral economy of AI must preserve human accountability. Institutions using AI must remain answerable for the ends, the data, the design, the deployment, the errors, the exclusions, and the consequences of their systems. The more powerful the system, the greater the obligation. AI should not become a new invisible hand, invoked to absolve human beings from judgment.

\section*{Digital Asceticism and the Formation of Desire}

The digital economy does not only reorganise infrastructure. It forms desire. Platforms compete for attention, recommendation systems shape curiosity, social media trains comparison, notifications fragment time, and metrics teach people what matters: likes and shares, views, rankings, streaks, followers, and engagement. Digital capitalism is a school of desire, and it operates on the same principle every tradition in this book identified: once desire is trained toward accumulation rather than service, it becomes hard to satisfy and harder to redirect.

This connects directly to the older moral traditions. Aristotle worried about desire without limit. Christianity worried about avarice, vanity, and the deformation of the soul. Protestantism worried about discipline, vocation, and Sabbath. Smith worried about admiration and the corruption of moral sentiments. The digital economy intensifies all of these concerns, since it can make vanity measurable, comparison constant, attention monetisable, rest difficult, outrage profitable, and addiction scalable. A moral economy of digital life therefore needs ascetic practices, not the rejection of technology, not nostalgia, not purity, but disciplined limits that protect attention, contemplation, friendship, worship, family, and embodied life from total colonisation. Sabbath becomes newly important. Silence becomes newly important. But personal discipline is not enough. Platforms should not be allowed to build addiction machines and then tell users to practise self-control. Moral responsibility belongs also to designers and companies, to investors and regulators. If a business model depends on capturing attention against the user's good, it is extractive. Capital that serves life must respect human attention. Attention is not raw material. It is part of the soul's capacity to love, think, pray, learn, and perceive reality.

\section*{The Return of the Commons}

The commons returns in the digital age because knowledge behaves differently from scarce material goods. If I give you a loaf of bread, I no longer have it. If I teach you an idea, we may both have it. If I share code, you can use it without preventing me from using it. Knowledge can be non-rivalrous, cumulative, and generative. This does not mean knowledge costs nothing, since creating, maintaining, verifying, documenting, and teaching it require real labour. But the logic of knowledge goods differs from the logic of ordinary commodities.

Enclosing knowledge can sometimes support creation, since intellectual property and proprietary products may incentivise investment. But excessive enclosure slows learning, concentrates power, duplicates effort, and creates dependency. When knowledge becomes infrastructure, enclosure becomes domination. The digital age therefore forces us to rethink property, not to abolish it but to ask of each good what should be privately owned, what publicly governed, what left open, what stewarded by communities, and what funded as public infrastructure. A moral economy cannot answer all these questions with one principle. But it can reject the assumption that maximum private enclosure is the natural endpoint of value creation. The commons is not the absence of ownership. It is ownership, governance, or stewardship ordered toward shared use. The future of freedom may depend on whether we can govern knowledge as a shared inheritance rather than a frontier of capture.

\section*{What Moral Economy Offers the Age of Machines}

The moral economy tradition gives the digital age several gifts that our current frameworks often lack. Most fundamentally, it gives the question of purpose: technology is not self-justifying, and AI, platforms, data systems, and digital infrastructure must be judged by the forms of life they create. It gives the distinction between use and accumulation, between digital tools that serve real human use and systems for the unlimited capture of attention, data, dependency, and claim. It gives the language of stewardship: digital infrastructure is entrusted power, and it must be governed for those who depend on it, not merely for those who own it.

It gives the critique of usury and rent, recognising that recurring claims on dependency and lock-in can become digital forms of old financial sins. It gives the concept of the commons, since shared knowledge, code, protocols, and public digital infrastructure require governance beyond both private extraction and state bureaucracy. It gives the category of moral formation, since technologies form desire, attention, judgment, and community, and are not neutral tools used by unchanged persons. And it gives the principle of enough. Not every system should maximise engagement, not every dataset should be captured, not every process automated, not every public good monetised, not every human capacity optimised for productivity.

These are not a complete policy programme. They are a moral grammar, and that is what we need most. The digital age is moving faster than our moral language. We reach for words like innovation and disruption, efficiency, alignment, and progress. These words matter, but they are not enough. We need older and deeper words: justice and stewardship, commons, vocation, Sabbath, dignity, solidarity, and the common good. Without such words, we will not know whether our machines serve us, or whether we have reorganised ourselves to serve them.

\section*{The Return of the Founding Question}

We have reached the end of the book's main argument. It began with a question: what is an economy for? That question now returns in digital form. What is software for? What is AI for? What is open source for? What is capital for?

If the answer is extraction, then the digital future will intensify the disorders of financialised capitalism. Platforms will become sovereign, AI will concentrate power, commons will be enclosed, public institutions will grow dependent, human attention will be harvested, and knowledge will become proprietary infrastructure. If the answer is stewardship, another future remains possible: digital infrastructure governed as a common good, open source funded and protected, AI augmenting human agency rather than replacing responsibility, companies mission-locked and steward-owned, capital receiving a fair return without enclosing the future, and workers and communities given a voice. Technology serving life.

This will not happen automatically. The default trajectory of financialised capitalism is enclosure. Stewardship requires interruption. It requires saying no to some forms of scale, some forms of capital, some forms of convenience, some forms of automation, some forms of enclosure. But it also requires saying yes: to shared infrastructure and patient funding, to public capacity and open standards, to mission-locked enterprise, to digital commons, to human dignity in the age of machines. The future of civilisation will not be determined by technology alone. It will be determined by the moral architecture through which technology is owned, financed, governed, and loved. The founding question is not behind us. It is ahead of us, and every new technology will ask it again.

\chapter*{Conclusion: Toward an Economy Ordered to Human Flourishing}
\addcontentsline{toc}{chapter}{Conclusion: Toward an Economy Ordered to Human Flourishing}

We began with a question.

What is an economy for?

The question is simple enough for a child to ask, and difficult enough for a civilisation to bury. It lies beneath every market and firm, every loan and investment, every technology, every ownership structure, every policy. It is answered not only in books, but in law and contracts, in capital flows and accounting systems, in board decisions and software licences, in debt covenants and land titles and labour practices, in platform architectures, and in the daily disciplines of work and desire.

An economy is never only a mechanism. It is a moral order. It teaches people what to love. It rewards some forms of life and punishes others. It shapes households and firms, cities and landscapes, institutions and souls. It decides what counts as success and what disappears as cost. It can honour work or degrade it, build trust or dissolve it, strengthen agency or monetise dependency, preserve common goods or enclose them. It can make capital a servant, or crown it as master.

This book has argued that we cannot understand our present crisis until we recover an older distinction: the distinction between wealth that serves life and wealth that feeds on it. That distinction has appeared under many names across the centuries traced in these pages. It is Aristotle's oikonomia and chrematistics, Christianity's stewardship and avarice, the Scholastics' commutative justice set against exploitation, the Franciscan question of whether capital could be productive without becoming usurious, Salamanca's revelation that coordination and justice are not the same thing, the Protestant warning about vocation without Sabbath, Smith's concern about the corrupting power of admiration for wealth, and the digital age's question of whether knowledge and infrastructure will be stewarded as commons or enclosed as capital's next frontier. Across this long story, the question keeps returning: does this serve life?

\section*{Commerce After Extraction}

To recover moral economy is not to reject commerce. Commerce can be honourable. Markets can coordinate human need and creativity. Prices can communicate knowledge. Entrepreneurs can discover possibilities that no planner could foresee, however impressive the planner's credentials and however many laminated strategy diagrams the planner carries. Firms can gather people around useful work. Finance can move resources across time. Profit can sustain enterprise, reward service, compensate risk, and fund future goods. The problem is not commerce. The problem is commerce detached from its proper end.

Commerce becomes extractive when gain no longer corresponds to service: when contract hides domination, when market price masks necessity, when ownership loses responsibility, when finance claims more than it contributes, when platforms monetise dependency, when growth is achieved through depletion, and when the category of enough disappears. An economy after extraction would not abolish markets. It would re-embed them within moral life, asking not only whether something can be sold, but whether the price reflects justice or power; not only whether investment can generate return, but whether the return is proportionate to contribution; not only whether a company can grow, but whether growth strengthens or deforms its purpose.

Commerce after extraction is commerce under discipline. It knows that scarcity is real, that incentives matter, and that coordination is difficult. It knows that good intentions can fail and that institutions must be financially viable. But it also knows that viability is not the same as virtue, and that profitability is not the same as legitimacy. Commerce after extraction asks every enterprise to give an account of its gain.

\section*{Freedom Without Domination}

Modern economies often speak of freedom: the freedom to choose, to contract, to invest, to start a company, to buy and sell. These freedoms matter. An economy without room for initiative becomes stagnant and coercive. People need space to build, trade, cooperate, create, and respond to local knowledge. Markets can protect pluralism by allowing many plans to coexist without requiring one centralised vision.

But freedom becomes thin when separated from power. A borrower may be free to accept a predatory loan because no better option exists. A worker may be free to accept degrading conditions because rent is due. A founder may be free to accept extractive capital because the company will otherwise die. A user may be free to leave a platform only by abandoning the social and professional network built around it. Formal choice is not always real freedom. Freedom requires protection from domination: real alternatives, intelligible terms, accountable power, and institutions that prevent vulnerability from becoming another person's business model. It requires markets, but it also requires laws and norms, commons and unions and cooperatives, public institutions, mission-locked firms, steward ownership, antitrust, social protections, and moral culture. Capital that serves life increases agency. Capital that feeds on life narrows it, creating dependency and lock-in, debt pressure, opacity, artificial scarcity, and forced exit. An economy ordered toward human flourishing must therefore seek freedom without domination.

\section*{Entrepreneurship as Stewardship}

Entrepreneurship has often been celebrated as disruption. Sometimes disruption is necessary. Existing systems may be unjust, stagnant, captured, or incapable of meeting real needs. Entrepreneurs can reveal new possibilities, build tools and organisations that improve life, take risks others avoid, and coordinate resources in new ways. But disruption is not a moral category. Some disruption heals and some destroys. Some liberates and some encloses. Some challenges unjust incumbents, and some merely replaces one form of domination with another, usually after an energetic launch campaign and a heroic quantity of branded tote bags.

Entrepreneurship needs a better name for its highest form: stewardship. The steward-entrepreneur does not ask only what can be built, but what has been entrusted. Not only whether this can scale, but what happens if it does. Not only whether this can attract capital, but what kind of capital it can accept without losing its soul. Stewardship does not make entrepreneurship timid. It makes it accountable. A steward can be ambitious, building institutions and technologies, companies and funds and infrastructure. But a steward knows that power is entrusted for purpose, that control exists for service, and that surplus exists for mission, resilience, and future capacity. Growth must remain ordered. Exit is not the highest proof of success. The company is not prey. The commons is not raw material. Users are not resources to be harvested, workers are not costs to be minimised, and the future is not collateral. Entrepreneurship as stewardship is one of the central tasks of the age.

\section*{Markets Re-Embedded Within Moral Life}

Markets cannot answer the question of the good by themselves. They can reveal willingness to pay, but not whether a desire is well ordered. They can coordinate scarcity, but not determine justice. They can reward innovation, but not tell us whether the innovation serves life. They can generate prices, but not decide what should never be priced. They need moral surroundings: families that form trust, schools that teach judgment, religious and philosophical traditions that discipline desire, laws that restrain fraud and monopoly, public institutions that protect common goods, workers with voice, communities with agency, accounting that reveals hidden depletion, governance that protects purpose, capital structures that define enough, commons that cannot be casually enclosed, and Sabbath practices that interrupt the total claim of production.

This is not anti-market. It is the condition for markets that remain humane. A market without moral culture becomes predatory. Moral culture without economic competence becomes ineffective. A state without subsidiarity becomes overbearing, subsidiarity without solidarity leaves the weak exposed, and solidarity without institutional form dissolves into sentiment. Moral economy requires all of these in relation, and no single tradition is sufficient to provide them. Catholic Social Teaching gives us dignity, subsidiarity, solidarity, structural sin, and the common good. Protestant vocation and stewardship give us work as calling and property as entrusted responsibility. Neo-Calvinist sphere sovereignty protects institutional plurality. Ordoliberalism reminds us that markets require a moral and legal order. The cooperative and mutualist traditions show that ownership can be shared and purpose-built. Austrian economics warns against centralised arrogance and teaches the value of dispersed knowledge. The commons tradition reminds us that not all goods should be governed by private capture. The task is not to collapse these traditions into one system but to let them correct and deepen one another.

\section*{Recovering the Question of the Good}

The question ``what is an economy for?'' cannot be answered once and for all by a formula. It must be asked repeatedly, at every level: what is this company for, this investment, this technology, this platform, this law, this fund, this surplus, this growth, this freedom? The answer cannot be merely more. More growth and more return, more scale and more efficiency, more engagement, more liquidity, more data, more automation, more speed. More is not an end. It is a direction, and without a telos, more can become monstrous.

The good of an economy is human flourishing. Not as a vague slogan, but as a governing standard. Human flourishing includes material sufficiency and meaningful work, freedom from domination, stable households, trustworthy institutions, ecological responsibility, civic participation, spiritual life, friendship, beauty, knowledge, rest, and the capacity of communities to shape their own futures. No economy will realise this perfectly. Every institution will remain flawed because human beings remain flawed. Greed, vanity, fear, and self-deception do not disappear through better structures, as history has taken some pains to demonstrate. But structures can make some vices easier and some virtues harder. They can reward extraction or stewardship. They can train appetite or restraint. They can make justice ordinary or heroic.

The question of the good is therefore not abstract. It becomes law and governance, ownership and accounting, procurement, lending terms, software licences, board duties, labour protections, data rights, repayment caps, mission locks, and commons governance. The good must take form.

\section*{Capital That Serves Life}

Capital can serve life. After diagnosing extraction at such length, it would be easy to speak as if capital itself were the enemy. But that would be false. Capital is stored capacity, the ability to build across time. It can fund what does not yet exist, preserve what would otherwise collapse, help people take risks, cultivate land, build institutions, support workers, and pass gifts to future generations. Capital is dangerous because it is powerful. But power can be ordered.

Capital serves life when it is connected to real goods, governed by purpose, limited by justice, accountable to communities, patient with time, and willing to receive a fair return without claiming sovereignty. It feeds on life when it extracts from dependency, dominates through debt, encloses commons, externalises harm, captures public value, manipulates desire, demands liquidity at any cost, and treats every form of life as material for return. The difference is written into structure: who controls and who benefits, who bears risk and who can leave, who decides when enough has been reached, who protects mission, who accounts for hidden harm, who speaks for the commons, and who remains responsible after capital exits. These questions reveal whether capital is servant or master. Non-extractive capital is one attempt to answer them institutionally. Moral economy must become concrete enough to enter term sheets and bylaws, repayment schedules, governance rights, ownership forms, and reporting systems. A beautiful principle that cannot survive contact with capital markets is not yet enough. The principle must be built.

\section*{The Work Ahead}

Recovering moral economy will require many kinds of work at once, and it is worth naming them plainly rather than assuming that one kind of effort will be sufficient.

There is intellectual work, in retrieving forgotten traditions, distinguishing legitimate profit from extraction, and developing a richer theory of productive capital. There is legal work, in creating and protecting forms of enterprise where mission, stewardship, and ownership are structured differently. There is financial work, in designing funds, repayment mechanisms, and accounting systems that let capital receive a fair return without dominating enterprise. There is entrepreneurial work, in building companies that serve real needs without surrendering to exit logic. There is political work, in resisting monopoly, rent-seeking, regulatory capture, housing speculation, and digital dependency. And there is technological work, in building open and interoperable infrastructure, sustaining open source, and governing AI as a matter of public and moral concern.

There is also cultural and spiritual work: changing what we admire, celebrate, teach, fund, and call success, and recovering limits, Sabbath, humility, gratitude, mercy, and freedom from the endless hunger for more. No single person or firm, no one fund or church or state, can do all of this. That is why moral economy must be an ecology. It needs philosophers and accountants, theologians and lawyers, founders and workers, investors and regulators, engineers and contemplatives, policymakers and maintainers, households and universities, churches and cooperatives, public institutions and private enterprises. It needs all the places where economic life is actually made.

\section*{The Final Question}

The old question returns at the end because it was never only old. It is the question beneath the next investment and the next company, the next platform, the next AI system, the next housing policy, the next public procurement decision, the next open-source licence, the next debt instrument, the next board meeting, the next act of work.

What is this for?

If the answer is extraction, we should not be surprised by the world that follows. We will have more activity and less life, more connection and less communion, more intelligence and less wisdom, more wealth and less security, more innovation and less agency, more growth and less flourishing.

But if the answer is life, then another economy becomes possible. Not perfect. Not pure. Not free from conflict, scarcity, failure, or sin. But better ordered, more truthful, more accountable, more humane. An economy where capital is useful but not ultimate, where markets coordinate but do not rule, where ownership carries obligation, where work can be vocation, where technology remains servant, where commons are protected, where finance is disciplined by surplus and limits, where enterprise is a form of stewardship, and where human beings are not required to justify their existence before the demands of return.

That is the economy this book has tried to imagine.

\begin{center}\itshape An economy ordered toward human flourishing.\end{center}

\begin{center}\itshape An economy in which capital serves life.\end{center}


\end{document}
